\documentclass[a5paper,10pt]{book}
\usepackage[brazilian]{babel}
\usepackage[T1]{fontenc}
\usepackage[utf8]{inputenc}
\usepackage{graphicx}
\usepackage{geometry}
\usepackage{longtable}
\usepackage{hyperref}
\usepackage{indentfirst} % Para 'indentar' o primeiro parágrafo de cada seção
\geometry{verbose,tmargin=2cm,bmargin=2cm,lmargin=2cm,rmargin=2cm,headsep=5mm,footskip=5mm} % 'headsep' é a separação entre o cabeçalho e o texto e o 'footskip' é a distância entre as notas de rodapé e o texto.
\pagestyle{plain}
% \usepackage[normalem]{ulem} -- Pacote para sublinhar texto.
% o \begin{abstract} -- ambiente 'resumo'
% \newline \newpage \linebreak
% ambientes: flushright, flushleft, center, description, itemize, quotation, quote, enumerate
\renewcommand{\baselinestretch}{1.25}
\setcounter{secnumdepth}{0}
% \maketitle \tableofcontets
%
% para colar uma figura:
%
% \begin{figure}[!htb]
% \centering
% \fbox{\includegraphics[scale=1]{FIGURA.jpeg}}
% \caption{descriç}
% \label{fig1}
% \end{figure}
% 
% O h significa here, t top, b bottom e o ! dá prioridade ao aparecimento h
%
% para tabelas
%
% \begin{tabular}{|c|l|r|}
% ítem 1x1 & ítem 1x2 & ítem 1x3 \\
% ítem 2x1 & ítem 2x2 & ítem 2x3 \\
% end{tabular}
%
% O c, o l e o r definem o alinhamento da coluna. O & separa colunas na mesma linha e o \\ vai pra próxima linha
% 
\title{Guia Definitivo de Customização do UNO}
\author{Peterson Silva}
\date{}

\begin{document}

\maketitle

\begin{center}
\url{http://petercast.net}
\end{center}

\vfill

\begin{center}
Esta obra é de \textbf{Domínio Público}.
\end{center}

\begin{figure}[!h]
\centering
\fbox{\includegraphics[scale=1]{dp.png}}
\end{figure}

\begin{center}
Para mais informações, visitar: \\

\url{http://creativecommons.org/publicdomain/zero/1.0/}
\end{center}

%\vspace{1cm}

%\begin{center}
%Organização: Peterson Silva

%Colaboradores:
%\end{center}

\tableofcontents

\section*{Introdução}
\addcontentsline{toc}{section}{Introdução}

Este guia é destinado a todas aquelas pessoas que já se encontraram em situações difíceis por ocasião das regras deste jogo incrível e flexível que é UNO.

À bem da verdade, o UNO foi uma encarnação nova para um jogo antigo: a ideia é muito óbvia e simples para acharmos que ele é um jogo original. Entretanto, ao invés de usarmos as cartas de figura do baralho normal (Rei, Rainha, Valete) para representar as ações como "pular o próximo jogador", "comprar duas cartas", etc, o baralho de UNO foi criado especificamente para este tipo de jogo. Esse é um apelo visual que torna o jogo ainda mais coerente e atraente, e não apenas algo inventado, usando-se de gambiarras com o velho e cinzento baralho normal.

Esta invenção, contudo, não deve ser vista como muito recente; ela existe desde a década de 70 nos Estados Unidos, e atualmente existem dezenas de variações de UNO --- versões em que você pode jogar na piscina, ou usando uma aparelho que ``cospe'' as cartas. Não tratamos de nenhuma dessas aqui: nos concentramos no bom e velho jogo de cartas e como suas regras podem ser flexibilizadas.

Apesar de não ser um jogo novo e revolucionário, UNO é divertido, viciante, \textit{encantador}. Sua simplicidade é um duplo fardo: mesmo ao comprarmos um baralho novo em folha, passamos os olhos pela folha de regras que o acompanha, damos uma olhada em como são as regras em francês, e então jogamos a folha no lixo ou em alguma gaveta escura. Isso faz com que muitas regras sejam distorcidas, e uma miríade de novas regras entrem em campo, fazendo com que quase nunca se jogue o mesmo UNO em grupos de jogadores diferentes.

Entretanto, isso de forma alguma é um problema; o jogo não se modifica tanto a ponto de não ser ``o velho UNO''. É por isso que esse guia foi escrito: para ajudar a fazer uma escolha consciente quanto às diferentes regras que podem ser aplicadas para tornar o jogo mais agradável de acordo com as preferências dos jogadores.

Ele possui algumas seções distintas: em primeiro lugar teremos contato com as regras oficiais do jogo. Como ele deve ser jogado. Logo depois veremos como a interpretação pode entrar em cena, já bifurcando levemente tipos distintos de jogo.

Entraremos então em contato com a essência estratégica do jogo. O que se deve fazer para ganhar? Quais as melhores artimanhas, qual a lógica a ser seguida?

Então poderemos discutir as regras alternativas, cada uma em sua seção: regras que mudam o modo de jogar ou adicionam funcionalidades às cartas, ou mesmo mudam as condições de vitória do jogo. Cada uma é discutida em suas consequências para a dinâmica e estratégia do jogo.

Bom divertimento!

\chapter{Regras Oficiais}
\label{oficiais}

A seguir serão apresentadas as regras oficiais do jogo. Elas serão fielmente reproduzidas (\textit{Verbatim}) do folheto original que acompanha as embalagens brasileiras do jogo\footnote{Folheto encontrado em uma embalagem comprada em 2010.}, distribuído pela empresa Mattel.

\section{Regras em Verbatim}

\vspace{0.5cm}

\textsc{\large{Conteúdo}}

\textbf{108 cartas, sendo}

\textbf{19 Cartas azuis} --- 0 a 9

\textbf{19 Cartas Verdes} --- 0 a 9

\textbf{19 Cartas Vermelhas} --- 0 a 9

\textbf{19 Cartas Amarelas} --- 0 a 9

\textbf{8 Cartas Comprar Duas Cartas} --- 2 de cada em azul, verde, vermelho e amarelo

\textbf{8 Cartas Inverter} --- 2 de cada em azul, verde, vermelho e amarelo

\textbf{8 Cartas Pular} --- 2 de cada em azul, verde, vermelho e amarelo

\textbf{4 Curingas}

\textbf{4 Curingas Comprar Quatro Cartas}

\vspace{0.5cm}

\textsc{\large{Objetivo do Jogo}}

Ser o primeiro jogador a marcar 500 pontos. Os pontos são marcados livrando-se de todas as cartas que estão na sua mão antes do adversário. Você marca pontos pelas cartas que ficaram nas mãos dos seus adversários.

\textsc{\large{Como Jogar}}

Cada jogador compra uma carta. A pessoa que tirar a carta mais alta faz a distribuição. Cartas de ação contam como zero nesta parte do jogo.

Quando as cartas estiverem embaralhadas, são distribuídas 7 cartas para cada jogador.

As cartas restantes são colocadas viradas para baixo para formar o MONTE. A carta de cima do MONTE é virada para cima para começar uma pilha de DESCARTE. Se uma Carta de Ação for a primeira a ser aberta no MONTE, são adotadas regras específicas (veja FUNÇÕES DAS CARTAS DE AÇÃO).

O jogador à esquerda do distribuidor começa o jogo. Ele precisa jogar uma carta que combine com a que está na pilha de DESCARTE em valor, cor ou naipe. Por exemplo, se a carta for um \textit{7 vermelho}, o jogador precisa descartar uma carta \textit{vermelha} ou um \textit{7 de qualquer cor}. Alternativamente, o jogador pode descartar um \textit{Curinga} (Veja FUNÇÕES DAS CARTAS DE AÇÃO).

Se o jogador não tiver uma carta que combine com a que está na pilha de DESCARTE, ele \textit{precisa} comprar uma carta do MONTE. Se a carta comprada servir, o jogador pode baixá-la na mesma vez. Caso contrário, a vez passa para o próximo jogador.

Os jogadores podem optar por não jogar uma carta válida que esteja na mão dele. Nesse caso, o jogador precisa comprar uma carta do MONTE. Se esta puder ser jogada, a carta pode ser descartada na mesma vez, mas o jogador não pode jogar uma carta que estiver na sua mão depois de ter feito a compra.

\textsc{\large{Funções das Cartas de Ação}}

As funções das Cartas de Ação, e quando podem ser jogadas, estão definidas abaixo.

\textbf{Carta Comprar Duas Cartas} --- Quando esta carta é jogada, o próximo jogador \textit{precisa comprar 2 cartas} e \textit{perderá a vez}. Esta carta só pode ser jogada sobre uma de mesma cor e sobre outras cartas Comprar Duas Cartas. Se ela for aberta no começo do jogo, as mesmas regras se aplicam.

\textbf{Carta Inverter} --- Ela simplesmente inverte a direção do jogo. Se o jogo estiver girando para a esquerda, passa a girar para a direita, e vice-versa. A carta só pode ser jogada sobre uma da mesma cor ou sobre outra carta Inverter. Se esta carta for aberta no começo do jogo, o distribuidor começa e, depois, o jogo se desloca para a direita em vez de se deslocar para a esquerda.

\textbf{Carta Pular} --- O jogador seguinte àquele que descartou essa carta perde a vez --- é ``\textit{pulado}''. A carta só pode ser jogada sobre uma carta da mesma cor ou sobre outra carta Pular. Se uma carta Pular for aberta no começo do jogo, o jogador à esquerda do distribuidor é ``pulado'', e o jogador à esquerda deste último começa o jogo.

\textbf{Curinga} --- A pessoa que jogar esta carta exige \textit{qualquer} cor para continuar a jogar, inclusive a cor que está sendo jogada no momento, se assim desejar. Um Curinga pode ser jogado a qualquer hora --- mesmo que o jogador tenha na mão outra carta que possa ser jogada. Se um Curinga for aberto no começo do jogo, o jogador à esquerda do distribuidor determina a cor que dará sequência ao jogo.

\textbf{Curinga Comprar Quatro Cartas} --- Esta é a melhor carta. A pessoa que jogar esta carta exige a cor que continua o jogo. Além disso, o próximo jogador precisa \textit{comprar 4 cartas} do MONTE e \textit{perde a vez}. Infelizmente, esta carta só pode ser jogada quando o jogador que a possui não tiver na mão uma carta de mesma cor da que está na pilha de DESCARTE. Entretanto, se o jogador tiver cartas com o mesmo número ou Cartas de Ação, o Curinga Comprar Quatro Cartas pode ser jogado. Um jogador que possuir um Curinga Comprar Quatro Cartas pode tentar blefar e jogar uma carta inapropriada, mas certas regras se aplicam se ele for apanhado (veja PENALIDADES). Se esta carta for aberta no começo do jogo, é devolvida ao baralho e uma outra carta é aberta.

\textsc{\large{Saindo}}

Quando só restar uma carta ao jogador, ele precisa gritar ``UNO'' (que significa ``uma''). Se ele não o fizer, terá de comprar 2 cartas do MONTE. Entretanto, isto só é necessário se ele for apanhado por um dos outros jogadores (veja PENALIDADES).

Quando o jogador não tiver mais nenhuma carta, a partida acabou. São marcados os pontos (veja PONTUAÇÃO) e o jogo recomeça.

Se a última carta jogada em uma partida for a carta Comprar Duas Cartas ou o Curinga Comprar Quatro Cartas, o próximo jogador \textit{precisa} comprar 2 ou 4 cartas, respectivamente. Estas cartas são contadas na totalização dos pontos.

Se nenhum jogador estiver sem cartas quando o MONTE acabar, o baralho é novamente embaralhado e o jogo continua.

\textsc{\large{Pontuação}}

O primeiro jogador que se livrar das cartas recebe pontos pelas cartas que \textit{ficaram} nas mãos dos adversários, como segue:

Todas as cartas numeradas (0--9) --- Valor Nominal

Comprar Duas Cartas --- 20 Pontos

Inverter --- 20 Pontos

Pular --- 20 Pontos

Curinga --- 50 Pontos

Curinga Comprar Quatro Cartas --- 50 Pontos

\textsc{\large{Vencendo o Jogo}}

O VENCEDOR é o primeiro jogador que fizer 500 pontos. Entretanto, os pontos do jogo podem ser marcados mantendo um total parcial dos pontos com que cada jogador é apanhado no fim de cada partida. Quando um jogador alcançar os 500 pontos, o jogador com menos pontos é o vencedor, ou seja, aquele que não deixar sobrarem muitas cartas na mão.

\textsc{\large{Penalidades}}

O jogador que esquecer de dizer ``UNO'' antes que a penúltima carta toque a pilha de DESCARTE, mas se lembrar (e gritar ``UNO'') antes que qualquer outro jogador ``perceba'', está salvo e não está sujeito à penalidade. Os jogadores não podem ser apanhados por não dizer ``UNO'' depois que o próximo jogador começar a sua vez. ``Começar a vez'' é definido como comprar uma carta do MONTE ou tirar uma carta da mão para jogar.

Jogadores que derem palpites aos outros jogadores \textit{precisam} comprar 2 cartas do MONTE.

Se um jogador jogar uma carta errada e isto for notado por qualquer um dos demais jogadores, ele ``\textit{precisa}'' pegá-la de volta e comprar mais 2 cartas extras do MONTE. O jogo continua com o próximo jogador.

Se um Curinga Comprar Quatro Cartas for exibido indevidamente (isto é, se o jogador possuir uma carta da mesma cor da que está na pilha de DESCARTE) e a pessoa que estiver jogando for desafiada, a mão \textit{precisa} ser mostrada primeiro ao jogador que fez o desafio.

Se o Curinga Comprar Quatro Cartas tiver sido jogado indevidamente, o jogador responsável precisa comprar 4 cartas. Se a carta foi jogada corretamente, o desafiante precisa comprar 2 cartas, além das 4. O desafio só pode ser feito pelo jogador obrigado a comprar 4 cartas depois que o Curinga Comprar Quatro Cartas tiver sido descartado na mesa.

\textsc{\large{Jogo em Dupla, Parceiros \newline \indent e Torneios em Várias Mesas}}

\textbf{Jogo em Dupla (jogo UNO com dois jogadores)}

Joga-se com as seguintes regras especiais:

\begin{enumerate}
\item{Jogar uma carta Inverter funciona como Pular. O jogador que joga a carta Inverter pode imediatamente jogar outra carta.}
\item{A pessoa que está jogando uma carta Pular pode imediatamente jogar outra carta.}
\item{Quando uma pessoa joga uma carta Comprar Duas Cartas e o outro jogador tiver comprado as 2 cartas, o jogo volta ao primeiro jogador. O mesmo princípio se aplica à carta Curinga Comprar Quatro Cartas. As regras comuns do jogo de cartas UNO se aplicam aos demais casos.}
\end{enumerate}

\textbf{Parceiros}

Os parceiros sentam-se de frente um para o outro. Quando um dos parceiros ficar sem cartas, a partida terminou. Todos os pontos nas mãos dos dois parceiros adversários são somados e passados para o time vencedor.

\textbf{Variações}

Com quatro jogadores podem ser jogadas quatro partidas e em cada uma mudam os parceiros. Todos os jogadores mantêm-se a par dos pontos feitos em cada uma das suas parcerias. Podem ser jogadas diversas rodadas e a pessoa que fizer mais pontos é declarada a vencedora.

Com oito jogadores, podem ser jogados dois jogos separados em duas mesas e cada jogador tem outro jogador como parceiro durante quatro partidas (um total de 28 partidas). Pontuação como acima.

\textbf{Jogo de Cartas DESAFIO UNO}

Este jogo é pontuado mantendo-se um total parcial das cartas que estão na mão de cada jogador. Quando um jogador alcança um valor combinado, possivelmente 500, o jogador é eliminado do jogo. Quando restarem apenas dois jogadores, eles jogam de igual pra igual. O jogador que alcançar ou superar o valor combinado perde. O ganhador da última partida é declarado o vencedor do jogo (veja regras especiais para JOGO EM DUPLA).

\textbf{Divirta-se com UNO --- e que vença o melhor ou aquele com mais sorte!}

\subsection{Ambiguidade das Regras Oficiais}

Agora que as regras oficiais do jogo foram expostas, podemos analisar com profundidade alguns de seus aspectos e verificar que, na verdade, algumas questões não foram completamente respondidas. Vamos a elas.

Depois de cada ambiguidade analisada, uma seção com opções --- escolhas a serem feitas entre um modo de jogar e outro --- será apresentada. Cada opção possui um código para que você possa fazer referência às suas escolhas quanto ao jogo.

\subsubsection{Procedimentos Pré-jogo}

A regra diz que:

\begin{quotation}
Cada jogador compra uma carta. A pessoa que tirar a carta mais alta faz a distribuição. [...] Quando as cartas estiverem embaralhadas, são distribuídas 7 cartas para cada jogador.

[...] A carta de cima do MONTE é virada para cima para começar uma pilha de DESCARTE. [...] O jogador à esquerda do distribuidor começa o jogo.
\end{quotation}

A questão é que o jogo parece ter tornado mais simples (e bem diferente) a tradição de distruibuição de cartas aplicada a jogos como Truco, Canastra e Caxeta (ou ``Pife'').

Na maioria dos jogos, uma pessoa deve embaralhar as cartas. Depois, o distribuidor em questão deve passar as cartas para o jogador à \textbf{esquerda}, que então \textit{corta} as cartas. Durante o corte, este jogador deve pegar uma carta do local onde o corte foi feito e virá-la para a cima --- no caso do Truco ou da Caxeta, esta carta em questão define o \textit{curinga}. No caso do UNO, apenas determina quais cartas poderão ser jogadas pelo jogador que inicia a partida.

Após feito o corte, as duas partes do baralho encontram-se novamente na ordem em que estavam antes, e o distribuidor passa então a distribuir as cartas, começando pelo jogador imediatamente à sua direita e prosseguindo no sentido anti-horário até que dê as suas próprias cartas\footnote{Algumas regras podem ser mais complexas; no Truco, se o jogador que corta escolher a carta de baixo do monte para definir o curinga, o distribuidor deve distribuir as cartas de baixo do monte. Além disso, podem existir regras específicas quanto ao número de cartas dadas aos jogadores a cada volta no sentido anti-horário até que todos os jogadores tenham o número de cartas estipulado para começar o jogo.}.

Após isso, o jogador à \textit{direita} inicia o jogo.

Tudo isso parece ter sido apenas subentendido ou ignorado nas regras do UNO; considerando, contudo, que eles não fazem menção alguma à ideia do jogador que faz o corte, e explicitam que o jogo deve começar pela \textit{esquerda}, é mais provável que toda a convenção tenha sido ignorada e o que ocorra seja um procedimento mais simples (até pela natureza comercial, que tende a simplificar as coisas, do jogo): o jogador que tira a carta mais alta faz a distribuição, na ordem que quiser, qualquer um pode virar uma carta do monte para formar a pilha de descarte, e o jogador à esquerda inicia o jogo.

O que a regra não parece clarificar, contudo, é quantas vezes este procedimento precisa ser repetido: apenas no início do \textit{jogo} ou a cada \textit{partida} dentro do mesmo jogo? Considerando que se o procedimento for feito a cada início de partida, a sorte acaba por determinar quem começa o jogo mais vezes. O curso normal do jogo geralmente faz com que a vantagem de começar uma partida em geral seja pequena, mas ainda assim é uma vantagem. Portanto, para maior equilíbrio, seria mais natural fazer com que, uma vez definido quem dá as cartas \textit{na primeira partida do jogo}, os distribuidores sejam escolhidos através de um ciclo no sentido horário (à esquerda), dando assim a todos a chance de embaralhar as cartas e iniciar a partida.

\paragraph{Opções}

\begin{itemize}
\item{Adotar os procedimentos pré-jogo simplificados do UNO. Código 1.}
\item{Adotar os procedimentos pré-jogo tradicionais. Código 2.}
\item{Adotar os procedimentos pré-jogo simplificados do UNO apenas no começo do jogo, e depois partir para o uso dos tradicionais. Código 3.}
\end{itemize}

\subsubsection{Correções}

\begin{quote}
Se um jogador jogar uma carta errada e isto for notado por qualquer um dos demais jogadores, ele ``precisa'' pegá-la de volta e comprar mais 2 cartas extras do MONTE. O jogo continua com o próximo jogador.
\end{quote}

As regras não explicam o que ocorre quando o erro é percebido algumas rodadas adiante (mesmo duas ou três). Quem seria ``o próximo jogador''? O próximo, que iria jogar caso o erro não tivesse sido descoberto, ou o jogador que jogou logo após o outro que cometeu o erro?

Além disso, que tipo de correção deve ser feita ao erro? O jogador que o cometeu deve pegar a carta de volta, comprar duas cartas e jogar alguma outra, desta vez apropriada? Ou deveria apenas pegar sua carta de volta, comprar mais duas e o jogo continua sem a participação desse jogador?

Tudo isto parece uma decisão que se relaciona mais ao valor que se tem em relação ao jogo: o que é mais priorizado? Que o jogo seja justo, correto, ou que seja mais rápido e dinâmico? Este é um tipo de decisão (dentre outras) que se encontra no cerne de praticamente todas as variações e personalizações do jogo contidas neste guia.

Suponha que você esteja jogando dentro de um sistema de apostas. É primordial, assim, que o jogo seja o mais justo possível; logo, o que deve ser feita é uma correção imediata e o jogo deve reiniciar a partir do primeiro jogador afetado, que precisa refazer a sua jogada.

Contudo, alguns jogadores podem ter se beneficiado do erro alheio, e para que o jogo não seja totalmente refeito, o que se pode fazer é continuar o jogo normalmente, mas com o jogador que errou comprando duas cartas e recebendo suas cartas da jogada equivocada de volta.

Há um ponto em comum, contanto: erros causam consternação, dúvida, problemas de estratégia não-necessários. Puni-los para não incentivá-los é uma prática razoável; não faz sentido, pois, que o jogador tenha chance de jogar outras cartas (desta vez, corretas) caso tenha errado. Ainda assim, é uma opção.

\paragraph{Opções}

\begin{itemize}
\item{O erro é corrigido com o jogador pegando de volta a carta que jogou erroneamente, com a partida reiniciada daquele ponto. Código 1.}
\item{O jogador errado é punido, mas as jogadas subsequentes não são alteradas. A partida prossegue. Código 2.}
\end{itemize}

\subsubsection{O Desafio do Curinga}

A regra diz:

\begin{quote}
Se um Curinga Comprar Quatro Cartas for exibido indevidamente (isto é, se o jogador possuir uma carta da mesma cor da que está na pilha de DESCARTE) e a pessoa que estiver jogando for desafiada, a mão precisa ser mostrada \textit{primeiro} ao jogador que fez o desafio.
\end{quote}

A ênfase dada ao termo \textit{primeiro} parece sugerir que as cartas devam ser mostradas a todos os jogadores, mas primeiramente para o desafiante --- o que certamente é desnecessariamente prejudicial ao jogador desafiado, uma vez que todos passariam a ver suas cartas, embora apenas uma pessoa já basta para determinar se a jogada foi legítima ou não.

\paragraph{Opções}

\begin{itemize}
\item{O jogador desafiado deve mostrar as suas cartas a todos. Código 1.}
\item{O jogador desafiado deve mostrar suas cartas apenas ao desafiante. Código 2.}
\end{itemize} 

\textbf{Além disso}, se um blefe for descoberto, a carta jogada ainda é válida? O jogador que a jogou ainda tem, pelo menos, o direito de exigir a cor? Me parece um modo de não tornar o blefe uma ação muito arriscada, estimulando-a e tornando o jogo mais divertido e complexo.

\paragraph{Opções}

\begin{itemize}
\item{Se o blefe for descoberto, o desafiado pelo menos exige uma cor. Código 1.}
\item{Se o blefe for descoberto, a carta torna-se inválida. Código 2.}
\end{itemize}

\subsubsection{A Cor do Descarte}

A regra diz o seguinte:

\begin{quote}
[...] Esta carta \textit{[Curinga Comprar Quatro Cartas]} só pode ser jogada quando o jogador que a possui não tiver na mão uma carta de mesma cor da que está na pilha de DESCARTE.
\end{quote}

Esta é curiosa, e até mesmo filosófica: suponhamos que um jogador jogue esta carta. O próximo deve comprar quatro, e perde a vez. O jogador que a jogou escolhe a próxima cor: amarelo. O jogador que vai jogar agora possui uma carta amarela, mas possui também um Curinga Comprar Quatro Cartas.

Teoricamente, a \textit{carta que está pilha de DESCARTE} não tem cor. Ou poderíamos dizer que ela é ``preta'', por ser um Curinga. Nesse caso, você tem uma carta da mesma cor (o Curinga) pra jogar, e portanto deve jogar esta carta --- o que é redundante, porque esta carta é o Curinga.

Contudo, não faz sentido considerar os Curingas como sendo de uma cor própria; eles assumem cores diferentes. Se o jogador anterior exigiu a cor amarela, o Curinga deve ser tratado como uma carta amarela. Dessa forma, a pessoa tem uma carta amarela pra jogar e ``não pode'' (i.e a não ser que blefe) jogar o Curinga Comprar Quatro Cartas.

\subsubsection{Fim das Cartas}

A regra diz que:

\begin{quote}
Se nenhum jogador estiver sem cartas quando o MONTE acabar, o baralho é novamente embaralhado e o jogo continua.
\end{quote}

Implica-se, de boa fé, que as instruções na verdade sejam estas: ``Se nenhum jogador estiver sem cartas quando o MONTE acabar, o DESCARTE é novamente embaralhado e a \textit{partida} continua.''. Se interpretarmos a regra literalmente (principalmente se considerarmos a distinção entre partida e jogo), contudo, vemos a evidente ambiguidade: ou a pilha de descarte é embaralhada para que a partida continue, ou a partida termina e uma nova começa (dando continuidade ao \textit{jogo}).

\paragraph{Opções}

\begin{itemize}
\item{Se o MONTE acabar, o DESCARTE é embaralho e a partida continua. Código 1.}
\item{Se o MONTE acabar, a partida termina e outra é iniciada. Código 2.}
\end{itemize}

\subsubsection{Desafio UNO}

Esta é uma variação oficial do jogo. A regra explica:

\begin{quote}
Este jogo é pontuado mantendo-se um total parcial das cartas que estão na mão de cada jogador. Quando um jogador alcança um valor combinado, possivelmente 500, o jogador é eliminado do jogo. Quando restarem apenas dois jogadores, eles jogam de igual pra igual. O jogador que alcançar ou superar o valor combinado perde. [...]
\end{quote}

Há um problema com essa variação de jogo: como exatamente se conta os pontos do total \textit{parcial} das cartas? Alguma pessoa, que não está jogando, deve manter um registro de cada jogador?
P
E como é a dinâmica de jogo? Como é possível \textit{alcançar} uma pontuação nas cartas à mão, se a dinâmica de jogo leva o jogador a descartar cartas, não a acumulá-las?

Este é um problema que as regras não resolvem, e a interpretação, portanto, torna-se livre para a execução desta modalidade de jogo.

\subsection{Os Próximos Exemplos}

Os próximos exemplos do livro tratarão do UNO em sua variante que considerei a mais direta e mais simples, com as interpretações mais baseadas no bom senso possíveis (se é que isto é possível): a variante \emph{11211}.

O código é lido da seguinte forma: ele possui cinco números porque nesta seção existem 5 seções com opções (com escolhas a fazer sobre o jogo). Os números estão ordenados de acordo com a ordem de aparecimento dessas seções. O primeiro número é 1; a primeira seção é Procedimentos Pré-jogo; ou seja, das duas opções contidas naquela seção, a de código 1 foi escolhida. E assim por diante.

Nesta variação, os procedimentos pré-jogo simplificados do UNO são adotados, as correções são precisas e o jogo todo precisa ser reiniciado, o jogador desafiado deve mostrar suas cartas apenas ao desafiante, e aquele que for pego no blefe pelo menos exige a próxima cor, e se o MONTE acaba, a partida segue com o DESCARTE embaralhado.

\section{Estratégias do Jogo Oficial} 

Há quem diga que UNO é um jogo de pura sorte, e há pouco espaço para a intervenção estratégica no jogo; que há muitas variáveis e que, portanto, é inútil tentar agir conscientemente em relação ao jogo. Acredito que não, e que existem várias estratégias a serem empregadas.

Esta seção se divide em três: Em primeiro lugar, daremos uma olhada em estratégias gerais, que são importantes guias de ação. Quando se se joga em dupla, entretanto, algumas diferenças devem ser levadas em conta; olharemos estas estratégias em particular e, por fim, pensaremos em estratégias e na dinâmica do jogo de parceiros.

\subsection{Estratégia Geral}

A dinâmica do jogo em suas regras oficiais é simples: você deve descartar suas cartas o mais rápido possível, não apenas para ganhar o jogo, mas para evitar conceder ao eventual vitorioso um grande número de pontos.

Ao mesmo tempo, você tem a sua disposição cartas que prejudicam os outros adversários, e deve procurar usá-las também para impedir que eles vençam. Existem cartas que não só prejudicam os adversários, mas também te dão benefícios essenciais, como as cartas Curinga.

Jogar estas cartas ofensivas no começo do jogo não é a melhor estratégia possível, uma vez que as cartas de todos os jogadores ainda estão equilibradas, e isso pode acabar dando também poder ofensivo a eles. Jogar cartas de ação ou vantajosas no começo do jogo também não é bom, pois perde-se a oportunidade de usá-las quando uma delas for essencial para impedir a vitória alheia ou alcançar a própria, mais adiante na partida.

Contudo, se as cartas forem conservadas por muito tempo, uma melhor estratégia de um adversário pode prevalecer, e muitos pontos serão entregues a ele.

Neste cenário, torna-se óbvio que criar uma boa estratégia, num jogo individual, significa uma \textbf{escolha do tempo certo, em face da situação dos adversários, para jogar as cartas certas e proporcionar uma continuidade ao seu jogo}. Aspectos da formação dessa estratégia serão discutidos a seguir.

\subsubsection{Continuidade do Jogo}

Uma característica importante a notar é tentar o máximo possível provocar a \textit{continuidade} do próprio jogo. Isto significa jogar de um tal modo que a sua próxima vez seja vantajosa para sua estratégia geral.

Ao jogar uma carta, você deve fazê-lo esperando que, quando for sua vez de jogar novamente, você seja capaz de jogar outra carta (e até planejar qual carta será) ao invés de não ter mais opções e ser obrigado a comprar outra carta.

\subsubsection{Escolha das Cartas em Virtude da Continuidade}

Suponha um jogo em que você tenha as seguintes cartas à mão\footnote{Cartas escolhidas aleatoriamente, a partir de um baralho de UNO real.}:

\vspace{0.5cm}

\textbf{Pular Azul \hspace{1cm} 1 Amarelo \hspace{1cm} 2 Azul \hspace{1cm} 6 Azul \hspace{1cm} 6 Azul \hspace{1cm} Inverter Azul \hspace{1cm} Comprar Duas Cartas Verde}

\vspace{0.5cm}

Supondo que todos os seus adversários tenham também 7 cartas, e que a carta na pilha de descarte seja um \textbf{2 Azul}. Que carta deveria ser jogada?

Você tem à sua disposição 5 cartas azuis. Duas delas são cartas de ação, que certamente serão muito mais úteis em rodadas vindouras, uma vez que ninguém está em real posição de vencer o jogo neste momento. Neste caso, sobram apenas três alternativas: dois 6 e um 2.

Pode parecer uma escolha irrelevante, mas imagine a próxima rodada. Se um 2 for jogado, mesmo não sendo azul (na verdade, é impossível que ele seja azul, afinal você sabe onde estão os dois 2 azuis do baralho e sabe que não existem outros), você terá um dois para jogar, e assim você não precisará comprar uma carta se for um 2 vermelho, por exemplo. Mas, se você jogar um dois agora, nesta rodada, não terá essa possibilidade mais adiante.

Contudo, se nesta rodada você jogar um 6 azul, na próxima rodada não só você terá um 2 azul para jogar caso um 2 vermelho seja posto no descarte, como também terá um 6 caso algum 6 seja jogado. Ao ter duas cartas iguais, você pode jogar uma delas e ainda assim conservar a outra, aumentando as chances de você poder dar continuidade ao seu jogo.

Quis o futuro\footnote{A carta escolhida foi aleatória, assim como a mão de exemplo, e a partir de agora não farei mais esta observação. Todas essas cartas foram escolhidas aleatoriamente.} que na sua próxima vez, a carta do descarte seja um 2 verde. Portanto, com efeito, o fato de você ter jogado um 6 azul, e não um 2 azul, faz com que você tenha opções: pode descartar um 2 azul ou um Comprar Duas Cartas Verde. O adversário à sua esquerda (neste sentido do jogo, vamos considerar que ele é o próximo a jogar depois de você) tem 6 cartas; talvez não seja a melhor hora para jogar esta carta. Se a sorte e as estratégias dos adversários vão favorecer esta espera, só o tempo dirá.

Você joga o dois azul. Na sua próxima rodada, a carta é um 6 amarelo. Aqui chegamos a uma decisão crucial, a decisão entre um número e uma cor.

Você poderia optar por jogar o 6 azul ou o 1 amarelo. Ao jogar o 1 amarelo, você deixa de ter qualquer carta número 1, e qualquer carta amarela. Ao jogar o 6 azul, deixa de ter qualquer 6 --- mas conserva suas cartas azuis.

Em qualquer mão, sem considerar outros atributos da situação, escolher a cor sobre o número é sempre o certo a fazer. Há muito mais chance de a próxima carta na sua vez ser alguma carta amarela do que um 6 (existem 25 cartas amarelas no jogo, sem contar os curingas; em contrapartida, existem oito cartas 6 no baralho), portanto conservar o 1 e jogar o 6 seria o mais lógico.

Entretanto, há sempre outras coisas a levar em consideração. Aqui, ao descartar um 6 amarelo, todas as suas cartas azuis serão cartas de ação. Não é um problema absoluto, mas tudo depende sempre da situação. Em outros jogos (não envolvendo este exemplo) jogar uma carta azul pode ser apropriado. Por exemplo, se o próximo jogador estiver com apenas uma ou mesmo duas cartas na mão, e você \textit{deduziu} (devido a observações das rodadas anteriores; mais sobre isso adiante) que ele não tem cartas azuis na mão, você deve jogar a carta azul para impedir que ele se aproxime ainda mais da vitória. Entretanto, em situações normais de simples escolha entre número e cor, onde você acaba eliminando uma cor ou um número da sua mão, \textbf{escolha eliminar um número}.

O 6 azul é jogado. A próxima carta é um 5 azul; você ou inverte o jogo ou pula o próximo jogador. Note que, em qualquer situação, o efeito acaba sendo evitar que a próxima pessoa jogue. Contudo, outras variáveis precisam ser observadas.

Se você acredita que o próximo jogador na sequência atual possua cartas perigosas pra você, você não deve inverter o jogo, porque você passaria a ser o próximo depois da jogada \textit{dele}.

Por outro lado, se o jogador que joga após o próximo possua uma ou mesmo duas cartas, pular o próximo pode não ser uma boa ideia; afinal de contas, isso vai fazer com que aquele jogador com poucas cartas jogue diretamente. Permitir que o próximo jogador faça uma interferência, mudando, talvez, a cor da carta no descarte, ou mesmo invertendo o jogo, fazendo com que ele compre mais cartas, etc, pode ser uma boa ideia. Pode não ser uma boa ideia se o próximo jogador \textit{também tenha poucas cartas}.

Se realmente não for desejável pular o jogo ou invertê-lo (ou se ainda você deseja conservar a carta especial para depois), lembre-se que você pode optar por não jogar uma carta que você possa jogar, e compre uma carta ao invés disso (só que, uma vez comprada uma carta, você não poderá jogar uma carta que já possuía antes, apenas a nova carta).

Vamos continuar, supondo que você tenha invertido o jogo. A próxima carta é um 1 verde. Considerando que a sua mão agora é formada por um Comprar Duas Cartas Verde, um 1 Amarelo e um Pular Azul, você tem novamente o dilema entre número e cor --- exceto que agora qualquer jogada sua vai eliminar uma cor da sua mão. Lembre-se, contudo, que provavelmente os seus adversários também agora possuem poucas cartas na mão. Seria uma boa ideia fazer o próximo jogador comprar duas cartas.

Neste ponto do jogo, é difícil criar estratégias com base na continuidade; imagine a situação corrente. Você possui apenas um Pular Azul e um 1 Amarelo. Uma carta não oferece continuidade à outra; o máximo a fazer é esperar que a dinâmica de jogo faça sua mágica e mude substancialmente a carta no descarte, a ponto de o seu amarelo virar um azul, ou a sua carta azul se transformar no amarelo. Entretanto, perceba: se você jogar o Pular Azul, menos pessoas jogarão na rodada, diminuindo a chance de que a cor na pilha do descarte mude.

\paragraph{Conclusão}

\begin{itemize}
\item{Vimos que é importante jogar de um modo a maximizar as suas chances de continuar jogando, sem parar para comprar mais cartas.}
\item{Vimos também que é importante descobrir a hora oportuna de descartar as cartas não-numéricas, sendo desvantajoso o uso precoce, porém perigoso o contínuo atraso do uso.}
\item{A chance de formar uma estratégia coerente baseada no princípio da continuidade diminui com o número de cartas. Portanto, quanto maior o número de cartas e, em especial, o de cartas de ação, maior o poder de gerenciar o jogo. Pode-se deduzir que uma estratégia vitoriosa passa por tentar \textit{gerenciar} o jogo no início, mas não se deixar levar pela ilusão de que controlar as coisas basta: quanto mais você tenta controlar o jogo, mais os adversários lutam para fugir a esse controle, e as estratégias de fuga são fundamentais para levar à vitória. Saber a hora de abandonar o gerenciamento para resistir ao gerenciamento alheio é, em resumo, a qualidade da estratégia geral do jogo.}
\end{itemize}

\subsection{O Jogo dos Outros}

A continuidade, levando em conta apenas as próprias cartas e as probabilidades absolutas, e não relativas, sobre que cartas podem determinar a sua próxima jogada, acaba contando pouco se não for complementada pelo conhecimento dos \textit{outros} jogadores: suas cartas e suas prováveis estratégias. Apenas escapar ao gerenciamento não funciona na maioria das vezes; gerenciar o jogo é fundamental.

Saber, diretamente, as cartas dos outros, é impossível (exceto quando alguém desafia um jogador que acabou de jogar um Curinga Comprar Quatro Cartas na mesa). Entretanto, sempre existem modos criativos de roubar no jogo --- e modos inteligentes de descobrir, indiretamente, as possíveis cartas dos seus adversários. Esses métodos serão explorados a partir de agora.

\subsubsection{Segurança}

Não vou, neste guia, aconselhar a desonestidade; além disso, explicá-la é trivial e tentar catalogar todos os métodos para ver as cartas dos adversários, impossível. O que se pode fazer é advertir os jogadores quanto a condições de jogo que favorecem a atitude desonesta:

\begin{description}
\item[Superfícies Espelhadas]{É preciso um cuidado especial quanto à escolha do lugar do jogo: espelhos e vidros podem proporcionar a certos jogadores posições privilegiadas em que podem ver as cartas de outros jogadores menos afortunados.}
\item[Interrupções Inocentes]{É natural que os jogadores sintam sede e fome, mas levantar no meio de uma partida para pegar água, quando para isto é necessário passar por detrás de vários jogadores, é suspeito: é possível instituir uma regra em que a partida é pausada para que o jogador possa fazer o que quer, ou, no mínimo, os jogadores devem ser espertos o suficiente para abaixar as cartas quando um deles passar por lugares específicos.}
\item[Postura e Relaxamento]{Uma coisa, contudo, é muito comum, especialmente em jogadores novatos e casuais: uma má postura dos braços ou do próprio corpo, combinada com uma má disposição de lugares para os jogadores, pode fazer com que, mesmo acidentalmente, outros jogadores vejam suas cartas. É preciso ter consciência disso e cuidado com o ângulo de visão dos adversários. Em geral, é possível prevenir isso se acostumando a deixar as cartas sempre na mesa, viradas para baixo, enquanto não for sua vez de jogar. Isso, contudo, pode ser problemático quando usado em conjunto com certas regras alternativas apresentadas mais adiante. Tome cuidado, também, com superfícies espelhadas e mesas transparentes, como foi dito anteriormente.}
\end{description}

\subsubsection{Leitura dos Jogadores: Psicologia}

Ao ``ler'' os jogadores e os ``sinais'' que eles mandam, é possível saber não que cartas exatamente possuem, mas determinar uma probabilidade de que sejam cartas boas ou ruins, ou de que a carta de destaque lhes agrade ou não.

Em geral, você não precisa ser um \textit{jogador} experiente para entender tais sinais: você precisa captar bem as pessoas e sua linguagem física. Isso pode ser aprendido, mas algumas pessoas têm esse talento de decifrar emoções e intenções observando as feições e atitudes das pessoas. O objetivo desse guia não é ensinar linguagem corporal a jogadores de UNO; se quiser aprender isso, deverá ler outros livros. Contudo, há algumas constantes e correlações a que devemos estar atentos:

Em geral, você deve primeiro determinar o tipo de jogador com o qual você está lidando. Algumas dessas perguntas podem ajudá-lo a desvendar aspectos essenciais da psicologia dos adversários:

\begin{itemize}
\item{Ele parece competitivo? Ou seja, a competição o estimula, de forma que ele queira sempre se esforçar para vencer? Ou isso lhe é, em última instância, irrelevante?}
\item{Ele parece habilidoso? Tem noção de estratégias de jogo, de tudo aquilo que ele pode fazer com as cartas que têm à mão? Ou joga de um modo quase arbitrário, aleatório?}
\item{Quando fica quieto e não propenso a risadas, falatórios e brincadeiras, seu jogo está ruim e sua estratégia, precária? Ou é apenas seu estado natural, seu estilo de jogo?}
\item{Por fim, tudo isso é o que ele aparenta ser ou o que ele é? Ele é capaz de dissimular emoções ou você pode extrair de seus sinais físicos imprudentes a verdade sobre sua situação no jogo?}
\end{itemize} 

Existe uma teoria acerca das \textit{5 Personalidades do UNO}: Elas descrevem cinco tipos de jogadores, e a partir deles podemos refletir sobre os adversários que encontraremos:

\begin{description}
\item[O Aldeão]{Este entra no jogo com a boca aberta e a mente fechada; sua mente não percebe as possibilidades à frente, e sua estranheza em relação aos procedimentos do jogo só é excedida por sua ingenuidade. Com relativa frequência, jogará cartas que não podem ser jogadas como se fossem válidas, e vai perguntar muitas vezes sobre as regras do jogo. Se fosse competitivo, provavelmente não estaria jogando, então ele provavelmente só joga por diversão sem grande esperança de vitória. Não é habilidoso, e costuma ser bastante transparente quanto à qualidade das cartas que tem à mão. Todas essas características estarão erradas, contudo, se ele for um \textit{Enganador}.}
\item[O Classe-baixa]{Este geralmente joga por motivos obscuros; não gosta do jogo e o considera barato e simples demais. Acaba sendo arrogante e violento na estratégia, o que pode tanto beneficiá-lo quanto (na maioria das vezes) ser sua perdição. Costuma criticar o jogo e outros jogadores gratuitamente, em especial quando faz algo que não é permitido no jogo, mas gostaria que fosse. Pode ser competitivo, mas vê a si mesmo como superior ao jogo e a qualquer derrota que venha a sofrer. Não necessariamente é habilidoso; em geral, apenas belicoso. É muito variável a relação entre seu humor temporário e a qualidade das suas cartas, uma vez que ele em geral vai estar mau-humorado pelo fato de estar jogando. Sua arrogância faz com que ele não tente enganar ou dissimular, mas, novamente, ele pode ser um \textit{Enganador}.}
\item[O Enganador]{Ele é o que o nome sugere: um enganador. É difícil defini-lo, mas em geral é competitivo (motivo pelo qual engana), habilidoso e, definitivamente, não se pode confiar em seus sinais físicos --- ele os cria e manipula a bel-prazer, a não ser que não seja habilidoso nisso e o faça muito mal, mas não costuma ser o caso. Em sua forma mais pura, ele é um jogador comum, embora inverta muitas coisas: fique animado com um jogo ruim, dando a impressão de estar preparando algo maior; fica triste com um jogo bom, parece nervoso mesmo com o jogo ganho, etc. Ele, entretanto, pode ir além de enganar suas ações e enganar mesmo sua essência: pode fingir ser um Aldeão (``Caramba, ganhei de novo! Acho que é sorte!''), um Falador (``Gente, vocês querem continuar jogando mesmo? As minhas cartas estão \textit{muito} boas\ldots'') ou um Classe-baixa (``Esse jogo é muito fácil, nem tem graça\ldots'') para atrair menos atenção dos adversários e ser um alvo menos prioritário. É preciso tomar um cuidado especial, contudo: ele pode ir ainda mais além e também roubar no jogo, procurando esconder cartas na roupa ou ver as cartas dos adversários.}
\item[O Falador]{Esse não para de falar. Vai reclamar do refrigerante, que está quente, ou de câimbras nas pernas, ou então que só tem cartas azuis nas mãos. Seu monólogo interno vai ser ininterruptamente transmitido aos outros jogadores. Suas características são muito variáveis, mas com o tempo fica fácil perceber se ele fala demais porque é um Enganador que não sabe disfarçar seus sinais físicos, e pretende enganar os outros com palavras, ou apenas alguém imprudente e não-competitivo que apenas gosta de falar durante o jogo e \textit{sobre} o próprio jogo.}
\item[O Sábio]{Aparenta ser calmo, contido, inteligente. O Sábio de tudo sabe, tudo vê. É a personalidade do mestre do UNO  e é aquela que deve ser buscada: competitivo a ponto de não ser imprudente; dissimulado quando necessário, mas acima de tudo impenetrável a ponto de, antes de passar aos adversários conceitos falsos, não passar conceito algum. Lembre-se de que você não deve se perguntar qual tipo você é: ao ler esta parte das estratégias do jogo, você busca por um conhecimento que leve à vitória. Mas o conhecimento das cartas por si só não fará isso: a vitória depende da sua atitude durante o jogo. Você não deve se perguntar qual tipo você é. Você deve ser este, o Sábio. Ou, no mínimo, o Enganador.}
\end{description}

Uma vez que você conhece seus oponentes (o que é mais fácil quando você joga com pessoas que conhece bem; amigos e familiares) você pode tentar decifrá-los --- se sua estratégia foi reforçada ou precisa agora ser reestruturada; se seu jogo é favorável e propício ou desarticulado --- em alguns momentos-chave:

\begin{enumerate}
\item{Ao comprar uma carta: como reagem?}
\item{Quando a cor da carta que eles jogaram na sua última vez muda durante a rodada: como eles lidam com a transformação? O que estavam esperando?}
\item{Uma mudança sutil, muito relativa e complexa: como eles reagem às reações à carta que jogaram na última rodada? Um jogador está tentando, em suma, \textbf{gerenciar} o jogo ou \textbf{escapar} ao gerenciamento? Por exemplo, ao jogar uma carta amarela, ele o fez por querer que um outro jogador comprasse uma carta ou por outras razões, que dizem mais respeito às cartas que ele mesmo possui?}
\end{enumerate}

Esse tipo de observações pode levá-lo a descobrir quais jogadores possuem boas cartas no momento --- ao gerenciar o jogo, você deve procurar atacá-los, impedindo que vençam. Ao entender quais jogadores estão procurando controlar o jogo, você pode querer evitá-los, ou mesmo enganá-los, fazendo-os acreditar que você não está bem no jogo --- ou fazendo-os acreditar nisso com atitudes como comprar uma carta, mesmo podendo jogar alguma da sua mão.

\subsubsection{Leitura dos Jogadores: Cartas}

Agora, deixando de lado os aspectos psicológicos dos jogadores e pensando na leitura fria e objetiva das possibilidades das \textit{cartas} que os adversários têm, chegamos a um pequeno conjunto de procedimentos e observações que podem nos ser úteis.

Em primeiro lugar, é sempre importante memorizar o momento em que os jogadores \textit{compram} cartas. Uma compra dissimulada é sempre arriscada (mais adiante), então aprenda a tomar as compras como ótimos indicadores das cartas do adversário. Se um adversário compra uma carta quanto, por exemplo, um 8 vermelho está na pilha de descarte, isso significa que, naquele momento, ele não possui oito algum, tampouco alguma carta vermelha. Também é relativamente seguro supor que a carta comprada não tenha sido 8, nem vermelha, caso contrário ela teria sido jogada pelo adversário. Isso é importante porque, à medida que o jogo prossegue, pode ser você o jogador a usar um curinga, e o adversário em questão pode ter uma ou duas cartas na mão. Se você souber que ele não possui nenhuma carta vermelha na mão, pode definir o curinga como vermelho e ele então não poderá jogar (a não ser que conserve um curinga na mão).

Esse tipo de memorização é crucial. Entretanto, quanto menos preocupados formos em ganhar, menos nos preocuparemos com detalhes: muitas vezes é difícil nos lembrar da cor que o adversário não tem, quanto mais dos números. Mas se você treinar e estiver interessado, poderá desenvolver uma melhor técnica para se lembrar desse tipo de detalhe.

Outro momento importante que pode passar despercebido é o momento em que um Curinga Comprar Quatro Cartas é jogado: ele (quase) equivale a uma compra, no sentido de que evidencia que \textbf{cor} de carta o adversário não possui --- a não ser que ele esteja blefando, é claro.

\subsubsection{Memorização das Cartas}

A informação sobre a cor de carta --- e os números ou símbolos --- que um adversário não possui é importantíssima, mas ela pode ser muito mais funcional se combinada com uma memorização completa de todas as cartas do jogo, analisando, assim, todas as possíveis mãos dos adversários e, à medida que o jogo progride, tentando descobrir especificamente quais cartas eles têm e não têm.

Como pode ser lido nas regras oficiais do jogo, o baralho de UNO possui 108 cartas. Memorize o número de cartas de cada tipo, e a cada carta decartada (e as que você sabe que estão na sua mão, ou as que você sabe que seu oponente não tem) você pode formar estatísticas: o que é mais \textit{provável} que seu adversário não tenha? Várias cartas vermelhas foram jogadas no descarte, mas poucas verdes. É mais provável que ele não tenha uma vermelha. E lembre-se de um detalhe que muitas pessoas não notam, mas pode fazer a diferença: todas as cartas numéricas possuem duplicatas (existem duas cartas 3 Verde, por exemplo) --- exceto as cartas de número \textit{0}.

\paragraph{Conclusão}

\begin{itemize}
\item{É essencial, na estratégia, balancear a continuidade e o tempo de gerenciamento de jogo (quando você gerencia e quando você não gerencia.}
\item{É importante não só pensar nas próprias cartas, mas nas dos adversários.}
\item{É importante prevenir contra desonestidades no jogo.}
\item{É importante poder ``ler'' os adversários.}
\item{E também positivo lembrar de certas cartas para poder fazer cálculos sobre o jogo.}
\end{itemize}

\subsection{O Estímulo Passional}

UNO é um jogo em que, quando jogado com leveza e por diversão, as emoções costumam vir à tona com a surpresa das reviravoltas inesperadas, a irritação do esquecimento da fugaz palavra ``UNO'', a frustração da falta de sorte com as cartas compradas\ldots

Quando o jogo não é composto majoritariamente por ``sábios'' (ver ``Personalidades do UNO'' acima) temos que as várias personalidades envolvidas no jogo tendem a agir passionalmente muitas vezes, a levar o jogo para o lado pessoal e a ter atitudes irracionais. Vamos procurar entender esta dinâmica, e como pode ser identificada, estimulada, e até certo ponto \textit{controlada}.

Como foi dito anteriormente, existem dois motivos para a ação consciente no jogo do UNO: o impulso de gerenciar o jogo e o impulso de escapar ao gerenciamento alheio; o foco no poder, no controle --- ou na velocidade e na sagacidade para se esgueirar pelas frestras das estratégias alheias e, com um pouco de sorte, dar o golpe final nos adversários.

Esses são impulsos \textit{conscientes}, atitudes de quem sabe o que está fazendo e como deve proceder para aumentar suas chances de vitória. Lado a lado com estas atitudes conscientes (exercidas em diversos níveis pelos mais variados tipos de jogadores; jogadores experientes preveem as jogadas e adivinham as cartas com mais clareza, mas iniciantes também possuem rudimentos de estratégia) caminham no jogador impulsos inconscientes, \textit{maneirismos ao jogar}, \textit{posturas de jogo} que podem, durante momentos de falta de lucidez tática, transbordar e se transformar em \textit{modo de jogar}, \textit{escolhas sobre o jogo}.

Existem, em geral, dois impulsos inconscientes: a agressividade e a tranquilidade. Há vários tipos de alvos contra o qual eles surgem: uma outra dupla (quando o jogo é jogado em parcerias), um jogador em especial ou o próprio azar (alvo que não deve ser ignorado).

Pode-se dizer que a agressividade, a vontade de atacar e reagir ao que o cérebro quase instantaneamente rotula como uma ``óbvia conspiração feita de ódio e ressentimento'' (Não pode ser! Comprar mais quatro, \textit{de novo}?) seria o lado sem foco, sem concentração e inteligência, do instinto de gerenciar o jogo, controlar a dinâmica da partida. Por outro lado, a tranquilidade, a sensação de que o jogo é conduzido de maneira equilibrada, seria a contrapartida instintiva e sem foco da tentativa de não controlar o jogo dos outros, mas o seu próprio, tentando fazer sua própria estratégia vitoriosa ao escapar da linha de fogo cruzado inimigo.

O foco, a atenção, a \textit{racionalidade} das escolhas separa os impulsos instintivos dos impulsos racionais. Os quatro impulsos podem levar à vitória de algum modo, afinal vencer neste jogo não passa somente por criar uma boa estratégia, mas esse é o único fator que se \textit{sabe} que ajuda e a medida dessa influência é visível e controlável. Contar com a sorte é possível, mas não creio que seja \textit{desejável} se você quer \textit{vencer}. E conhecer os impulsos irracionais faz parte da tentativa de incorporá-los à sua estratégia de jogo, seja ela no momento um impulso de gerenciamento ou não.

Ao passo que a tranquilidade é um impulso insípido no que concerne à nossa estratégia, é a agressividade que nos é útil --- entender seus efeitos sobre os jogadores é importante, pois um jogador sob o efeito da pura tranquilidade (que é levemente diferente do impulso racional correspondente) vai adotar uma estratégia que o conduz, calmo como um Buda, ao que lhe parece ser uma vitória --- ele apenas acompanha o fluxo que, caso não fosse interrompido pelas loucuras do jogo que todos bem conhecemos, levaria qualquer um, evidentemente, à vitória. De descarte em descarte, um dia se chega a ter nada na mão.

Já o jogador agressivo está sedento de sangue: ele pode até mesmo preferir sacrificar alguma oportunidade de melhorar seu jogo por qualquer chance que tenha de atacar seu agressor. Em outros casos, a raiva simplesmente o cega, impede que ele siga caminhos mais benéficos, que elevariam suas chances de vencer.

A agressividade tem origem em alguns pontos-chave do jogo:

\begin{description}
\item[Ficar muitas rodadas sem jogar]{Quando está prestes a jogar, o jogador anterior joga uma carta Pular, ou Inverter, e isso se repete fazendo com que ele fique muito tempo sem jogar. Em UNO, duas rodadas acaba significando uma terrível desvantagem --- muito tempo.}
\item[Ser obrigado a comprar cartas por não ter o que jogar]{Ter que comprar cartas por não ter na mão nenhuma carta passível de ser jogada. Note que isso é diferente de \textit{escolher} comprar cartas; a raiva surge da falta de escolha. A raiva é tanto maior quando a compra é infrutífera, e o número de cartas na mão aumenta.}
\item[Ser obrigado a comprar muitas cartas por causa de punições]{Esquecer de dizer ``UNO'' ou outro motivo que cause punições.}
\item[Ser obrigado a comprar cartas por obrigação]{Cartas de compra (Comprar Duas, Curinga Comprar Quatro) são verdadeiras causadoras de úlceras quando são jogadas repetidamente. A ideia de comprar 6 ou 8 cartas em duas ou três rodadas é um pesadelo para o jogador que não encara isso como uma oportunidade de gerenciar o jogo. Entretanto, quando há realmente pouca oportunidade de gerenciar o jogo, a raiva é mesmo muito difícil de controlar.}
\end{description}

Em alguns desses casos estímulos são necessários para que a situação desvantajosa gere uma raiva que possa obscurecer a mente do adversário ou fazer com que ele resolva atacar diretamente uma outra pessoa. Esse estímulo pode ser uma risada quando o incauto azarado comenta que vai mal, ou, mais fundamentalmente --- quando isso é possível --- comentários que indiquem a \textit{culpa} de alguém. Por exemplo, quando uma pessoa compra várias cartas em virtude de cartas Comprar Duas, é possível tripudiar da má fortuna do jogador e sugerir que o atacante que jogou essas cartas o fez de \textit{propósito}, por \textit{maldade}. Às vezes isso é verdade; às vezes o jogador até preferiria guardar as cartas, mas isso não seria possível sem ter que \textit{comprar} cartas, o que seria mais desvantajoso.

Os alvos, e os efeitos que causam, são:

\begin{description}
\item[Outro jogador ou outros jogadores]{O adversário em situação de infortúnio vai simplesmente tentar se vingar da melhor maneira que tem à disposição. Se ele possuir uma carta de Inversão e uma Comprar Duas Cartas, vai tentar fazer de tudo pra inverter o jogo e poder agora fazer seu algoz comprar as cartas --- e isso por vezes inclui comprar cartas desnecessariamente e fazer escolhas de jogo que até mesmo quebrem a própria continuidade.}
\item[O próprio azar]{Ao ter raiva do próprio azar, muitas vezes o jogador desiste de formar uma estratégia e simplesmente ``joga qualquer coisa'', descrente de que possa vencer o jogo. Desanimar o jogador, incutir nele a ideia de que a partida está perdida --- ou de que alguém é culpado pela sua compra de cartas, o que faz com que voltemos ao primeiro ítem desta lista --- é uma forma de se livrar, sem muito esforço, de um adversário que, sem tal desânimo, poderia tentar atrapalhar os outros sempre que possível.}
\end{description}

Para evitar se tornar alvo de jogadores vingativos, procure fazer o seguinte:

\begin{description}
\item[Não concentre ataques]{Tente não atacar um mesmo jogador seguidamente; dê no mínimo um espaço de uma rodada entre uma carta desvantajosa para ele e outra.}
\item[Não chame atenção]{Não faça um ataque declarado a um adversário ou comente suas intenções de fazer com que alguém compre mais cartas ou se atrapalhe no jogo, ou alguém experiente sequer precisará incitar um jogador contra você.}
\item[Justifique-se]{Há um velho ditado que diz que você nunca deve se explicar: seus inimigos não acreditarão, e quem for seu amigo não exigirá explicações. No UNO, declarar inocência pode funcionar. Dizer que era a única escolha (ou mesmo o mais verossímil ``era a única escolha sensata'') pode amenizar a situação.}
\item[Não blefe]{Blefar tem um desdobramento lógico perceptível para todos os jogadores: se ele mentiu, é porque \textit{queria} que o oponente comprasse as cartas. Isso chama a atenção e causa agressividade. Blefar, uma hora ou outra, torna-se importante, mas se você não quiser um comportamento agressivo contra você, prefira blefar raramente ou não blefar.}
\end{description}

É importante lembrar duas coisas: mesmo nas situações mais informais, alguém que tenta sempre jogar uma pessoa contra outra logo será reconhecida por tal característica, e mais adiante tudo que disser nesse sentido será ignorado. Portanto, como quase toda tática não-discreta contida neste guia, o uso repetido leva à ineficiência.

Além disso, é crucial lembrar que a agressividade mencionada aqui é puramente lúdica; UNO é apenas um jogo e as rivalidades devem durar apenas enquanto o jogo durar, devem ser encaradas como um aspecto natural e \textit{divertido} do jogo.

\paragraph{Conclusão}

\begin{itemize}
\item{Existem aspectos de um jogador focado --- gerenciamento e não-gerenciamento --- e impulsos menos racionais --- agressividade e tranquilidade.}
\item{A agressividade é útil pois pode enfraquecer os impulsos racionais no jogador. Ela pode ser incitada durante momentos particulares, contra um outro jogadores, outros jogadores ou o próprio azar.}
\item{É possível evitar ser um alvo corrente de agressividade.}
\end{itemize}

\subsection{Estratégias Específicas}

Existem certas estratégicas --- mais provavelmente \textit{táticas} --- que são empregadas durante um jogo oficial que vale a pena discutir e apresentar, ainda que não façam parte de um grande sistema de pensamento acerca da administração de uma partida de UNO.

\subsubsection{O Curinga Surpresa}

É muito comum querer guardar uma carta curinga para a última rodada. As vantagens são óbvias: se for possível chegar a um ponto em que você tenha apenas uma carta na mão, e essa seja um curinga, a única coisa que pode impedir você de ganhar o jogo é uma carta Comprar Dois (ou Curinga Comprar Quatro) do jogador anterior a você. Caso contrário, você pode vencer com qualquer carta na pilha de descarte.

O problema, entretanto, divide-se em dois problemas distintos: em primeiro lugar, você vai precisar de um pouco de sorte para chegar a ter apenas o curinga na mão. Como mostramos anteriormente, o seu poder de dar continuidade ao próprio jogo pode ser enorme quando você tem várias cartas, mas diminui progressivamente à medida em que você vai usando o seu poder, \textit{i.e.} descartando as cartas. Quando você chega a ter duas ou três cartas na mão, você acaba dependendo muito das jogadas dos outros jogadores; é preciso que a carta na pilha de descarte na sua vez seja perfeita para que você possa guardar o curinga até o final. Isso nem sempre acontece, e, quando não acontece, você é obrigado a comprar uma carta se quiser continuar esperando pelo momento oportuno.

Isso geralmente acaba mal, e aí temos nosso segundo problema: alguém com uma estratégia diferente vence o jogo antes de você e acaba ganhando 50 pontos por causa do seu curinga não descartado. Temos então que é uma estratégia bem arriscada, mas plausível e, quando bem executada, muito eficiente.

\subsubsection{``Olha o Avião'' Longo}

O cenário: você está com duas cartas na mão. Uma delas é o curinga; com o jogo do jeito que está, você precisa descartá-lo para seguir em frente, ou precisará comprar outra carta --- e outros jogadores já estão tão avançados quanto você no jogo. Você resolve não tentar ``O Curinga Surpresa'' (ver acima), e joga o curinga. Você possui, digamos, um 3 amarelo na mão. E agora você decide a cor.

Temos, antes de tudo, de ver um dilema anterior a essa tática: você deve evitar que os adversários ganhem; gerenciar o jogo. Dessa forma, se algum adversário já estiver com apenas uma carta na mão você deve exigir na sua carta curinga uma cor que você saiba que ele não tenha (ver ``Leitura dos Jogadores: Cartas'' acima).

Entretanto, se este não for o caso, você deve exigir a cor amarelo, na lógica da continuidade do jogo. Entretanto, dependendo do número de jogadores, esperar que na sua próxima vez a cor na pilha de descarte continue sendo amarela é, em geral, uma atitude que leva à frustração. Uma tática que parece uma aposta cega é exigir uma outra cor, não a que você tem na mão: exija azul, por exemplo. Quando a tática costuma funcionar? Quando você identifica um impulso gerenciador de jogo em algum dos seus adversários. Aquele que percebe que você pediu a cor azul, e que você, portanto, \textit{deve ter} uma carta azul na mão. Aquele que vai mudar a cor --- para, há uma grande chance, amarelo --- para impedir que você vença. Aquele que vai se irritar quando você revelar que o enganou.

Uma consequência curiosa do uso repetido e insistente dessa estratégia é que seus oponentes podem perder a fé na coerência das suas escolhas, e eventualmente prestarão menos atenção às suas atitudes, deixando de tentar ``ler'' a sua atitude e suas cartas --- porque julgarão isto uma tentativa inútil, já que você vai sempre pedir o que não tem. Contudo, essa estratégia, como mencionado, é arriscada: é mais provável que você perca mais do que ganhe ao usá-la repetidamente.

Essa é a variante ``Longa'' do ``Olha o Avião''. Existe uma variante que pode ser empregada em um jogo em parceiros, a ``Curta'', que veremos mais adiante.

\subsubsection{Jogue o Zero, Mantenha o Zero}

O zero é uma carta especial no jogo: existem menos zeros do que outros números no baralho. No final da partida, o zero não conta pontos para o adversário. A partir da primeira constatação, vemos que, teoricamente, descartar o zero é uma boa estratégia durante a partida: você vai ter uma melhor continuidade se conservar na mão cartas que te darão maior possibilidade de continuar jogando. Contudo, se um adversário ganhar, uma de suas cartas não vai lhe dar nenhum ponto.

O que fazer? Bem, é uma questão de momento. Se você tem mais esperanças de vencer o jogo, está numa situação favorável, não arrisque e livre-se do zero. Caso a vantagem seja do oponente, conserve-o.

\subsubsection{O Consumismo}

Comprar cartas é algo que, em teoria, deve ser evitado a todo custo. Afinal, o objetivo do jogo é se livrar das cartas à mão; você compra cartas quando não há carta que você consiga jogar no momento (o que parece uma irônica punição da justiça ``uno-versal'' a uma possível estratégia ruim) ou --- o que deixa óbvio quão ruim isso é --- quando é punido por algo como não dizer ``uno'' ou jogar errado, etc.

Entretanto, há momentos em que comprar uma carta pode ser uma boa estratégia: nos momentos em que você assume o gerenciamento do jogo, ou seja, pretende controlar, interferir nos outros jogadores.

Se a vez é sua e o próximo jogador tem apenas uma carta na mão --- e você não tem nenhuma carta que possa impedi-lo de jogar --- por que não arriscar e comprar uma carta? Você se distancia um pouco mais da vitória, mas caso você compre algo útil, também se afasta da derrota iminente\footnote{Os riscos de se fazer isso, tentando prever as chances de conseguir uma carta útil, podem ser calculados por cima (ou de forma precisa, dependendo da pessoa), como já visto na seção ``Memorização das Cartas'' acima.}.

Comprar cartas pode ser útil para se munir de cartas de ação, que ajudam a gerenciar o jogo, ou mesmo de cartas que melhorem a sua disponibilidade de cores à mão. Você pode fazer isso enquanto a partida ainda está ``fria'' --- quando uma carta a mais não é assim tão relevante.

Outro motivo para comprar cartas pode ser também a vontade de reservar uma carta para mais depois --- mais ou menos como ocorre com a estratégia do ``Curinga Surpresa'' mencionada acima --- mas, em contrapartida, não tem nada para jogar a não ser a carta que deseja guardar. A regra do jogo diz que uma compra é necessária; no mínimo você pode jogar a carta que comprou caso ela seja compatível com a carta da pilha de descarte no momento.

Assim, existem vários motivos para comprar cartas espontaneamente, não apenas por necessidade. Esteja alerta para estes motivos, que podem dizer muito sobre o jogo dos seus adversários --- não é porque eles compram uma carta que eles ``não têm a cor pra jogar''.

\subsubsection{Trébuchet Longo}

Em um jogo de muitas pessoas, o Trébuchet (como encontrado na internet) é o uso de uma situação desvantajosa em geral\footnote{Na verdade, alguns jogadores fazem uma distinção entre o Trébuchet, que é a situação de azar em que todos os jogadores não tem o que jogar e são obrigados a comprar, e um falso Trébuchet, em que um jogador pode jogar, mas não quer fazê-lo. Esse último caso é estrategicamente valioso, portanto é este de que vamos tratar agora.}, mas não exatamente para si mesmo, para gerenciar o jogo.

Ocorre que, em um número surpreendentemente considerável de partidas (sendo mais frequente com menos jogadores, como apenas 3 ou 4), vai haver um momento em que a última carta de uma cor e de um número ou tipo serão jogados; nenhum jogador vai ter em mãos alguma carta que possa ser jogada. Dessa forma, todos os jogadores comprarão cartas durante uma, duas, três ou até mesmo quatro rodadas até que finalmente alguém consiga comprar uma carta útil e possa jogá-la.

Ocorre que essa situação, tecnicamente desvantajosa para todos, acaba sendo mais desvantajosa para alguns do que para outros; pessoas que têm alguma estratégia, alguma esperança de continuidade (ainda aliada, talvez, à estratégia ``Olha o Avião'' acima) acabam se prejudicando mais do que outras que não teriam nenhuma jogada que não fosse guiada pela sorte caso o jogo continuasse. Se você é uma dessas pessoas, essa estratégia --- ou falta de uma --- é ideal.

Suponhamos que uma rodada já tenha se passado desde que aquele 4 vermelho foi jogado na pilha de descarte; ninguém consegue tirar ele dali. É uma partida de quatro jogadores, você vai comprar a segunda carta, iniciando a segunda rodada de compras --- e, coincidentemente, você compra um quatro azul. Mas os outros jogadores, você percebe, estão em melhor situação que você. Admita, você não tem uma boa estratégia formada. Então, por que ajudar os outros jogadores? Por que dar a eles a chance de continuar, retomar suas estratégias? Ignore que você comprou uma carta que possa ser jogada. Continue comprando --- é claro, contanto que os outros também comprem. Duas ou três rodadas depois todos vão estar em pé de igualdade novamente, e você pode refazer sua estratégia, com cartas que abram o seu leque de possibilidades.

Esta é a variante longa do Trébuchet; existe uma variação ``Curta'' (para jogos apenas com dois jogadores) e uma variante ``Dupla'' (para jogos com parceiros), que serão abordadas mais adiante.

\subsubsection{Pontos Extras}

Uma estratégia interessante consiste em deixar para a última carta a ser descartada um Comprar Dois (ou um Curinga Comprar Quatro, o que de certa forma caracteriza um Curinga Surpresa, mas esta estratégia é um pouco diferente), pois, de acordo com as regras oficiais, se uma vitória for alcançada através de uma carta que force o próximo jogador a comprar, o próximo jogador em questão \emph{deve} comprar as cartas para que elas possam ser contadas no final da partida.

\subsection{Jogo em Dupla}

Como visto anteriormente nas regras oficiais, o jogo em dupla conta com regras levemente diferentes. Afinal, as cartas ``Pular'' e ``Inverter'' não fariam sentido num jogo de duas pessoas. Nesse caso, essas cartas recebem a nova função de fazer com que o jogador que as joguem apenas jogue novamente.

A estratégia do jogo torna-se um pouco diferente: o foco na continuidade agora é ainda maior, levando em conta que não é tão provável que a cor da carta mude entre duas jogadas suas (em que apenas uma pessoa joga). Por outro lado, o fato de que as cartas especiais provocam a repetição do seu turno nos deixam com três tipos básicos de estratégia: uma assertiva, uma oportuna e uma flexível, e as três (especialmente as duas primeiras), apesar de mutuamente exclusivas, podem ser combinadas ao longo de uma partida (enquanto reação a uma estratégia bem sucedida do adversário). Além disso, ainda possuímos uma equivalente à ``Trébuchet'' analisada na seção anterior.

\subsubsection{Assertividade e Oportunismo}

Considerando que as cartas Pular, Inverter, Comprar Duas Cartas e Curinga Comprar Quatro Cartas fazem com que você jogue novamente, você deve ter em mente que qualquer quantidade desse tipo de cartas na sua mão faz toda a diferença do mundo, especialmente se elas puderem ser combinadas (i.e tiverem a mesma cor ou forem do mesmo tipo). Assim, ter três cartas Pular de cores diferentes é certamente melhor do que ter uma carta Pular, uma carta Inverter e uma carta Comprar Duas Cartas, as três de cores diferentes --- isso porque, quando aquela primeira série de cartas é jogada, é jogada em conjunto, uma atrás da outra, e faz com que você fique três rodadas à frente do jogador.

Contudo, a questão é: como e quando fazer uso dessas cartas que te deixam à frente? Já no começo da partida, para ganhar vantagem o mais cedo possível e liquidar o jogo? Mais perto do fim, arquitetando assim um ataque fulminante para qual o adversário não está preparado? Ou nenhum dos dois?

A primeira opção é a da Assertividade. Pode parecer óbvia a vantagem de começar o jogo já em vantagem; afinal, você sai na frente antes que os percalços do jogo façam com que você acabe comprando cartas e fique uma ou duas rodadas atrás do adversário --- o que pode ser fatal caso ele esteja se preparando para atacar de forma oportuna, o que veremos adiante.

O ataque assertivo possui uma função de desarmar o jogador adversário; se ele possui alguma combinação de cartas que possa dar esse tipo de vantagem a ele, ele vai ter que usar já para recuperar as rodadas que você já tomou dele; dessa forma, ele não pode mais guardar nenhuma carta na manga.

Guardar cartas na manga é o espírito do ataque oportuno. Construir as jogadas, basear-se na continuidade, até que a situação seja propícia e você despeje em sequência todas as cartas que façam você jogar novamente, ganhando, assim, o jogo, mesmo tendo três cartas na mão, por exemplo.

É claro que tanto a estratégia assertiva quanto a oportuna têm seus pontos fracos. Se a estratégia assertiva força o oponente a ``sair da toca'', você pode não gostar do que vai sair de lá; ele pode ter cartas muito mais vantajosas, e virar o jogo de um jeito espantoso, fazendo com que você fique numa situação muito pior do que se estivesse esperado, talvez, para realizar um ataque oportuno (embora se as cartas do adversário forem realmente boas talvez não se pudesse fazer muito quanto a isso).

Por outro lado, a estratégia oportuna funciona muito bem com um adversário passivo e que não tem um jogo bom. Você precisa de uma ótima continuidade de jogo para prosseguir constantemente rumo à utilização das cartas de ação que possam levá-lo à vitória. Um ataque oportuno do seu adversário, dado antes do seu, fará você perder. Um ataque assertivo vai forçar você a abandonar o ataque oportuno --- ou não, embora isso seja aconselhável.

Tanto o ataque assertivo quanto o oportuno se referem às cartas de ação jogadas como um contínuo; mas nem sempre essa jogada é a melhor saída (ou mesmo é possível. Nesse caso, precisamos falar de uma terceira estratégia.

\subsubsection{Flexibilidade}

Por vezes, as suas cartas de ação não podem ser jogadas uma em cima da outra (proporcionando, assim, várias rodadas de vantagem de uma vez só), ou simplesmente não é vantajoso fazer isso. Nesse caso, adotamos uma estratégia flexível.

Por vezes esta estratégia não é bem uma tática, e sim a falta de uma; uma pessoa com pouca habilidade não percebe direito as consequências de ataques assertivos e oportunos, e apenas joga as cartas de ação quando julga necessário --- mas não necessariamente em conjunto, e sim de uma em uma, usando-as com moderação. Inadvertidamente ou de propósito, esse pode acabar sendo, justamente por seu equilíbrio e sua prudência, o melhor tipo de estratégia.

É preciso reconhecer o momento certo de atacar para obter pequenas vantagens por dois motivos distintos: Ou o adversário já está em vantagem, e é preciso dirimi-la ou diminui-la; ou ele não está em vantagem, mas obter uma vantagem grande (de duas ou três rodadas) fará com que ele fique desesperado (como comentado acima) e ataque você também. Dessa forma, conquistar pequenas vantagens, carta a carta, é uma estratégia que evita o susto e o jogo baseado no medo; você não pode deixá-lo pensar que as coisas estão ficando ruins demais e, portanto, escapando a seu controle (incentivando-o, assim, quem sabe, a continuar com sua estratégia de ataque oportuno quando deveria imediatamente abandoná-la).

Lembrando que o jogo baseado na flexibilidade também oferece maiores oportunidades de continuidade (o que é um aspecto muito positivo) e, por vezes, é o único possível.

\subsection{Estratégias Específicas do Jogo em Dupla}

Em um jogo entre duas pessoas, podemos dizer que existe uma estratégia específica a ser debatida:

\subsubsection{Trébuchet Curto}

Assim como no jogo de 3 ou 4 pessoas, você pode chegar a uma situação muito desvantajosa para ambos os jogadores ao jogar em dupla. O Trébuchet é a estratégia em que você, mesmo podendo jogar uma carta e fazer o jogo continuar, prefere continuar comprando (contanto que o seu adversário também esteja comprando) para refazer o jogo.

Isso, é claro, não faz nenhum sentido a não ser que você esteja em situação de desvantagem. A compra de cartas durante duas, três ou quatro rodadas pode fazer uma diferença gigantesca --- tanto para o seu bem quanto para o seu mal.

Esta é a variante ``Curta'' do Trébuchet. Para a variante Longa (jogos com vários jogadores), veja ``Trébuchet Longo'', acima. Para a variante ``Dupla'' (jogos com parceiros), veja mais adiante.

\subsection{Jogo em Parceiros}

As regras oficiais do UNO são particularmente displicentes quanto ao jogo em parceiros: ela apenas explicita que os parceiros devem sentar-se um à frente do outro. Mas o que é permitido que façam?

Essa ambiguidade não foi referida na seção que seria teoricamente apropriada pois não é uma ambiguidade entre duas coisas: a questão é muito aberta para ser posta em termos de fazer alguma coisa ou outra coisa. A questão é que não se sabe como funciona um jogo em parceiros; eles podem mostrar suas cartas uns aos outros, ou só podem fazê-lo clandestinamente, correndo o risco de fazer com que os adversários vejam as cartas? Podem fazer sinais e pedir por cartas e cores abertamente, ou isso é considerado um ``palpite'' (acarretando, assim, penalidades)?

Como um jogo em parceria não faria sentido se as equipes fossem descoordenadas, vamos considerar que a comunicação entre elas é aberta. Elas podem, portanto, inventar símbolos e formas secretas de comunicação para melhor coordenar suas ações --- afinal, de que adianta jogar em equipes se isso não acaba constituindo uma união tática de forças, mas sim apenas uma contagem conjunta de pontos?

Contudo, ver as cartas --- assim, de forma fácil, o tempo todo --- pode ser um pouco exagerado. Formar uma estratégia, gerenciar o jogo, pode ser, se assim for feito, fácil demais; além disso, a dinâmica do jogo pode ser um pouco prejudicada se o tempo inteiro os jogadores ficarem olhando as mãos uns dos outros para pensar em um jeito de atacar os adversários --- de forma implícita, o tempo (e o fato de que você deve jogar rápido para deixar o jogo interessante) faz parte da dificuldade em formar e executar uma estratégia.

Então, partindo do princípio de que esse é um manual de customização, crie suas próprias regras; eu, particularmente, recomendo que não sejam aplicadas penalidades para a comunicação entre membros da equipe, mas essa comunicação não deve ser constante e absoluta.

\subsubsection{Dupla Continuidade}

Se antes o importante era jogar de modo que a sua próxima vez seja vantajosa para sua estratégia geral, agora você tem uma outra variável: para que a sua parceria seja vitoriosa (contra uma outra), vocês precisam de uma única estratégia, utilizando-se do poder das duas ``mãos'' e jogadas --- duas pessoas em sintonia. A continuidade passa a ter outro sentido: o sentido de ser posta em ação em prol da estratégia geral, não apenas da própria.

Considerando, pois, um jogo de duas duplas, em que os jogadores A e C jogam contra os jogadores B e D\footnote{vamos considerar essas duplas com letras de ordem alternada, uma vez que as duplas não podem ficar em posições consecutivas na mesa.}, vamos imaginar um jogo-exemplo para darmos uma melhor ideia dessa questão da dupla continuidade (a continuidade que serve não apenas a si mesmo, mas à dupla).

\subsubsection{Escolha das Cartas em Virtude da Dupla Continuidade}

Pensar em escolhas de cartas que favoreceram a uma estratégia geral parece complexo se você não sabe qual é essa estratégia: tendo em vista que o UNO não é explícito quanto ao modo de jogar um jogo de dupla\footnote{Ver seção ``Jogo em Parceiros'' acima}, vamos adotar um procedimento cauteloso em que é possível se comunicar, de forma secretiva, com seu parceiro --- mas não de forma aberta, excluindo, por exemplo, trocas de mãos para verificação das cartas.

Assim sendo, é impossível chegar a esse tipo de consenso explícito quanto a o que é melhor fazer. Contudo, a própria dinâmica de jogo pode acabar por levar a uma situação em que seja claro que um dos jogadores possui uma melhor chance de vitória: nesse momento, é crucial entender como a parceria pode fazer a diferença. No UNO jogado apenas em um, você deve sempre escolher entre ter uma atitude que gerencia o jogo ou uma que foge ao (potencial) gerenciamento alheio. Você deve escolher entre tentar manipular ou simplesmente correr para a vitória, com o risco de, ao fazer isso, baixar a guarda e algum outro jogador te atacar, evitando sua vitória. Com o jogo em parceria, entretanto, você \emph{não tem} que fazer essa escolha; se um dos jogadores estiver com, por exemplo, 3 cartas, e o outro com 10, não faz sentido para o jogador com 10 jogar para tentar livrar-se de suas cartas, ou jogar tentando dar continudade ao próprio jogo. Deve, acima de tudo, tentar atrapalhar os adversários ou (se possível, simultaneamente) ajudar o companheiro a se livrar de suas cartas, facilitando a vitória da dupla.

É claro que, nesse contexto, ter um jogador que acaba acumulando, a um certo tempo de jogo, várias cartas especiais é um tanto perigoso; um ataque surpresa (vindo de uma estratégia como o Isolamento ou o Isolamento Fatal, vistos mais adiante) poderia dar à equipe vencedora um número absurdamente alto de pontos em apenas uma jogada. Isso, entretanto, é apenas uma meia-verdade que não resiste à investigação: um jogador acumula cartas para que o outro possa livrar-se dela. De nada adiantaria uma política de controle equilibrado: ambos os jogadores teriam cartas especiais, de 20 ou 50 pontos, e acabariam encrencados da mesma forma após uma derrota inesperada e surpreendente.

De qualquer forma, podemos identificar quatro dinâmicas de duplas:

\begin{description}
\item[Falta de Foco]{Quando não há uma estratégia definida, e ambos os jogadores tentam gerenciar o jogo: se preocupam com o jogo alheio e, sem objetivos claros, não conseguem formar um ataque sólido, embora consigam, até certo ponto, atrapalhar os adversários. Essa é, portanto, uma dinâmica, em geral, \emph{negativa}.}
\item[Descoordenação]{Quando não há uma estratégia definida, mas ambos os jogadores desistem de gerenciar o jogo e se preocupam em tentar vencer, dando continuidade ao próprio jogo. Embora isso possa parecer suficiente, uma equipe descoordenada pode virar presa fácil para uma coordenada --- e a tentativa de dar continuidade apenas ao próprio jogo pode acabar atrapalhando o próprio outro membro da equipe. Essa é, portanto, uma dinâmica, em geral, \emph{negativa}.}
\item[Divisão de Tarefas]{Quando há uma estratégia implícita; um jogador se responsabiliza por tentar, de todas as maneiras, gerenciar o jogo --- algo como um ``sacrifício'', enquanto o outro deve comprometer-se unicamente em ganhar o jogo o mais rapidamente possível, livrando-se das cartas à mão, contando com o máximo de apoio do companheiro. Essa é, portanto, uma dinâmica, em geral, \emph{positiva}.}
\item[Neutralidade]{Quando não há estratégia delineada, mas há coordenação na dupla: os jogadores possuem a sensibilidade --- e são capazes de se comunicar de forma eficiente, ainda que minimalista --- de identificar quem tem mais poder de gerenciamento de jogo e assim conseguem, em uma ou duas rodadas, formar uma estratégia. Essa é, portanto, uma dinâmica, em geral, \emph{positiva}.}
\end{description}

Voltamos, portanto, ao jogo-exemplo mencionado no final da seção anterior. Vamos acompanhar a dupla dos jogadores A e C. É início de jogo. O jogador A possui em sua mão as cartas:

\vspace{0.5cm}

\textbf{Inverter Verde \hspace{1cm} 9 Vermelho \hspace{1cm} 7 Azul \hspace{1cm} 5 Azul \hspace{1cm} Comprar Duas Cartas Azul \hspace{1cm} 7 Vermelho \hspace{1cm} Curinga}

\vspace{0.5cm}

O jogador C possui em sua mão as cartas:

\vspace{0.5cm}

\textbf{3 Verde \hspace{1cm} 2 Vermelho \hspace{1cm} 2 Azul \hspace{1cm} 3 Amarelo \hspace{1cm} 3 Amarelo \hspace{1cm} 7 Verde \hspace{1cm} 3 Vermelho}

\vspace{0.5cm}

A primeira coisa a se perguntar é: é possível, desde já, formar alguma estratégia, e evitar acabar jogando dentro de uma dinâmica de duplas negativa? Há sim algo a fazer, desde que se saiba de antemão como identificar quem possui um jogo mais capaz de \emph{gerenciamento}.

Dos dois jogos, o primeiro é o mais capaz disso: o jogador A possui um curinga, uma carta Comprar Duas Cartas e um Inverter, de modo que ele pode prestar atenção ao jogo e coordená-lo de acordo com sua vontade (até certo ponto). O jogador C, contudo, não tem todo esse poder bélico. Se o jogador A for capaz de comunicar que seu jogo é um jogo bom em termos de gerenciamento, então está formada uma estratégia de Divisão de Tarefas.

Caso nenhum dos dois jogos se mostrasse obviamente superior em termos de gerenciamento (seja porque ambos os jogos são ruins para tal ou bons), então deve-se esperar uma rodada, duas, ou mais, para ver quem deverá correr e quem deverá dar cobertura.

O jogador A começa o jogo. Suponha que ele saiba que ele vai ser o parceiro a gerenciar o jogo, facilitando as jogadas do jogador C. A carta na pilha de descarte é o 2 azul; você deve jogar o 5 azul, o 7 azul ou o Comprar Duas Cartas Azul. Como você deve se lembrar, jogar uma carta de compra logo no começo não é a estratégia mais recomendada, e uma vez que o jogador A já possui um outro 7, ele pode jogar o 7 Azul para auxiliar a continuidade do seu próprio jogo --- afinal, não é porque você deve priorizar a continuidade da dupla que você deve sempre esquecer da sua. Considere também que o jogo está apenas começando.

É a vez do jogador C, e a carta na pilha é um 5 azul. Não há outra opção, se ele quiser continuar jogando (e lembre-se que esse deve ser o objetivo dele: continuar jogando, inexoravelmente, até a vitória --- pelo menos enquanto a estratégia for esta). Ele joga sua única carta azul (um 2).

Como estamos considerando um exemplo de \emph{boa estratégia}, vamos considerar que o jogador C consiga emitir um sinal para o jogador A (e se fazer entender) demonstrando que esta é sua última carta azul. Lembre-se: o jogador A é responsável por tentar fazer, o quanto puder, que o jogador C continue jogando. Esta é, portanto, uma informação muito valiosa --- e seria péssimo que esse sinal fosse interceptado pelos adversários.

A carta agora é um 2 Amarelo, e é a vez do jogador A. Isso é bom; não é uma carta azul. Entretanto, ele não tem o que jogar, a não ser um Curinga. Como o Curinga pode ser muito mais útil em um estágio mais avançado do jogo, ele compra uma carta. Que mal pode haver? Será uma oportunidade de, quem sabe, ganhar mais uma carta de ataque, que pode ser mais útil.Infelizmente, é apenas um 8 vermelho.

Nossa dupla é surpreendida: o jogador C é pulado com uma carta Pular Amarela. Depois que o último adversário jogou, é a vez do jogador A novamente. Há um 9 amarelo na pilha de descarte. Você joga seu 9 vermelho, na esperança de, quem sabe, atacar algum plano do jogador B.

O jogador B joga um 9 Amarelo e o jogador C pode jogar: a estratégia continua. Ele joga um 3 Amarelo (uma vez que até possui outro). O jogador D inverte o jogo --- sorte, uma estratégia inadvertidamente prejudicial ou falta de coordenação: qualquer uma das coisas se sucede e o jogador C jova mais uma vez. Outro 3 amarelo e agora as coisas começam a se complicar; o jogador C faz um sinal para o jogador A: as cartas amarelas acabaram.

O jogador B não consegue jogar e agora é a vez do jogador A, que deve mudar a cor da carta para tentar garantir que o jogador C continue jogando. O jogador A joga o Curinga e pede pela cor verde. O Jogador D joga um 2 verde, e o jogador C pode então jogar um 3 Verde.

Podemos, antes de analisar as próximas jogadas, fazer um balanço do jogo e confirmar a importâncida coordenação da equipe. Se a divisão de tarefas não fosse empregada (e se não houvesse ocorrido uma excelente comunicação entre parceiros), o jogador A poderia fazer diversas jogadas visando a própria continudade, o que atrapalharia o jogador C --- ou ter deixado de jogar para ajudar o próprio parceiro, o que teria atrasado mais a dupla. Nesse momento, o jogador A tem 5 cartas, e o jogador C, 3. O jogador B, prestes a jogar, tem 5 cartas, e o jogador D, outras 3.

O jogador B pula o jogador A; o jogador D pula o jogador C e o jogador B usa mais uma carta Pular (foi uma sequência de cartas Pular; uma verde, uma vermelha e uma azul) para pular o jogador A e agora o turno é do jogador D novamente. Ele joga uma carta Inverter e fica com apenas uma carta na mão (não esquecendo de declarar-se ``UNO'').

O jogador A se vê numa situação complicada: ele é o responsável por gerenciar o jogo, e deve fazer algo quanto a isso. O próximo a jogar é o jogador B. Ele pode inverter o jogo, mas então o próximo a jogar seria o jogador D --- justamente aquele que possui apenas uma carta. Seria bom se o jogador A pudesse saber qual é a carta do jogador D, pois se ela não fosse verde (ou um Inverter) ele forçaria o jogador D a comprar uma carta. Entretanto, ao não prestar atenção no jogo (contando cartas) e ao não ter ocorrido um momento em que o jogador D se mostrou carente de cartas verdes (o que incorreria numa possibilidade alta de sucesso no ataque), não há como saber disso. É arriscado demais voltar o jogo para o jogador D.

Em compensação, se ele deixar o jogo correr, o jogador C não vai possuir nenhuma carta que possa prejudicar o jogador D.

Há outras duas cartas que possam ser jogadas: um 5 Azul --- inócuo nesse cenário --- e um Comprar Duas Cartas Azul. A questão é que jogar um Comprar Duas Cartas Azul faria com que o jogador \emph{B} --- e não o \emph{D} --- comprasse duas cartas, o que não seria exatamente o ideal. Além disso, a carta no topo da pilha de descarte seria azul, e o jogador C não possui mais cartas azuis.

Entretanto, isso, vai forçar o jogador C a comprar uma carta, e quem sabe essa carta não será uma carta que possa fazer com que o jogador D perca? Embora não tenha contado, com precisão, as cartas do jogo, o jogador A percebeu que \emph{nenhuma carta Curinga Comprar Quatro Cartas ou Comprar Duas Cartas, de qualquer cor, foi usada}. Dessa forma, há mais chance de que essas cartas ainda estejam no monte de compras. O Jogador A, portanto, decide por fazer com que o jogador B compre duas cartas.

É a vez do jogador C e, infelizmente, ele compra um 4 Vermelho. Com sorte, o jogador D não consegue vencer o jogo com a carta que tem à mão (note que na pilha de descarte está o Comprar Duas Cartas Azul). O jogo continua.

Para aqueles que possam pensar que de nada adianta a técnica se a dupla quase perdeu, lembre-se que o UNO também incorpora o fator da sorte e do acaso. E, ademais, estamos acompanhando apenas a trajetória de uma das duplas, não sabendo que relevância teve a estratégia da \emph{outra} dupla na configuração deste cenário.

O jogador D compra uma carta, entretanto, que ele pode jogar: um 4 Azul. Mais uma vez ele possui apenas uma carta na mão. Sabendo que ela não é azul, há agora mais chances de ela ser verde: jogar o Inverter Verde tornou-se ainda mais arriscado (sem falar que, no momento presente, impossível devido ao 4 Azul na pilha de descarte).

Declaradamente, a estratégia não mudou: o jogador A continua a tentar gerenciar o jogo. Jogar uo 5 Azul seria inútil, além de não mudar a cor da carta (atitude preferível, uma vez que o jogador C continua sem cartas azuis). O jogador A resolve, portanto, tentar a sorte mais uma vez: ele compra uma carta. É um Comprar Duas Cartas Amarelo que, infelizmente, não pode ser jogado agora.

O jogador B joga um 8 Azul. O jogador C compra um 4 Amarelo; não pode jogar. Mais uma vez, o jogador D não consegue ganhar e compra uma outra carta, que ele joga: é um Curinga. Ele pede a cor Vermelha.

O jogador A se vê em maus lençóis: as únicas cartas vermelhas que tem são inofensivas, e ele não pode jogar nada diferente frente ao Curinga na pilha de descarte. Ele compra mais uma carta: um 1 Vermelho. Não é o suficiente.

O jogador B joga um 8 Vermelho. Com a situação crítica, é mais do que sensato abandonar a divisão de tarefas e tentar gerenciar o jogo: é claro que o jogador C possui cartas vermelhas para jogar (3 delas), mas ele resolve comprar mais uma carta no que será a derradeira tentativa de impedir a vitória da outra dupla: ele compra um 3 verde, que não pode ser jogado.

O jogador D joga, enfim, seu 9 Vermelho, garantindo a vitória da dupla B e D. Como lição deste jogo-exemplo, vemos que embora a estratégia geralmente decida o jogo, a sorte também é importante. Por outro lado, não acompanhamos a dinâmica de dupla da dupla vencedora, que pode ter sido simplesmente superior. Além disso, em um jogo de parceiros existem vários outros sinais, e o fato de que alguns sinais podem ser interceptados --- ou, pior, interceptados, mas mal-interpretados --- adicionam ao jogo um fator de aleatoriedade e uma gama de possibilidades impossível de replicar em jogos-exemplo como esse.

\subsection{Sinais}

Já discutimos, no jogo-exemplo anterior, que tipos de sinais podem ser úteis aos parceiros durante o jogo. Vamos discutir mais em profundidade estes sinais nesta seção.

Existem uma característica fundamental de todo sinal: a profundidade de detalhes com a qual relata para o outro jogador o próprio jogo. Portanto, um sinal extremamente raso indica apenas se um jogo está ``bom'' ou ``ruim''; um sinal extremamente profundo diz ao parceiro exatamente quais são as suas cartas. Sinais rasos são difíceis de interpretar (por parte dos adversários ou, se não houver combinação prévia, por parte do próprio parceiro), ao passo que quanto mais profundo um sinal, mais difícil será de fazê-lo de forma discreta, ou de criptografá-lo a ponto de criar dificuldades para que os adversários o entendam. Alguns sinais, contudo, não precisam ser explícitos: podem ser mascarados como cacoetes, como coçar a bochecha ou a cabeça, e são mais eficientemente mascarados se você não olha para o seu parceiro enquanto os executa (embora isso exija um nível de cumplicidade enorme: você não terá certeza se o parceiro entendeu o sinal ou mesmo se o viu).

Seja como for que o sinal seja feito, é útil saber que tipo de informação, não muito rasa, mas também não muito profunda a ponto de necessitar de sinais complexos, pode ser útil para o parceiro em um jogo de duplas:

\begin{description}
\item[Número de Cartas de Ataque]{Quantas cartas de ataque (Inverter, Pular, Comprar Duas, Comprar Quatro) você possui (ou não possui, caso não possua nenhuma). Pode ser útil para definir quem gerencia o jogo.}
\item[Número de Cartas Pular]{Quantas cartas Pular você possui. Útil para uma eventual formação de um Isolamento (ver adiante).}
\item[Comprei uma Carta de Ataque]{Informe ao parceiro que você comprou algo interessante. Se possível, diga que tipo de carta, ou apenas diferencie entre Pular/Inverter e alguma carta que force o adversário a comprar cartas --- ou diga se é um Curinga ou não.}
\item[Sem esta Cor]{Indicar, de alguma forma, quando você não tem mais uma cor em sua mão. Algo como uma coçada na nuca quando a cor estiver na pilha de descarte já funciona.}
\item[Que Cor?]{Quando o parceiro puder escolher uma cor (por meio de um Curinga), e você tiver alguma preferência, pode indicar que cor prefere.}
\item[Mude a Cor]{Se por algum motivo (não necessariamente por ser a cor que você não tem) seja uma boa ideia mudar a cor na pilha de descarte, você pode informar ao parceiro.}
\end{description}

\subsection{Estratégias Específicas de Jogo de Parceiros}

A maioria das estratégias específicas, empregadas em um jogo individual, pode ser traduzida em estratagemas análogos, porém utilizando todas as possibilidades do jogo em duplas.

\subsubsection{O Curinga Facilitado}

Quando o Curinga Surpresa (ver acima) é empregado, um dos jogadores precisa contar com certa sorte: precisa jogar de modo a fazer com que sua continuidade dê certo, sem poder controlar o jogo, guardando o Curinga para que ele seja a última carta.

Num jogo de duplas, contudo, quando uma equipe possui alguém que claramente é o gerenciador do jogo e o outro que é aquele que deve se livrar das cartas (e este último possui um Curinga), essa estratégia passa a precisar menos de sorte, e mais de companheirismo e perspicácia. No Curinga Facilitado, o jogador gerenciador controla o jogo de modo a permitir que o outro jogador da mesma equipe conserve seu curinga e jogue-o na última rodada. 

\subsubsection{``Olha o Avião'' Curto}

O Olha o Avião longo, empregado no jogo individual, pode ser empregado também em um jogo de duplas em uma situação muito particular.

Imaginemos novamente um jogo-exemplo, com os jogadores A, B, C e D (duplas A e C, B e D). O jogador A escolhe a cor para a próxima jogada. O jogador B é o próximo a jogar. O jogador C, da equipe do jogador A, pode ganhar o jogo (ou tem duas, três cartas na mão).

Se for provável que o jogador B seja o jogador gerenciador na dupla adversária --- ou que, naquele momento, esteja, de alguma forma, claro que ele tentará inverter a cor para que o jogador C não consiga jogar, uma boa estratégia é apostar na requisição de uma cor estranha, que mesmo o jogador C não tenha (ou que não seja a melhor escolha possível). O jogador B inverte as cores, escolhendo uma que talvez seja justamente a melhor cor para a situação.

\subsubsection{Zero em Duplas}

O mesmo conselho pode ser dado quanto ao zero, tanto em jogos individuais como em duplas: conserve se o momento não for favorável, descarte o mais rápido possível em nome da continuidade. Entretanto, em um jogo de parcerias, se um dos jogadores estiver em posição de gerenciar o jogo, e o zero estiver nas mãos do outro parceiro, então a situação não é tão ameaçadora, e o zero pode ser conservado por mais tempo. 

\subsubsection{Trébuchet Duplo}

Assim como os tipos de Trébuchet vistos anteriormente, essa estratégia, que na mairoia das vezes surge de forma acidental, também pode ocorrer em um jogo de duplas. É possível que ninguém possua uma carta adequada para jogar, e então todos comprem durante duas ou três rodadas, o que pode ser muito vantajoso para uma dupla, embora não para outra (embora nem sempre este seja o caso).

No Trébuchet Duplo, contudo, vale a pena destacar algumas diferenças: Em primeiro lugar, é preciso que a dupla esteja unida nessa estratégia; se um dos jogadores a vê e a utiliza, é preciso que seu parceiro também colabore, se não de nada adiantará o esforço individual do primeiro.

Além disso, o fato de que as jogadas dos dois jogadores de uma dupla forcem os outros jogadores a comprar não constitui um Trébuchet; um Trébuchet é quando \emph{todos} compram. Estratégias deste outro tipo serão abordados de agora em adiante.

\subsubsection{Abastecimento}

A estratégia do abastecimento pode ser empregada mediante uma série de fatores e é difícil de funcionar bem por um longo tempo. Em geral, é antes um princípio de ação do que uma estratégia, uma vez que se pode agir de modo a tentar ``abastecer'' um adversário, mas dificilmente se poderá manter essa dinâmica de jogo sem prejuízo à estratégia maior da dupla, que é, evidentemente, a vitória.

O abastecimento consiste em jogadas feitas com o objetivo de forçar um adversário (ou até mesmo os dois, embora isto seja muito mais complexo) a comprar cartas. O alvo preferencial será escolhido de acordo com seu número de cartas (quanto menor, mais importante é que seja abastecido) e também de acordo com a característica das suas cartas no momento.

O abastecimento envolve, portanto, a coordenação da dupla em prol de jogadas que forcem um jogador a comprar cartas. Vamos voltar ao modelo de jogo exemplo com o qual temos trabalhado: duplas entre os jogadores A e C, e entre os jogadores B e D.

O jogador A joga uma carta esperando, de alguma forma, facilitar uma jogada do jogador C em que ele force o jogador D a comprar --- ou seja, o jogador A tenta proporcionar ao jogador C a change de jogar uma carta que force o jogador D a comprar duas ou quatro cartas, ou ainda uma carta de uma cor que o jogador D seguramente não tem (aumentando as chances de que o jogador D tenha que comprar uma carta).

O abastecimento, portanto, depende das condições do adversário alvo: se ele está vulnerável ao não ter diversos tipos de cartas para jogar. Também depende de sua possibilidade de ganhar o jogo; gastar tempo abastecendo um jogador cheio de cartas é prejudicial, uma vez que isso pode eventualmente ajudá-lo a assumir uma posição de gerenciamento do jogo que acabará por prejudicar a dupla que o abasteceu.

\subsubsection{Isolamento}

O isolamento, ao contrário do abastecimento, não consiste em fazer com que um dos jogadores da dupla adversária compre mais cartas, mas, ao invés disso, tenta fazer com que eles não joguem. O isolamento, por tanto, é uma técnica que se utiliza de cartas Pular.

O Isolamento simples é o mais fácil de ocorrer: em uma dupla, um dos jogadores possui duas ou mais cartas Pular. Vamos supor que o jogador A possua três cartas Pular; ele pula o jogador B. O jogador C joga de modo a proporcionar ao jogador A a chance de jogar outra carta Pular, que de fato é jogada; o mesmo acontece na próxima rodada e o jogador B deixa de jogar por três rodadas seguidas. Um jogador foi ``isolado'' do jogo.

Isso é extremamente benéfico pois faz com que um jogador continue com muitas cartas na mão (não apenas para afastar dele a chance de vitória mas também para, em um caso de vitória da dupla de ataque, muitos pontos forem obtidos), faz com que o jogo esteja ainda mais sob controle por parte da dupla de ataque, e deixa sem reforços o outro jogador da dupla. Por outro lado, o que é um aspecto que pode ser tanto positivo quanto negativo, o jogador isolado pode se deixar levar pelas emoções, pretendendo, portanto, atacar de forma mais violenta (como visto acima na seção ``O Estímulo Passional'').

Se um jogador possui três cartas Pular e as joga em diferentes momentos do jogo, isso não é considerado um Isolamento. O isolamento é o uso repetitivo, seguido, de cartas Pular para, de forma prolongada, impedir um jogador específico de jogar. Ele, portanto, se diferencia do uso avulso destas cartas de ação por essa característica tática não existente no emprego aleatório (ou com base em outro princípio) das cartas em questão.

\subsubsection{Isolamento Sistemático}

O isolamento sistemático é uma variante mais rara do isolamento, em que ambos os jogadores das duplas possuem uma certa quantidade de cartas Pular, e as usam de forma simultânea para descartar várias cartas e, ao mesmo tempo, impedir que a outra dupla jogue.

Ao final de um isolamento sistemático, percebe-se que é como se o jogo tivesse sido paralisado e retornado ao normal, porém com uma das duplas com várias cartas a menos. É, portanto, um ataque muito vantajoso, porém difícil de ser empregado por causa das condições raras em que ele pode ser executado.

\subsubsection{Isolamento Fatal}

O isolamento fatal é uma espécie de isolamento sistemático, em que a mesma lógica do ``Curinga Surpresa'' é aplicada: a dupla deixa, como as últimas cartas em suas mãos, cartas Pular que são usadas para realizar um Isolamento Sistemático, de forma que assim, com três cartas na mão de um jogador e até mesmo quatro na mão de outro, uma dupla pode vencer de modo surpreendente com este ataque fulminante, ao qual é absolutamente impossível reagir uma vez que tenha sido iniciado. É ainda mais difícil de realizar do que um Isolamento Sistemático, uma vez que precise que ambos os jogadores tenham no mínimo duas cartas Pular, e tenham tido condições de deixá-las por último.

Tanto o abastecimento quanto o isolamento não foram discutidos no jogo de cunho individual por não fazerem sentido em um jogo deste tipo: em um jogo com 5 pessoas, uma pode até mesmo ``isolar'' um jogador (o próximo na sequência de jogadas, caso ela não se altere) ao impedi-lo repetidamente de jogar, mas isso não possui um efeito estratégico muito interessante, pois ele ainda precisa lidar com outros 3 jogadores, cujas estratégias de nada dependem desse único adversário afetado. No jogo em parcerias, contudo, atacar um dos jogadores constitui um ataque direto à estrutura de uma espécie de adversário único que é a outra dupla.

\section{Regras Alternativas: A Customização do UNO}

Depois de analisar as estratégias para o jogo oficial (na sua variante 11211), podemos agora finalmente pensar em flexibilizar, customizar o UNO; \textbf{modificar} suas regras ou \textbf{inserir} novas, e considerar o \textbf{impacto estratégico} dessas modificações em \textbf{jogos normais}, com \textbf{apenas dois jogadores} ou com \textbf{parceiros}.

Além disso, considera-se também \textbf{como proceder} em casos em que a regra nova parece não se adaptar bem a certas regras oficiais --- ou \textbf{como ela funciona} em conjunto com \textbf{outras regras alternativas}.

Mesmo essas regras (muitas delas) possuem muitas flexibilizações internas e escolhas a serem feitas; Considerei, dentro do possível, todas essas diferenças e expliquei por que escolhi um modo de implementar essas regras, e não outro. Se o motivo não for razoável o suficiente, as opções ficam em aberto (e deixo isso claro no final de cada regra).

Todas as regras contêm \textit{siglas}; isso ajuda depois a definir o seu tipo preferido de jogo e então criar um código para lembrar de todas as regras que ele encerra.

Você já deve ter tido contato com muitas regras, mas provavelmente não com todas. Divirta-se!

\subsection{Regras que Modificam o Objetivo do Jogo}

Essas regras modificam o Objetivo do Jogo. Lembre-se, do \textit{jogo}, não da \textit{partida}; a diversão ainda está em livrar-se das cartas na mão, mas o que você faz \emph{assim} que isso acontece pode ser modificado.

\subsubsection{Número de Pontos}

Uma observação inicial a ser feita é que uma das coisas mais usualmente flexibilizadas é o número de pontos no UNO a serem alcançados para que um jogador possa ser considerado vencedor: os 500 geralmente se tornam menos caso os jogadores não queiram que uma partida torne-se muito longa.

Mesmo sem adotar uma das regras que modificam o objetivo do jogo, tenha em mente que esse número pode sempre ser flexibilizado para as suas vontades (princípio que vale também para, por exemplo, a Regra dos Pontos Reversos --- ver abaixo). 

\subsubsection{Regra do Torneio de Vitórias [TV]}

Ao invés de acumular 500 pontos, você pode abandonar de vez a ideia de pontos e jogar apenas por número de vitórias; em um jogo com várias pessoas, o primeiro a vencer \textit{x} vezes é o campeão; em um jogo de duas pessoas, uma melhor de 3, 5 ou 7\footnote{Uma ``melhor de três'' é uma expressão utilizada para definir uma sequência de três jogos em que aquele entre dois adversários que vence mais partidas é o vencedor. É particularmente interessante pois se há apenas dois jogadores e três partidas, aquele que vence duas é automaticamente o campeão, não precisando que uma terceira partida seja jogada. Em uma ``melhor de cinco'', por exemplo, aquele que vence três partidas é o campeão (a mesma lógica de um jogo de vôlei, por exemplo). Essa ideia, contudo, não serve para jogos com mais de dois jogadores.} pode ser o esquema usado. De uma forma ou de outra, a contagem de pontos sequer é feita.

\paragraph{Consequências}

Esta regra, embora possa parecer inócua, possui profundas consequências na dinâmica de jogo e em aspectos externos do jogo. Vamos começar por este último caso: no Torneio de Vitórias, os pontos não são contados ao final da partida. Isso pode levar a partidas mais rápidas. Por outro lado, as diferenças estratégicas provocadas por essa regra, como veremos adiante, pode acabar compensando esse tempo ao tornar mais longo o jogo.

No quesito estratégia, temos que tudo que importa é \textit{vencer}; nesse caso, muitas das estratégias estudadas deixam de trazer consigo qualquer espécie de \textbf{risco}. É importante que se entenda isso: se um jogador perde, não importa mais com que cartas foi pego na mão; isso é irrelevante agora. A carta zero, por exemplo, deixa de ser uma carta valorizada como tática defensiva em momentos de tensão, e o Curinga Surpresa é uma estratégia que passa a ser empregada com uma frequência absurdamente maior. Esse tipo de modificação na dinâmica do jogo ocorre em qualquer tipo de jogo UNO: individual, com dois jogadores ou em duplas.

Como dito anteriormente, isso também possui impacto no tempo de jogo: sem a pressa para livrar-se das cartas e com a paciência para explorar todos os tipos de estratégia que envolvem \textit{manter} as cartas na mão, o jogo pode vir a demorar um pouco mais --- fazendo com que, no final das contas, um jogo pautado por esta regra não seja, afinal de contas, mais rápido.

\paragraph{Variações}

Não só as vitórias podem ser contadas, mas a exemplo da Regra dos Pontos Reversos (a ser vista logo abaixo), as derrotas podem ser a base da contagem de pontos: todos os jogadores podem começar com um número \emph{x} de pontos e ir perdendo um ponto a cada derrota. 

\paragraph{Indicado}

para pessoas que queiram experimentar mais livremente com possibilidades estratégicas, sem sentir a pressão da pontuação, ou para pessoas que queiram simplesmente jogar sem ter que contar pontos.

\subsubsection{Regra dos Pontos Reversos [PR]}

Ao invés de acumular 500 pontos, todos os jogadores \textit{começam} com 500 pontos e, ao longo do jogo, vão perdendo-os com base nas cartas que permanecem em suas mãos após a vitória de algum dos jogadores.

\paragraph{Consequências}

Esta regra pode, novamente, parecer não fazer muita diferença, mas na verdade faz: em um jogo onde o objetivo é acumular pontos a partir dos pontos remanescentes nas mãos dos perdedores, um objetivo secundário do jogo torna-se, muitas vezes, fazer com que os jogadores \textit{comprem} cartas e \textit{mantenham} essas cartas na mão.

Ao mesmo tempo em que o fator da continuidade tem sido apontado ao longo do livro como a estratégia mais relevante para a vitória de uma partida e, consequentemente, do jogo, o UNO é um jogo de flexibilidade em que várias táticas podem ser seguidas (por vezes simultaneamente), com variados graus de sucesso. Para ganhar mais rapidamente, fazer com que os jogadores comprem cartas não deixa de ser importante.

Em um jogo onde esses pontos remanescentes não se combinam para formar a pontuação de alguém, mas apenas são subtraídos do total individual, o importante é tentar ganhar com cautela, para não perder muitos pontos caso alguém seja mais rápido. Enquanto o jogo normal favorece o ataque e artimanhas que vão na contramão da continuidade apenas em nome da ofensividade, este tipo de jogo acaba incentivando uma atitude mais defensiva, em que tudo que é feito o é exacerbando a importância de se livrar de cartas de ação para evitar ter muito prejuízo com a vitória alheia.

Além disso, essa dinâmica de jogo faz com que o jogo fique mais longo: quando alguém perde seus 500 pontos, sai do jogo, mais os outros continuam a jogar até que reste apenas um jogador que ainda possui pontos. Dessa forma, o jogo demora mais para chegar ao fim (embora não o mesmo tempo para todos os envolvidos).

\paragraph{Indicado}

para pessoas que prefiram um jogo em que apenas os mais prudentes (não necessariamente os \emph{melhores} jogadores) vão resistindo ao tempo, ajudando a moldar uma espécie de ``partida final''.

\subsubsection{Regra da Morte Súbita [MS]}

Nesse jogo, não só ficar com nenhuma carta na mão é importante, mas também continuar a jogar. Esta é, \emph{par excellence}, a variação da continuidade: o jogador que for obrigado a comprar uma carta (exceto por meio de cartas Comprar Duas / Comprar Quatro) é eliminado do jogo, e os outros continuam até sobrar apenas um ou alguém descartar todas as cartas da mão. Note-se que por ``ser obrigado a comprar uma carta'' deve-se entender que o próprio ato de compra leva à eliminação --- o jogador que não puder jogar com as cartas que tem à mão, quando for sua vez, é eliminado.

\paragraph{Consequências}

Em primeiro lugar, as partidas se tornam mais rápidas, pois quem é obrigado a comprar não permanece jogando. Em segundo lugar, nenhum outro aspecto e impulso do jogo se torna mais primordial do que a continuidade: fazer o outro comprar não é mais relevante, o que é importante é planejar o jogo de modo a aumentar as suas chances de \emph{continuar jogando}. É um jogo de dinâmica rápida e, por isso mesmo, mais tenso que os outros.

\paragraph{Variações}

\begin{description}
\item[Morte Súbita Sem Proteção (MS1)]{Qualquer jogador, desde o princípio do jogo, está sob o risco de ser eliminado.}
\item[Morte Súbita Com Proteção (MS2)]{Apenas depois de um número \textit{x} de rodadas a Morte Súbita começa a valer. Número sugerido: duas rodadas.}
\item[Morte Súbita Com Trava (MS3)]{Apenas após algum evento do jogo a Morte Súbita começa a valer; o primeiro Curinga a ser jogado, a primeira carta de um determinado número a ser jogada, a quinta carta de ação a ser jogada, etc.}
\end{description}

\paragraph{Possíveis Contradições e Conexões}

\begin{description}
\item[Pontuação]{Para fazer a pontuação, vários métodos anteriormente vistos podem ser usados: quem for eliminado pode ter suas cartas revertidas como pontos para o vencedor (que ainda está por ser decidido), ou pode perder pontos (se o esquema de pontos reversos for utilizado), etc.}
\item[Regra do Passe ou do Inferno]{Se esta regra for aplicada em conjunto com qualquer variação da Regra do Passe ou com a Regra do Inferno, não ter uma carta Comprar Duas, Comprar Quatro ou carta de ataque que possa ser jogada garante uma eliminação.}
\item[Regra do Final Limpo]{Caso a regra do Final Limpo seja aplicada, caso a pessoa possua como última carta uma carta que ela não pode jogar, este jogador deve ser eliminado.}
\item[Regra do Corte]{Se aplicada conjuntamente com a regra do corte, pode gerar algumas dúvidas: supondo que o jogador A vá jogar, mas não possa (\emph{i.e} será eliminado), mas antes que \emph{declare} que será eliminado, outro jogador corta sua vez --- neste caso, o jogador A não é eliminado. Se ele já tiver declarado que será eliminado, o jogador B é quem será cortado, e o jogador A terá sido eliminado.}
\item[Regras da Paz]{As regras de paz impedem que cartas de ação sejam jogadas durante períodos específicos de tempo. Assim sendo, se não for possível jogar uma carta de ação num dado momento e o jogador, posto isso, \emph{não conseguir} jogar \emph{carta alguma}, a lógica permanece: o jogador é eliminado.}
\item[Regra do Descarte Duplo]{Se jogada em conjunto com a regra do Descarte Duplo: se o jogador A não tiver o que jogar e, para não ser eliminado, mentir e for pego mentindo, é eliminado (e os outros jogadores agem de acordo com os mecanismos da Regra do Descarte Duplo. Se jogar de forma oculta, o acusarem de mentira, mas sua jogada for legítima, novamente os mecanismos do Descarte Duplo se aplicam.}
\item[Outras Punições]{Sempre que um jogador for obrigado a comprar uma carta devido a uma \emph{punição} (como, por exemplo, a Regra do Silêncio ou a do Monte) isto não o elimina; um jogador é apenas eliminado se não puder jogar nenhuma carta que estiver em sua mão no momento --- por exemplo, se com a Regra de Chukcki aplicada um jogador não puder jogar uma carta (ou uma combinação de cartas) que altere a cor da pilha de descarte, ele será eliminado.}
\end{description}

\paragraph{Indicado}

para quem tem pouco tempo disponível para jogar ou para quem prefira um jogo mais imprevisível.

\subsection{Regras que Modificam a Dinâmica de Jogo}

Essas regras modificam o modo como o jogo é jogado, partindo de características como o número de cartas com que uma pessoa começa o jogo, passando por regras que determinam aquilo que pode ser descartado, e indo até regras que restringem o comportamento de um jogador.

\subsubsection{Regra do Final Limpo [FL]}

Os jogadores podem apenas vencer o jogo com cartas numéricas, cartas Pular e Inverter. Curingas e cartas que obriguem os jogadores a comprar não podem, portanto, ser usadas para vencer.

\paragraph{Consequências}

A consequência mais imediata é o fim da estratégia do Curinga Surpresa (já discutido anteriormente), uma vez que você não pode guardar mais o Curinga como última carta.

Esta regra costuma ser mais impactante, contudo, quando jogada em conjunção com variações da Regra do Passe (ver mais adiante). É uma estratégia comum (pois muito eficaz) guardar um Comprar Duas ou Comprar Quatro como última carta (neste último caso, é mesmo um Curinga Surpresa, mas com outras vantagens) pois caso alguém tente fazê-lo comprar cartas na última rodada, você pode se defender. Com a Regra do Final Limpo, ganhar o jogo desse jeito (que alguns podem considerar chato por ser até muito fácil) torna-se impossível.

\paragraph{Variações}

Pode-se estipular que só se pode ganhar o jogo com cartas numéricas; ou que não se pode ganhar o jogo com Curingas (mas cartas de Compra são permitidas, uma regra que faz sentido se o grupo não joga com Regras do Passe).

\paragraph{Possíveis Contradições e Conexões}

\begin{description}
\item[Impossibilidade]{À exemplo da Primeira Regra da Paz (ver adiante), essa regra não estipula punições, e sim determina a \emph{impossibilidade} de jogar certas cartas para ganhar o jogo.}
\item[Regra do Passe e do Inferno]{Se jogada em conjunto com a Regra do Passe ou a Regra do Inferno, significa que se um jogador A joga uma carta Comprar Duas e o jogador B estiver UNO (sua carta sendo um outro Comprar Duas ou uma carta que poderia ser jogada, mas não pode por causa do Final Limpo), ele não poderá usar essa carta para vencer o jogo; deverá comprar duas cartas.}
\item[Desafio Polivalente]{Se jogada em conjunto com a Regra do Desafio Polivalente, a lógica é a seguinte: Se um jogador A não pode descartar o Curinga (por exemplo) pois é a sua última carta, e precisa comprar, se alguém o desafia, deve comprar quatro cartas, pois ele de fato precisava comprar.}
\item[Regra do Descarte Duplo]{Se jogada em conjunto com a regra do Descarte Duplo: se o jogador A tiver como última carta uma carta que não possa ser jogada, ele pode tentar ganhar o jogo mentindo. No caso de ser descoberto, os mecanismos da Regra do Descarte Duplo se aplicam. Se não tiver mentido, os mesmos mecanismos se aplicam.}
\item[Regra do Corte]{Se um jogador A possui como última carta uma carta que não pode ser jogada para vencer, o jogador também não pode vencer com esta carta ao fazer um corte. Ele não pode jogá-la, seja numa jogada normal, seja num corte.}
\item[Regra do Descarte Rápido e Duplicatas]{Como os jogadores não podem \emph{vencer} com as cartas que forem combinadas no começo do jogo, eles não podem fazê-lo, mesmo se não estiverem UNO. Por exemplo: o jogador A não pode vencer com um Curinga, e também não pode vencer jogando dois curingas de uma vez só (sem ficar UNO).}
\end{description}

\paragraph{Indicado}

para quem acredita que assim o jogo se torna mais difícil e taticamente complexo, uma vez que a estratégia do Curinga Surpresa pode ser (dependendo de algumas combinações de regras específicas) muito fácil de ser realizada.

\subsubsection{Regra do Flash UNO [FU]}

Ao invés das oficiais 7 cartas, todos os jogadores começam o jogo com 4 cartas.

\paragraph{Consequências}

Com menos cartas iniciais, a duração de cada partida tende a ser menor. Em compensação, as estratégias tendem a ser mal formuladas ou não-formuladas \emph{at all}, incluindo certa confusão, desde o início, entre que papel se deve exercer no jogo (o de gerenciador ou de fugitivo do gerenciamento alheio). Dessa forma, é um jogo em que a estratégia é exercida de forma limitada, e pode apenas vir a exercer um papel mais proeminente na decisão da partida se todos os jogadores acabarem por comprar mais cartas.	

\paragraph{Variações}

Na verdade o número de 4 cartas é uma arbitrariedade; poderia ser qualquer número de 6 para baixo (com o 6 não fazendo muita diferença, e o 2 fazendo uma diferença verdadeiramente \emph{excruciante} em termos de importância tática).

\paragraph{Indicado}

para quem tem pouco tempo disponível para jogar; para aqueles que desejam um jogo mais rápido.

\subsubsection{Regra do Battlefield UNO [BU]}

O Battlefield UNO é o inverso do Flash UNO: Ao invés das oficiais 7 cartas, todos os jogadores iniciam o jogo com 9 cartas.

\paragraph{Consequências}

As possibilidades de gerenciamento são maiores e o jogo geralmente torna-se mais longo e intrincado.

\paragraph{Variaçãoes}

Assim como no Flash UNO, o número 9 é arbitrário: pode ser qualquer número de 8 para cima.

\paragraph{Indicado}

para quem deseja um jogo em que muitas reviravoltas e tramas estratégicas possam ocorrer. Indicado também especialmente para jogos com apenas dois jogadores, em que um número de apenas 7 cartas pode tornar o jogo rápido demais.

\subsubsection{Regra do Corte [C]}

Uma das regras alternativas mais populares, a Regra do Corte estipula que cartas \emph{idênticas} (iguais em número/símbolo e cor) possam ser jogadas consecutivamente por dois jogadores que não estão necessariamente na ordem normal do jogo. Na prática, se pensarmos em um jogo-exemplo entre os jogadores A, B, C, D, E e F (que jogam nesta ordem) poderemos ter uma ideia mais precisa do sentido da regra:

O jogador A joga um 5 Verde na pilha de descarte. O jogador B prepara-se para jogar, mas uma vez que tenha demorado muito, o jogador D (que não é o próximo a jogar) joga uma carta idêntica (outro 5 Verde, portanto) na pilha de descarte. Essa jogada é chamada de ``corte'', pois corta a ordem natural do jogo. Depois do corte, o jogo prossegue a partir do ponto de corte; o jogador B, cortado, não joga, e nem o jogador C: agora o jogador E continua, e depois dele o jogador F, e então o jogador A, e apenas aí o jogador B poderá jogar novamente.

Outros exemplos, utilizando-se dos mesmos jogadores: o jogador C joga um 7 Vermelho. O jogador A poderia cortar o jogo com seu 7 Vermelho à mão, mas o jogador D foi mais rápido e continuou o jogo; agora já não há mais possibilidade de corte. O corte só pode acontecer enquanto o próximo jogador não joga seu turno; lembre-se que comprar uma carta conta como jogar um turno. Além disso, se o jogador E joga um 2 Azul e o jogador F também possui um 2 Azul, ele pode jogá-la normalmente; é como um corte, porém já que o jogador que corta é o mesmo que normalmente jogaria, não há diferença alguma.

\paragraph{Consequências}

O jogo torna-se muito mais tenso e animado, uma vez que \emph{todos os jogadores devem prestar atenção ao jogo o tempo todo}. Isso ocorre porque agora qualquer jogada é oportuna e, ao mesmo tempo, na própria vez, você deve jogar rapidamente para evitar que outro corte sua vez.

O pensamento estratégico de longo prazo, contudo, é um tanto prejudicado nesse jogo, uma vez que existe uma pressão para que as jogadas sejam feitas rapidamente, sob o risco de perder a vez.

\paragraph{Possíveis Contradições e Conexões}

\begin{description}
\item[Momentos de Corte]{O corte pode ser feito por qualquer jogador, a qualquer momento do jogo em que se pode normalmente jogar. Um jogador pode ``ficar UNO'' através de um corte, e pode também ganhar o jogo a partir de um corte. Um jogador não pode cortar a si mesmo, contudo, quando uma carta de compra tiver sido jogada (para se livrar da punição de comprar o número estipulado de cartas).}
\item[Compra]{Se o jogador A jogar um 5 Vermelho e o jogador B precisar comprar uma carta, o jogador B pode ser cortado antes de comprar e depois de comprar. Neste último caso, quem está sendo cortado, na verdade, é o jogador C; se o jogador B comprar e puder \emph{jogar aquilo que comprou}, ninguém pode cortá-lo.}
\item[Jogadas Equivocadas]{Da mesma forma que jogadas equivocadas devem ser punidas com a compra de duas cartas, cortes equivocados devem ser penalizados da mesma forma.}
\item[Punições]{Um jogador não pode cortar a si mesmo quando é sua vez de receber uma punição, como dito acima; dessa forma, se o jogador A joga uma carta Comprar Duas Verde, o jogador B não pode cortar a si mesmo para que o jogador C compre quatro cartas. Entende-se, contudo, que algum \emph{outro jogador}, que não o B, possa cortar o jogo, porque quem ele cortará, na verdade, será o jogador C (de modo que o jogador B ainda precisa comprar as duas cartas).\footnote{A maioria das pessoas, contudo, joga com alguma variação da regra do passe, e vai compreender melhor como o corte se aplica ao ler o ítem logo abaixo deste, o ``Punições Acumuladas''}.}
\item[Punições Acumuladas]{Quem preferir jogar em conjunto com qualquer variação da Regra do Passe ou com a Regra do Inferno deverá jogar de acordo com a seguinte dinâmica: se o jogador A joga uma carta Comprar Dois Amarela, qualquer outro jogador pode cortar \emph{o jogador B}, então se o jogador C joga a outra carta Comprar Dois Amarela, o jogador \textbf{D} deve comprar quatro cartas. Contudo, há ainda outra opção: deixar que o jogador B compre as duas cartas e então cortar o jogador C, antes que ele jogue ou compre uma carta. A coisa fica ainda mais curiosa e interessante com as cartas Curinga Comprar Quatro, uma vez que se uma for jogada, três cortes podem ser feitos, e alguém pode ser obrigado a comprar 16 cartas.}
\item[Cartas de Ação com Cortes]{Quando cortes forem feitos envolvendo cartas de ação, algumas regras específicas se aplicam:
\begin{description}
\item[Carta Pular]{Se o jogador A joga uma carta Pular (fazendo com que o jogador B não jogue), qualquer corte feito sobre essa carta anula o efeito da carta cortada, e agora apenas o efeito novo sobra. Dessa forma, num jogo entre A, B, C e D, se o jogador A joga uma carta pular e for cortado pelo jogador C, quem na verdade é pulado é o jogador D, e o jogador A deve, portanto, jogar novamente.}
\item[Carta Inverter]{Dessa vez, de certa forma o corte considera o efeito anterior: se os jogadores A, B, C, D e E jogam nesta ordem, o jogador A joga uma carta Inverter e faz com que o jogador E seja o próximo. Antes que o jogador E jogue, contudo, o jogador B corta: se, então, a ordem do jogo acabara de se tornar A--E--D--C--B, a ordem volta a ser A--B--C--D--E, e como começa agora do jogador B, o jogador C é o próximo a jogar.}
\item[Curinga]{Se um Curinga for cortado, a escolha de cor anterior é ignorada e o jogador que fez o corte escolhe a cor. Como de praxe, o jogo prossegue do jogador que cortou, com a cor escolhida agora por ele.\footnote{Para entender melhor sobre a dinâmica do Curinga Comprar Mais Quatro, contudo, leia o ítem ``Punições''.}}
\end{description}
}
\item[Duas Pessoas]{Com UNO jogado entre duas pessoas, o efeito desta regra é diminuto, mas um pouco significativo; em um jogo oficial 11211, funciona como uma espécie de Regra das Duplicatas (o jogador que jogou uma carta pode jogar outra idêntica ao mesmo tempo). Além disso, como em um jogo entre duas pessoas uma carta Pular significa que o mesmo jogador joga novamente, cortar esta carta Pular significa impedir essa jogada.}
\item[Parcerias]{Em um jogo de parcerias, a regra do Corte costuma ser pouco benéfica. Imaginando um jogo exemplo entre os jogadores A, B, C e D, duplas A e C, B e D, e a ordem do jogo é ABCD, vamos acompanhar o ponto de vista do jogador A: cortar durante a vez do jogador B significa apenas jogar duas vezes. Bom, mas é, como foi discutido no ítem acima, como a Regra das Duplicatas. Cortar o próprio parceiro é péssimo: o jogador A jogaria novamente, mas o jogador B também. Cortar o jogador D faria com que o jogador A e o C jogassem novamente, mas também o C. O que se demonstra é que quanto maior o número de pessoas, mais efetivo o Corte pode ser. Em pequenos números e, em especial, em Parcerias, a regra não surte tanto efeito positivo para quem a aplica.}
\item[Chukcki]{Ao jogar em conjunto com a Regra do Chukcki, ocorre o seguinte: considere os jogadores A, B e C (que estão jogando entre mais jogadores). O jogador A joga uma carta numérica, e o jogador B joga uma carta de mesmo valor numérico. O jogador C deve, através de quaisquer mecanismos apropriados, mudar a cor da pilha de descarte. As jogadas do jogador A e B podem ser originadas de um Corte. Contudo, \emph{não considerando outras regras}\footnote{A Regra do Descarte Rápido tornaria um Corte possível para o jogador C, por exemplo.}, uma vez que o corte trata de cartas \emph{idênticas}, ninguém vai conseguir mudar a cor da pilha de descarte com um corte.}
\item[Duplo Descarte]{Ao jogar com a regra do Duplo Descarte, apenas as cartas \emph{declaradas} valem: não se pode cortar uma carta que tenha sido omitida (mesmo que depois se descubra qual carta é) e não se pode cortar de forma ``omitida''.}
\end{description}

\paragraph{Indicado} 

para jogos com \emph{muitos jogadores}, em que quando um joga, as pessoas do lado mais distante da mesa, que vão demorar para jogar novamente, não costumam prestar atenção ao jogo. Bom também para quem deseja um jogo mais imprevisível e complexo. Por outro lado, às vezes a tensão e a pressa ao jogar podem comprometer um estilo de jogo mais bem pensado, estratégico.

\subsubsection{Regra do Descarte Rápido [DR]}

Outra regra relativamente comum, a Regra do Descarte Rápido determina que cartas de mesmo número/símbolo possam ser jogadas na mesma rodada por um mesmo jogador. Dessa forma, o jogador A pode jogar três cartas 5 (uma verde, uma vermelha e uma azul), o jogador B pode jogar duas cartas (um 6 Azul e um 6 Verde), o jogador C pode jogar 4 cartas (Um 7 Verde, outro 7 Verde, um 7 Azul e um 7 Amarelo), e aí por diante. A única exigência para que uma carta possa ser jogada em cima de outra é que o número e/ou símbolo seja correspondente.

\paragraph{Consequências}

Adiciona-se ao jogo uma dinâmica nova em que os jogadores não apenas ficam para trás por não terem cartas para jogar, mas também avançam, tornam-se mais rápidos que outros, por conseguirem se livrar de duas ou mais cartas em uma mesma rodada. Isso altera estratégias e insere um elemento de surpresa a mais no jogo, o que, por sua vez, torna o pensamento tático mais conveniente. 

\paragraph{Possíveis Contradições e Conexões}

\begin{description}
\item[Diferenças]{Cartas Comprar Quatro não podem ser jogadas juntamente com cartas 4; cartas Comprar Duas não podem ser jogadas com cartas 2.}
\item[Acúmulo]{Cartas Comprar Quatro podem ser jogadas juntamente com outras Comprar Quatro (fazendo com que o próximo jogador obrigatoriamente compre o número acumulado de cartas. O mesmo vale para cartas Comprar Duas.}
\item[Ordem das Cartas]{No descarte rápido, é como se você jogasse a carta apropriada de acordo com a carta na pilha de descarte, e então adicionasse cartas de números ou símbolos iguais \textbf{acima} da carta apropriada. Dessa forma, se há um 3 Amarelo na pilha de descarte, e você pretende jogar, ao mesmo tempo, um 9 Azul, um 9 Amarelo e um 9 Vermelho, o 9 Amarelo deve ser jogado primeiro e, depois dele, você pode escolher a ordem em que quer deixar os outros 9s da jogada (primeiro o azul e depois o vermelho, ou o contrário).}
\item[Momento de Acúmulo e UNO]{As cartas podem ser jogadas de forma agrupada em qualquer momento em que uma única pode ser jogada normalmente. Dessa forma, uma pessoa pode ``ficar UNO'' ao jogar duas cartas ao mesmo tempo (tendo três cartas na mão e, após o descarte, tendo apenas uma), e uma pessoa pode vencer o jogo sem precisar se declarar UNO ao jogar suas duas últimas cartas de forma conjunta. Isso cunha um novo tipo de estratégia, similar ao Curinga Surpresa, conhecida como o \textbf{Grupo Surpresa} (ver abaixo).}
\item[Vez]{Se um jogador for jogar várias cartas ao mesmo tempo, deve fazê-lo de uma vez só, ao mesmo tempo. Se não o fizer, deve declarar em tom suficientemente alto que sua jogada ainda não acabou e que ele ainda tem mais cartas a descartar\footnote{Especialmente em jogos em que a Regra do Passe é aplicada, alguns jogadores podem ter dificuldade para organizar na mão um grande número de cartas, o que leva a problemas pra jogar várias cartas de uma vez só rapidamente.}; se o próximo jogador na ordem do jogo já tiver comprado uma carta ou jogado uma carta (\emph{i.e} iniciado sua vez), o jogador anterior não poderá jogar mais cartas.}
\item[Cartas de Ação em Conjunto]{Quando cartas de ação forem jogadas em conjunto, algumas regras específicas se aplicam:
\begin{description}
\item[Carta Pular]{Se duas cartas Pular forem jogadas, dois jogadores são pulados; se três cartas forem jogadas, três são pulados, e assim por diante. Em um jogo com $n$ jogadores, se $n - 1$ cartas Pular forem jogadas, o próprio jogador que as jogou joga novamente. Se qualquer número de cartas de $n$ para cima for jogado, o próprio jogador é ignorado e apenas os adversários contam na contagem de ``pulos'' até a verificação de quem é o próximo a jogar (exemplo: num jogo entre A, B e C o jogador A jogou quatro cartas Pular; pulou o B, o C, depois o B novamente e enfim o C novamente, de forma que o B é o próximo a jogar.)\footnote{Cuidado como o jogo em duplas: estratégias derivadas do Isolamento exigem que uma carta Pular seja jogada de cada vez}.}
\item[Carta Inverter]{Se um número par de cartas Inverter for jogado em conjunto, então a ordem do jogo não se altera. Se um número ímpar for jogado, a ordem do jogo se altera.}
\item[Curinga]{Se vários curingas forem jogados, contam como um: o jogador que os jogou escolhe uma cor e o jogo prossegue normalmente.}
\end{description}
}
\item[Regra do Corte]{Ao jogar conjuntamente com a Regra do Corte (ver acima), você pode cortar jogando duas ou três cartas ao mesmo tempo (uma delas, contudo, deve ser a carta responsável por possibilitar o corte, e ela deve ser a primeira na ordem das cartas, como visto acima).}
\end{description}

\paragraph{Indicado} 

para uso geral, por melhorar consideravelmente as complexidades estratégicas do jogo, sem necessariamente prejudicar outras características do jogo (nem deixá-lo muito mais curto).

\subsubsection{Regra das Duplicatas [DD]}

A Regra das Duplicatas possui o mesmo mecanismo da Regra do Descarte Rápido, porém apenas cartas \emph{idênticas} podem ser jogadas em conjunto.

\paragraph{Consequências}

As mesmas consequências da Regra do Descarte Rápido se aplicam (inclusive a possibilidade de usar a estratégia do Grupo Surpresa). Contudo, esta regra restringe a ocorrência do descarte em grupos, além de impossibilitar que mais que duas cartas sejam jogadas de cada vez, a não ser no caso das cartas Comprar Quatro.

\paragraph{Indicado} 

para aqueles que veem com desconfiança a estratégia do Grupo Surpresa, que oferece grande chance de vitória com um risco minimizado. Por esta razão, a Regra das Duplicatas é um meio de ainda possibilitá-la, mas dificultar sua ocorrência, e, ao mesmo tempo, também inserir o elemento tático da Regra do Descarte Rápido (porém, de modo semelhante, com mais restrições).

\subsubsection{Primeira Regra da Aritmética [AA]}

Com uma carta numérica na pilha de descarte, você pode jogar duas ou mais cartas numéricas de mesma cor que, ao serem somadas ou subtraídas, resultem no número daquela primeira carta. No caso da soma, mais do que duas cartas podem ser jogadas; no caso da subtração, apenas duas.

Dessa forma, se há um 8 Azul na pilha de descarte, um 1 Azul e um 7 Azul podem ser jogados (pois sua soma é 8). Um 10 Azul e um 2 Azul podem ser jogados também (pois sua subtração é um 8), assim como um 4, um 3 e um 2 Azuis (pois a soma deles resulta em 8).

\paragraph{Consequências}

Ela insere o mesmo tipo de complexidade estratégica que as regras do Descarte Rápido e das Duplicatas: é possível agora avançar no jogo por suas próprias jogadas, e não quando os outros ficam para trás.

Essa regra, contudo, torna essa possibilidade um pouco mais difícil de se concretizar, além de mais difícil ainda se jogada em conjunto com outras regras como a da Velocidade, em que não se pode demorar muito para pensar.

Outra consequência é a de que as jogadas tornam-se mais imprevisíveis, uma vez que mudar a cor da pilha de descarte se torna mais fácil.

\paragraph{Possíveis Contradições e Conexões}

\begin{description}
\item[Momento]{É possível ficar UNO ao usar do mecanismo da regra e é possível também ganhar o jogo usando-se do mecanismo.}
\item[Cor]{As cartas a serem jogadas para formar o resultado da soma ou da subtração devem ser da mesma cor, mas não da mesma cor da pilha de descarte do momento.}
\item[Vez]{Assim como na Regra do Descarte Rápido, todas as cartas devem ser jogadas ao mesmo tempo, de uma vez só. Se isto não ocorrer, o jogador deve declarar em tom suficientemente alto que sua jogada ainda não acabou e que ele ainda tem mais cartas a descartar; se o próximo jogador na ordem do jogo já tiver comprado uma carta ou jogado uma carta (\emph{i.e} iniciado sua vez), o jogador anterior não poderá jogar mais cartas. Na regra da Aritmética, esquecer de incluir o zero (que pode ser incluído em qualquer soma ou subtração, desde que seja da mesma cor que as outras cartas) pode gerar um problema do tipo.}
\end{description}

\paragraph{Indicado} 

para uso geral, em especial grupos que se identifiquem com uma versão um pouco mais complexa do Descarte Rápido --- esta é uma regra em que quase a mesma dinâmica é alcançada, mas é preciso pensar com cuidado para formar uma boa estratégia.

\subsubsection{Segunda Regra da Aritmética [AB]}

Você pode jogar uma carta numérica que seja a soma das duas últimas cartas da pilha de descarte, se elas (as cartas na pilha de descarte) forem numéricas e de mesma cor.

\paragraph{Consequências}

Ao contrário da Primeira Regra da Aritmética, discutida anteriormente, esta carta não possibilita que várias cartas sejam jogadas ao mesmo tempo, não alterando significativamente a dinâmica de jogo. Por outro lado, aumenta as possibilidades de jogo: em uma pilha de descarte cujas últimas cartas tenham sido um 2 e um 3 vermelhos (sendo a do topo o 3), você pode jogar um 3 de qualquer cor, qualquer carta vermelha, e agora também um 5 de qualquer cor.

\paragraph{Possíveis Contradições e Conexões}

\begin{description}
\item[Momento]{É possível ficar UNO com esta regra, e também vencer o jogo utilizando os mecanismos desta regra.}
\item[Duplo Descarte]{As duas últimas cartas da pilhas de descarte referem-se, evidentemente, apenas à pilha de descarte das cartas visíveis. Contudo, se for possível jogar uma carta devido a esta regra, esta carta pode ser posta na pilha de descarte das cartas omitidas (ela é considerada uma carta legítima, afinal.)}
\end{description}

\paragraph{Indicado} 

Para uso geral; a regra, contudo, não consegue alterar dramaticamente a dinâmica da partida por si mesma.

\subsubsection{Regra da Satisfação [SAT]}

Esta regra determina que um jogador que não tenha nenhuma carta que possa ser jogada na rodada compre não apenas uma carta, mas \emph{prossiga} comprando até que compre uma carta que possa ser usada (o que pode significar uma, duas, cinco ou vinte compras numa mesma rodada).

\paragraph{Consequências}

Do ponto de vista estratégico, essa regra faz com que a continuidade seja um princípio ainda mais valorizado, uma vez que é imprevisível o preço a ser pago por não jogar em uma rodada. Além disso, é uma estratégia que desestimula a tática do consumismo, enfraquecendo o jogo taticamente.

Além disso, é importante notar que a regra da satisfação torna impossível a execução de qualquer modalidade de Trébuchet, uma vez que nenhuma pessoa pode parar de comprar até que possa jogar novamente. Por outro lado, quando o Trébuchet não é aplicado e todos os jogadores devem passar por entediantes duas ou três rodadas sem jogar, a regra acaba por evitar isso, acelerando o ritmo do jogo (com o custo de um jogador comprando muitas cartas).

Como consequência tangencial (e prejudicial) para a questão da continuidade no jogo, é muito mais complicado esperar que a cor na pilha de descarte permaneça inalterada até que você jogue; por exemplo, em um jogo de quatro pessoas em que você \emph{sabe} que ninguém mais tem cartas azuis além de você, se você jogar uma carta azul pode ter muito mais ``esperança tática'' no fato de que na próxima rodada a cor do descarte ainda será azul. Contudo, com esta regra, os outros jogadores \emph{certamente} jogarão, e a cor da pilha de descarte tornar-se-á muito mais imprevisível .

\paragraph{Possíveis Contradições e Conexões}

\begin{description}
\item[Insatisfação]{Esta regra não pode ser aplicada junto à regra da Insatisfação; elas são diametralmente opostas.}
\end{description}

\paragraph{Indicado} 

para jogadores que queiram uma dinâmica de jogo mais acelerada xe simplificada, embora isto prejudique muito o pensamento tático do jogo, e portanto não seja recomendada para uso geral.

É recomendado que seja jogada com alguma regra em que várias cartas possam ser jogadas na mesma rodada (como Descarte Rápido, Duplicatas, Primeira Regra da Aritmética), pois às vezes uma pessoa pode acabar acumulando cartas demais. 

\subsubsection{Regra do Consolo [CON]}

Esta regra torna possível jogar com as cartas compradas de uma compra obrigatória (ou seja, provocada por cartas como o Curinga Comprar Quatro ou um Comprar Duas).

\paragraph{Consequências}

É uma regra que, em geral, acaba neutralizando uma importante faceta da estratégia do UNO, especialmente no jogo de parcerias: ao fazer alguém comprar, a jogada tem o potencial de possuir um duplo impacto --- o próximo jogador compra uma carta, e o outro, que jogaria depois do que foi ``punido'', tem que lidar com a cor que da carta Comprar Duas (ou a cor escolhida pelo jogador que jogou um Comprar Quatro). Com essa regra, porém, o jogador que primeiro sofreu o impacto da jogada tem a chance de jogar, subvertendo as intenções originais da carta quanto ao seu \emph{segundo} impacto.

É, portanto, em geral, uma regra que possui pouco sentido estratégico, e não costuma adicionar diversão ou outros atrativos ao jogo.

\paragraph{Possíveis Contradições e Conexões}

\begin{description}
\item[Regra do Corte]{Se jogada com a Regra do Corte, a seguinte dinâmica é seguida: o jogador A joga, por exemplo, um Comprar Duas Verde. O jogador B compra as duas cartas e, como agora tem direito a jogar, um outro jogador pode cortá-lo com um Comprar Duas Verde --- o jogador B foi cortado. Contudo, se estas duas regras forem aplicadas com a Regra do Passe, o jogador B pode ser cortado antes e depois de comprar as cartas.}
\item[Momento]{O jogador apenas pode jogar se tiver comprado as cartas em função de cartas Comprar Duas ou Comprar Quatro; punições vindas de outras regras (Regra do Silêncio, do Monte, da Velocidade, etc) não contam.}
\end{description}

\paragraph{Indicado} 

para jogadores que preferem o jogo desta forma; não recomendado para uso geral.

\subsubsection{Regra da Insatisfação [INS]}

Esta regra determina que, se um jogador não puder no momento jogar carta alguma, ele deve comprar uma carta do monte, mas não pode jogá-la.

\paragraph{Consequências}

Esta regra, ao contrário das regras da Satisfação e do Consolo, \emph{aumenta} o controle de jogo do ponto de vista tático: se o jogador A joga uma carta sabendo que o jogador B será obrigado a comprar uma carta, ele também sabe que o jogador B nada poderá fazer para modificar a pilha de descarte, e então poderá calcular o impacto da jogada também no jogador C.

Contudo, tamanho controle que mexe no equilíbrio das coisas pode ser negativo: impedindo que o hipotético jogador B tenha também alguma influência no jogo tem-se como o resultado uma diminuição do fator de imprevisibilidade do jogo que também o torna mais difícil e divertido.

\paragraph{Possíveis Contradições e Conexões}

\begin{description}
\item[Regra do Consolo]{Se aplicada com a Regra do Consolo, as cartas compradas por causa de uma punição (Comprar Duas ou Curinga Comprar Quatro) podem ser jogadas, como dita a Regra em questão.}
\item[Satisfação]{Esta regra não pode ser aplicada junto à regra da Satisfação; elas são diametralmente opostas.}
\end{description}

\paragraph{Indicado} 

Para quem prefira um jogo mais previsível e controlado; não recomendado para uso geral.

\subsubsection{Regra do Duplo Descarte [DED]}

Nesta variação bem diferente de UNO, existem duas pilhas de descarte: uma em que as cartas são descartadas viradas para baixo (de forma que ninguém pode ver que carta foi descartada) e uma que funciona como no jogo normal.

Os jogadores escolhem em que pilha depositarão suas cartas. Eles podem, portanto, mentir: descartar uma carta que não pode ser descartada na pilha de descarte normal, na pilha virada para baixo. Contudo, se os outros acharem que um jogador está mentindo (e descobrirem que está mesmo), o mentiroso compra quatro cartas e o desafiante joga fora uma carta qualquer de suas mãos. Se esse jogador não mentiu, contudo, quem o acusou compra quatro cartas e ele descarta mais uma.

\paragraph{Consequências}

Em primeiro lugar, esta regra insere, de certa forma, o mesmo tipo de complexidade estratégica inserida em outras regras já vistas como a do Descarte Rápido ou a do Corte: A possibilidade de se livrar de cartas mais rapidamente do que os outros. Contudo, esse objetivo, com essa regra, torna-se um pouco mais difícil para que possa ser considerado uma consequência direta da regra.

O que é mais observável é a mudança da dinâmica da própria partida: esta regra pode deixar o jogo mais divertido não por qualquer propriedade de cunho estratégico, mas porque os jogadores vão, a cada jogada ``oculta'', fazer julgamentos quanto à jogada de cada um. Descobrir quem está mentindo e quem está contando a verdade torna-se uma espécie de prioridade, e a diversão do jogo se concentra mais nesse ponto do que necessariamente na vitória em si.

Uma outra consequência estratégia extremamente importante é a possibilidade de descartar uma carta, mas \emph{não ter de lidar com os seus efeitos}: por exemplo, você pode jogar uma carta Inverter sem ter que \emph{de fato} inverter o jogo. Jogar um 7 sem deixar todos quietos (usando a Regra do Silêncio como exemplo de regras alternativas). Os cenários são inúmeros: a tática do consumismo não precisa mais ser empregada quando você não deseja lidar com as consequências de jogar uma carta.

\paragraph{Possíveis Contradições e Conexões}

\begin{description}
\item[Cartas Válidas]{Várias regras fazem referência à pilha de descarte ou à cor da pilha de descarte. Com a Regra do Duplo Descarte, a pilha de descarte em que as cartas são \emph{visíveis} é a válida para todas essas situações. Mesmo quando se descobre qual carta foi jogada na pilha de descarte oculta, essa carta não é transferida para a pilha de descarte oficial, e o jogo prossegue da última carta que foi jogada de forma visível.\footnote{Como visto anteriormente, é válido destacar que as consequências das cartas que foram omitidas não são aplicadas uma vez que se descobre que cartas eram.}}
\item[Vitória]{É possível ``ficar UNO'' e vencer o jogo descartando no monte de cartas ocultas. Ao vencer o jogo (com uma carta legítima) de forma oculta e ser desafiado, as cartas que o desafiante comprará a mais contarão pontos no final da partida. Ao vencer o jogo de forma oculta, mentindo, e for descoberto, a carta volta para sua mão, os mecanismos se aplicam, o jogo prossegue.}
\item[Vitória por Consequência]{É possível ganhar por consequência de uma acusação: se você tiver uma carta na mão, acusar uma mentira, e estiver certo (a pessoa de fato for mentirosa), você descarta a única carta que tem em mãos e vence o jogo.}
\item[Várias Cartas]{Ao jogar com regras como o Descarte Rápido ou a Primeira Regra da Aritmética --- regras em que, em geral, várias cartas podem ser jogadas ao mesmo tempo --- a primeira carta da sequência de cartas a ser jogada pode ser jogada na pilha visível, e as outras, na invisível, ou o contrário; qualquer combinação de partes é possível.}
\item[Regra do Passe e Inferno]{Se jogada em conjunto com a Regra do Passe ou do Inferno, a lógica é a seguinte: supondo que um jogador A jogue uma carta Comprar Duas. O jogador B pode comprar as duas cartas ou jogar uma carta que possa acumular com a carta já jogada; caso ele não possua tal carta, ele pode blefar, jogar uma carta de forma oculta e dizer que a carta é um Comprar Duas. Se o jogador C duvidar, os mecanismos se aplicam; se o jogador C não duvidar, ele deve comprar quatro (como se a carta do jogador B fosse, de fato, um Comprar Duas. O jogador C, por sua vez, pode escolher blefar também, dando continuidade ao ciclo de mentiras, até que alguém duvide e interrompa o acúmulo de cartas\footnote{Note-se que nem todas as jogadas subsequentes serão, necessariamente, mentiras; um outro jogador poderá, numa jogada futura, jogar de fato um Comprar Duas, mas de forma oculta, justamente esperando que alguém duvide dele. Além disso, note também que se este ciclo se prolongar, ter uma boa noção de quantas cartas Comprar Duas ainda estão em jogo (não foram descartadas até o começo deste ciclo de acúmulo) pode ser uma boa ideia.}}
\end{description}

\paragraph{Indicado} 

para jogos em grupos de amigos, que queiram um jogo mais divertido baseado na dinâmica de verdades e mentiras durante a partida.

\subsubsection{Regra do Gêmeo Traidor [GT]}

Essa é uma regra curiosa encontrada na internet: Se você jogar uma carta de valor numérico, e a próxima pessoa jogar uma carta \emph{idêntica}, você deve comprar o número de cartas indicado pelo valor numérico da carta em questão.

\paragraph{Consequências}

Esta carta não adiciona realmente nenhuma complicação estratégica no jogo além do fato de que ao jogar haverá mais a considerar se uma das possibilidades for jogar uma carta idêntica. A regra tampouco parece adicionar ao jogo um atrativo muito grande em termos de diversão.

\paragraph{Possíveis Contradições e Conexões}

\begin{description}
\item[Regra do Corte]{Se jogada com a Regra do Corte: realmente não importa quem jogar uma carta idêntica, se o próximo jogador na ordem natural ou um que corte o jogo --- você ainda deve comprar o número de cartas indicado pela carta.}
\item[Carta Única]{Se um jogador A jogar um 4 Vermelho, o próximo jogador deve jogar exatamente o outro 4 Vermelho no baralho para que o jogador A compre quatro cartas; ou seja, a carta idêntica a ser jogada precisa ser jogada, sozinha, apenas ela. Qualquer aplicação de regras como o Descarte Rápido ou a Primeira Regra da Aritmética anulam o efeito da Regra do Gêmeo Traidor.}
\end{description}

\paragraph{Indicado} 

para grupos de jogadores que se divertiriam com a regra.

\subsubsection{Regra do Passe [P]}

A Regra do Passe é, sem dúvida, a mais famosa, adotada e debatida regra alternativa para o UNO. Para alguns, é considerada um vulgarismo que leva a um jogo irregular e injusto; para outros, um elemento estratégico fundamental sem o qual o jogo torna-se simples e chato demais.

Existem quatro variações da Regra do Passe, e a diferença entre elas é a profundidade com a qual o princípio do Passe é aplicado. O princípio é simples: ao ser descartada uma carta que obriga o próximo jogador a comprar um determinado número de cartas, este próximo jogador não \emph{precisa} comprar as cartas, uma vez que pode adicionar a essa carta uma outra de natureza semelhante, passando o ``fardo'' da compra, agora mais ``pesado'', para o próximo jogador, que pode se defender da mesma maneira, e assim sucessivamente até que alguém seja obrigado a comprar o número \emph{acumulado} de cartas.

\paragraph{Primeira Variação (P1)} Todas as cartas Comprar Duas podem ser jogadas em cima de outras Comprar Duas. Nenhum Curinga Comprar Quatro pode ser jogado em cima de Comprar Dois ou outro Curinga Comprar Quatro, e nenhuma Comprar Duas pode ser jogada em cima de um Curinga Comprar Quatro. O Curinga Comprar Quatro deve ser jogado de forma isolada, sem acumular com nenhuma outra carta que tenha vindo antes, e o próximo jogador é obrigado a comprar as quatro cartas. As cartas Comprar Duas, contudo, podem acumular.

\paragraph{Segunda Variação (P2)} Cartas Comprar Duas podem ser jogadas em cima de outras Comprar Duas, e Curingas Comprar Quatro podem ser jogados em cima de Curingas Comprar Quatro, mas esses dois tipos de cartas não se misturam, de modo que um Curinga Comprar Quatro não pode ser jogado em cima de uma Comprar Duas, e uma Comprar Duas também não pode ser jogada em cima de um Curinga Comprar Quatro.

\paragraph{Terceira Variação (P3)} Cartas Comprar Duas podem ser jogadas em cima de outras Comprar Duas, Curingas Comprar Quatro podem ser jogados em cima de Curingas Comprar Quatro e em cima de Comprar Duas, mas Comprar Duas não podem ser jogadas em cima de Curingas Comprar Quatro.

\paragraph{Quarta Variação (P4)} Qualquer carta que exija que o próximo jogador compre cartas pode ser jogada em cima de qualquer outra carta desse tipo, sem restrições.

\paragraph{Consequências}

As consequências são profundas e a dinâmica de jogo se altera: se antes todos os jogadores são presas fáceis das cartas Comprar, agora elas adquirem um papel estratégico fundamental: o de \emph{defesa}. Ao mesmo tempo em que é bom que sejam descartadas caso haja o perigo de um adversário ganhar o jogo, é também fundamental guardá-las para poder se defender de outros Comprar Duas.

Além disso, ao saber que outras pessoas também podem se defender, você pode calibrar o \emph{alcance} dos seus ataques. Sem esta regra é impossível para o jogador A fazer o jogador C comprar cartas diretamente; o máximo que se poderia fazer é tentar incentivar o jogador B a usar uma carta Comprar Duas ou Comprar Quatro. O jogador A, agora, pode jogar uma Comprar Duas, se souber que o jogador B também possui uma carta semelhante, e então o jogador C comprará quatro cartas. As possibilidades de gerenciamento de jogo se expandem e ficam muito mais interessantes.

O uso da carta Comprar Duas torna-se muito mais complexo e sua importância aumenta exponencialmente. Ter várias cartas Comprar Duas significa poder desarmar os oponentes (lançar uma das cartas, apenas para que os oponentes sejam obrigados a usar suas eventuais cartas Comprar Duas) e então atacá-los mais tarde.

\paragraph{As Variações}

Além de a própria regra causar debate, suas próprias variações são constante motivo de disputa entre defensores: qual será a melhor maneira de arranjar a Regra do Passe?

Bom, isto claramente depende do objetivo do jogo. A Primeira Variação forja um equilíbrio interessante: às cartas Comprar Duas é reservado o poder de acumular, mas a carta Comprar Quatro continua mantendo seu poder de inefabilidade: quem a possui faz o próximo jogador comprar, e não há como \emph{contestar} o poder desta carta. A Quarta Variação renega este tratamento especial, e nenhuma carta mais possui esse poder especial, o que banaliza um pouco a característica de raridade e relevância do Comprar Quatro. Por outro lado, a Quarta Variação é a mais comumente responsável por provocar compras gigantescas de 18, 20, 24 cartas de uma só vez, o que pode criar um jogo muito divertido para um grupo de amigos.

Estrategicamente, entretanto, as três primeiras variações são as mais ricas e proveitosas, sendo a primeira a melhor delas. Isto porque, na segunda variação, temos que o Curinga Comprar Quatro pode ser combatido por outro Curinga Comprar Quatro (o que não seria tão frequente, considerando a quantidade dessas cartas no baralho), e na terceira variação, o Curinga Comprar Quatro poderia se sobrepor a uma carta Comprar Duas.

O que ocorre é que apenas na primeira variação valeria à pena guardar o Curinga Comprar Quatro para a última rodada\footnote{Estratégia ``Curinga Surpresa'', discutida acima.}, mas não tanto a ponto de outro Curinga Comprar Quatro (ou mesmo uma carta Comprar Duas) não conseguir anulá-lo; a ideia é não criar uma situação em que possuir dois Curingas Comprar Quatro, por exemplo, tornasse a vitória extremamente simples do que ela já acaba se tornando numa situação como essa, e ainda assim possibilitar a boa variedade estratégica que esta regra proporciona.

\paragraph{Possíveis Contradições e Conexões}

\begin{description}
\item[Descarte Rápido e Duplicatas]{Se jogada em conjunto com uma dessas regras, significa que várias cartas Comprar Duas e Comprar Quatro podem ser jogadas ao mesmo tempo --- tanto em uma jogada original quanto em uma que seja uma \emph{resposta}, um acúmulo a uma jogada anterior. Assim, se o jogador A joga um Comprar Duas Verde, o jogador B pode jogar duas cartas Comprar Duas Azul de uma só vez. Lembre-se, entretanto, da velha questão do equilíbrio entre descartar e guardar; com a Regra do Passe possui uma carta Comprar é muito mais importante e jogar todas que você possui de uma só vez é definitivamente algo não-sábio.}
\end{description}

\paragraph{Indicado} 

Indicado para uso geral, embora as três primeiras variações para um jogo mais estratégico, e a quarta para um jogo mais divertido e descontraído.

Como alguns jogadores podem acabar (em qualquer uma das variações) comprando muitas cartas, é recomendável aplicar também regras que permitem o descarte de mais cartas em uma mesma rodada, como a Regra do Descarte Rápido, a Primeira Regra da Aritmética, entre outras.

\subsubsection{Regra da Velocidade [V]}

Nesta regra, nenhum jogador pode demorar mais do que cinco segundos para descartar sua carta, contando a partir do momento em que o mais recente jogador descartou a dele. Caso demore, deve comprar duas cartas.

\paragraph{Consequências}

As consequências não são fortes no que tange àquilo que se \emph{faz} no jogo, mas sim ao \emph{modo} como se faz: sem tempo para pensar (ou com tempo, porém pressionado sob um limite artificial que faz o tempo parecer ínfimo), o jogador acaba errando, não fazendo as melhores escolhas ou fazendo as piores, enfim; esta regra cria um jogo displicente do ponto de vista tático, mas pode ser divertido para grupos que notam em seus jogadores a presença de alguns que costumam demorar muito; este jogo pode acabar servindo como um bom desafio para aqueles que queiram aprimorar a velocidade do pensamento estratégico em UNO.

Além disso, pode ser necessária a presença de alguém, que não esteja jogando, para marcação do tempo.

\paragraph{Possíveis Contradições e Conexões}

\begin{description}
\item[Tempo para Tudo]{Se a vez de um jogador consistir na compra de cartas, esta deve ser feita dentro dos cinco segundos também. Todas as atividades do jogo (como a escolha da cor se um Curinga for jogado ou a jogada de várias cartas simultaneamente por causa da Regra do Descarte Rápido, etc) devem também ser feitas dentro do mesmo limite.}
\end{description}

\paragraph{Variações}

O tempo de 5 segundos é arbitrário: outro limite de tempo pode ser estipulado.

\paragraph{Indicado} 

para treinar o pensamento rápido, ou verificar quem ainda é capaz de forma estratégica mesmo sob a pressão do tempo, ou para formular jogos rápido, para jogadores sem muito tempo.

\subsubsection{Primeira Regra da Paz [PAA]}

Com a aplicação desta regra, pode-se evitar um cenário que consegue se repetir com espantosa frequência: a partir da ação precoce de cartas Inverter, Pular e Comprar, um jogador pode demorar duas ou mesmo mais rodadas até que tenha a chance de jogar.

Para que todos tenham a chance de jogar pelo menos uma vez, esta regra impede que qualquer jogador jogue uma carta de ação na primeira rodada de cada partida.

\paragraph{Consequências}

Como consequência óbvia, todos terão a chance de jogar pelo menos uma vez no começo da partida. Isso possui uma conclusão um pouco difícil de perceber: o jogo fica \emph{um pouco mais} a cargo de escolhas e não de sorte; muitas vezes o fato de um jogador ter pego uma excelente mão no início do jogo faz com que os adversários comprem, não joguem, e depois fica muito mais difícil reverter essa vantagem.

Esta regra faz com que cada jogador tenha pelo menos uma chance de intervir no jogo antes que a sorte faça seu trabalho e, aquele que eventualmente tenha uma mão melhor, se usar de boa estratégia, ganhe (tendo mais chances de evitar isso).

É preciso lembrar também que cada ação tem suas consequências e, além de ser muito difícil vencer \emph{muito rapidamente} sem regras alternativas (como a do descarte rápido e a do corte, por exemplo), o fato de que outros jogadores sejam obrigados a comprar cartas no começo do jogo, por exemplo, pode acabar por ampliar o escopo deles em termos de gerenciamento de jogo --- o que não acontece tanto com a carta Pular, contudo.

\paragraph{Possíveis Contradições e Conexões}

\begin{description}
\item[Impossibilidade]{Esta regra \emph{não} determina uma \emph{punição} para quem joga uma carta de ação na primeira rodada do jogo, ela torna isso \emph{impossível}, de forma que se alguém não puder, com esta regra, jogar nenhuma carta na primeira rodada, deve comprar uma carta.}
\end{description}

\paragraph{Variações}

O limite de não poder jogar cartas de ação na primeira rodada é arbitrário; podem ser duas, três ou várias rodadas, ao gosto dos jogadores.

\paragraph{Indicado} 

Para jogadores que queiram adicionar um elemento tático a mais, ou que entendam que o jogo fique mais justo desta forma.

\subsubsection{Regra do Inferno [I]}

Nesta regra --- que faz jus a sua denominação --- qualquer carta de ação pode ser jogada em cima de qualquer carta de ação --- assim, uma Inverter pode ser jogada em cima de uma Comprar Duas ou de um Pular, que por sua vez podem ser jogadas em cima de um Curinga Comprar Quatro ou de outro Inverter, e por aí vai.

\paragraph{Consequências}

Em suma, o jogo torna-se mais bagunçado, acelerado, e menos refinado. A partir do momento em que cartas de natureza diferente possam ser jogadas uma em cima da outra, muitas possibilidades se abrem, e é difícil jogar taticamente dessa forma; o jogo torna-se pouco desafiante e muito imprevisível.

\paragraph{Possíveis Contradições e Conexões}

\begin{description}
\item[Regra do Passe implícita]{A Regra do Passe (especificamente, sua quarta variação) está implícita na Regra do Inferno.}
\item[Anulação das compras]{Se uma carta Comprar for jogada em cima de outra Comprar, seus valores são acumulados para o próximo jogador. Entretanto, este próximo jogador pode se defender utilizando uma carta Inverter ou Pular, anulando completamente o efeito das cartas anteriores. Cartas Comprar Duas podem ser jogadas em cima de qualquer carta Inverter ou Pular, não importando sua cor.}
\item[Descarte Rápido]{No caso da Regra de Descarte Rápido ser aplicada, não são todas as cartas de ação que podem ser jogadas juntamente; nesse caso, os mecanismos do Descarte Rápido recebem prioridade.}
\end{description}

\paragraph{Indicado} 

para grupos que se divertiriam com essa regra. Não recomendado para uso geral.

\subsubsection{Regra de Chukcki [CHU]}

Esta é uma regra que se refere a um detalhe pequeno que faz do UNO um jogo mais requintado, possuidor de uma espécie de encanto estratégico que apenas os jogos mais antigos podem se dar ao luxo de conter. Nesta elegante regra, se duas cartas de mesmo valor numérico forem jogadas por jogadores diferentes, porém consecutivamente, o terceiro jogador da sequência deve obrigatoriamente mudar a cor na pilha de descarte. O quarto, entretanto, caso o terceiro não tenha sido capaz de fazê-lo, pode continuar o jogo normalmente. 

\paragraph{Consequências} 

Esta regra adiciona um detalhe estratégico: algo que faz uma sutil diferença tática. Ao escolher jogar uma carta de valor numérico igual na pilha de descarte, você condena o próximo jogador a mudar a cor da pilha de descarte. 

\paragraph{Possíveis Contradições e Conexões}

\begin{description}
\item[Regras da Aritmética]{Se jogada em conjunto com quaisquer Regras da Aritmética, não importa que combinação ou carta seja jogada após uma sequência de duas cartas numericamente iguais (descartadas por jogadores diferentes), ela \emph{precisa} acabar por mudar a cor na pilha de descarte.}
\item[Duplo Descarte]{Se jogada em conjunto com a regra de Duplo Descarte, significa que se duas cartas de mesmo valor numérico, jogadas por jogadores diferentes, consecutivamente, estiverem na pilha \emph{visível} de cartas, a cor na pilha de descarte deve mudar. Esta carta (ou grupo de cartas, dependendo da combinação de regras alternativas) que muda a cor da pilha de descarte pode ser jogada na pilha visível ou na oculta. Assim, uma jogada não-legítima pode ser feita, mas se for descoberta os mecanismos da Regra do Duplo Descarte se aplicam.}
\end{description}

\paragraph{Indicado} 

Para quem deseja uma complexidade estratégica a mais no jogo.

\subsubsection{Regra do Batismo [BA]}

Assim como a Regra de Chukcki, a Regra do Batismo também adiciona ao jogo um elemento belo à dinâmica de jogo: Se uma carta Inverter for jogada depois de duas cartas de valores numéricos iguais (jogadas consecutivamente por jogadores diferentes), todos os jogadores compram uma carta.

\paragraph{Consequências}

Assim como a Regra do Carrossel (a ser vista mais adiante), esta regra acaba por adicionar um impedimento a uma carta; o jogador deve pensar antes de jogar uma carta Inverter em certas situações, ou também comprará uma carta. Ao mesmo tempo, fará todos comprarem uma carta, o que pode ser uma excelente estratégia às vezes.

\paragraph{Possíveis Contradições e Conexões}

\begin{description}
\item[Duplo Descarte]{Se jogada em conjunto com a regra de Duplo Descarte, a seguinte lógica é adotada: Duas cartas de mesmo valor numérico, jogadas por jogadores diferentes, consecutivamente, estão na pilha \emph{visível} de cartas. Se a carta Inverter for jogada na pilha oculta, nada acontece, mesmo se a carta for descoberta.}
\item[Descarte Rápido e Duplicatas]{Se várias cartas Inverter forem jogadas, uma carta por jogador deve ser comprada a cada carta Inverter.}
\item[Regra de Chukcki]{Se jogada em conjunto com a Regra de Chukcki, a carta Inverter a ser jogada precisa modificar a cor da pilha de descarte (algo que, portanto, só pode ser feito se a Regra do Descarte Rápido também estiver valendo).}
\end{description}

\paragraph{Indicado}

Para quem deseja uma complexidade estratégica a mais no jogo.

\subsubsection{Regra da Revelação [REV]}

Neste jogo, a estratégia torna-se crucial, porque agora as cartas dos adversários serão cada vez menos um mistério: o Monte é virado \emph{para cima}, de forma que todos os jogadores podem ver o que um jogador compra quando o faz.

\paragraph{Consequências}

As consequências são extremamente profundas e extensas: todos saberão a carta que um jogador comprará caso o faça; com memória o suficiente, é possível que boa parte do jogo de todos os jogadores seja conhecida caso o jogo se arraste por um bom tempo.

É possível saber qual será a próxima carta a ser comprada, o que significa que algumas pessoas se aproveitarão disso para comprar propositalmente --- outras, sabendo que o risco é diminuído de um adversário conseguir comprar algo prejudicial, \emph{correrão mais riscos}.

É realmente difícil imaginar todas as consequências disso durante um jogo --- e também considerando todas as regras alternativas apresentadas aqui.

\paragraph{Possíveis Contradições e Conexões}

\begin{description}
\item[Visibilidade]{Todas as ações que são feitas com o Monte são feitas da mesma forma; não é porque ele está voltado para cima que os jogadores devem ver \emph{todas as cartas}. Ao comprar várias cartas, por exemplo, um jogador está autorizado a contar o mais discretamente possível as cartas que deve comprar e pegá-las sem que os outros jogadores consigam necessariamente ver todas elas. Usar de artifícios externos para proibir a visualização do monte, contudo (como panos para encobrir a visão) já não é permitido, bem como tentar atrapalhar, a qualquer momento, a visualização da carta do topo do monte.}
\end{description}

\paragraph{Indicado} 

Para jogadores que queiram experimentar um tipo de jogo diferente. Não recomendado para uso geral.

Recomenda-se também o uso da Regra do Desafio Polivalente (ver mais adiante) ou da Insatisfação (ver acima), para evitar que as compras sejam feitas de modo muito indiscriminado.

\subsubsection{Regras do UNO Sexy}

Existe uma lenda na internet que diz que se algo pode ser encontrado na rede mundial de computadores, então existe uma versão pornô da coisa em questão. A pesquisa não foi conduzida em profundidade (o que não é um trocadilho), mas é digno de nota mencionar a óbvia possibilidade de transformar o UNO em um jogo de caráter erótico.

Os meios para se alcançar tal resultado são variados e, por aleatórios e múltiplos (e não recomendado para todas as idades) que são, não serão abordados aqui. Porém, é preciso deixar registrado que são possíveis e, possivelmente, até mesmo populares por fãs mais ardentes (sem trocadilho, mais uma vez) do jogo.

\subsection{Regras que Modificam ou Inserem Funções em Cartas}

Estas regras se referem não a modificações na dinâmica da partida, mas a funções novas e resultados inusitados para cartas específicas. Nesta seção se encontram algumas das regras mais populares do jogo e suas variações.

\textbf{É necessário dizer}, contudo, que as cartas escolhidas para receberem as funções novas foram escolhidas arbitrariamente, ou em termos de popularidade (a regra do silêncio é, indiscutivelmente, muito relacionada com a carta 7. Isto pode ser percebido com uma rápida pesquisa na internet). Todos esses números podem ser modificados; contudo, é importante notar que usar os números 6 e 9 para qualquer uma dessas regras não é aconselhável para evitar mal-entendidos e erros no jogo, pois as cartas são muito similares, e regras que exigem pensamento rápido (e a maioria delas o faz) podem levar a erros frequentes se usarem esses números. 

\subsubsection{Regra do Silêncio [S]}

Uma regra muitíssimo popular. Com ela, todos os jogadores devem permanecer em silêncio a partir do momento em que um 7 for descartado. O silêncio define-se pela não produção de sons que tomem a forma de palavras: grunhidos, onomatopeias acidentais, gemidos, suspiros, roncos, cafungadas, flatulências e outros sons de cunho não integralmente linguístico são permitidos. Quem não cumpre a regra deve comprar duas cartas. Os jogadores só podem voltar a falar se outro 7 for descartado, anulando este efeito do primeiro.

Existem variações, contudo, quanto às penalidades, à duração do silêncio e à questão da palavra ``UNO'', que é (deveria ser) dita toda vez que um jogador fica com uma carta na mão. As variações são apresentadas na forma de escolhas, portanto cada ítem da lista a seguir, se escolhido, adiciona sua sigla à S que identifica a Regra do Silêncio\footnote{Algumas escolhas são mutuamente exclusivas.}.

\paragraph{Variações}

\begin{description}
\item[Penalidade Progressiva (PP)]{Se esta característica for escolhida, o jogador não apenas compra duas cartas quando fala, mas compra uma carta a cada palavra falada.}
\item[Silêncio Temporário (T)]{O silêncio dura apenas uma (ou uma quantidade a gosto dos jogadores) rodada. A rodada inicia a partir da vez do jogador subsequente àquele que jogou o 7, e termina quando o jogador que jogou o 7 terminar sua vez. Qualquer frase que este jogador diga para sinalizar que a rodada de silêncio terminou é permitida, e não é punida. Qualquer 7 jogado dentro desta rodada extende o período, criando uma nova rodada de silêncio, iniciando agora por um jogador diferente.}
\item[Exceção para o UNO (U))]{Dizer ``UNO'' é permitido --- ou então combina-se previamente um sinal visual que indicará que a pessoa só tem uma carta na mão (como bater três vezes na mesa ou algo assim). Se dizer UNO não for permitido, a pessoa deverá comprar duas cartas sempre que ficar com uma carta na mão (caso os jogadores percebam, é claro).}
\item[Silêncio Específico (SE)]{Os jogadores podem falar, mas não podem falar uma \emph{palavra específica}. Palavras frequentes ou importantes podem ser escolhidas, como ``não'', ``carta(s)'', ``ponto(s)'', etc.}
\item[Silêncio Soturno (SS)]{Nesta regra os jogadores até podem falar, mas não podem rir.}
\end{description}

\paragraph{Consequências}

Não existem muitas consequências estratégicas com esta regra alternativa; ela basicamente torna o jogo mais divertido. O aumento de atenção dos jogadores que ela provoca, entretanto, é comparável à da regra do corte: um jogo com muitas pessoas pode se beneficiar dessa regra para captar a atenção de todos, o tempo todo. Contudo, a regra do silêncio não possui o efeito negativo da regra do corte de provocar pressa ao jogar, conservando, assim, a possibilidade de se fazer jogadas mais bem arquitetadas.

\paragraph{Possíveis Contradições e Conexões}

\begin{description}
\item[Descarte Rápido e Duplicatas]{Em variações da regra que dependam da carta 7 para ativar e desativar o efeito, um número par de cartas 7 jogadas não causa efeito algum, e um número ímpar ativa o efeito (ou desativa se ele já estava ativado).}
\end{description}

\paragraph{Indicado}

para grupos que queiram um jogo mais divertido.

\subsubsection{Regra do Carrossel [CAR]}

Nesta regra, toda vez que uma carta 0 for jogada, a mão de cada jogador deve passar a ser agora a mão do próximo jogador, na sequência momentânea da partida; todos, portanto, trocam de mão.

Usando do jogo-exemplo entre A, B, C, D e E (que jogam nesta ordem), se qualquer um deles joga um 0, a mão do jogador A passa para o jogador B, a do B passa para o C, e assim por diante, até que a mão do jogador C passe a ser a do jogador A. Caso o sentido esteja, naquele momento, invertido, a mão do jogador E passa a ser a mão do jogador D, e assim por diante.

\paragraph{Consequências}

O jogo torna-se um tanto imprevisível, uma vez que qualquer zero jogado por \textit{outro} jogador quebra toda a sua estratégia, que deverá ser recomeçada do zero.

Essa regra, contudo, não é de forma alguma integralmente prejudicial; quando os \emph{outros} possuem um zero, isso é motivo de preocuopação. Possuir um zero, contudo, é uma vantagem gigantesca, e possibilita a estratégia do Roubo de Cena, vista mais adiante.

Outra questão interessante dessa regra é que cada jogador saberá quais são as cartas do jogador à sua frente, ajudando assim na tarefa de gerenciamento do jogo (por algumas rodadas), caso ela seja assumida por algum dos jogadores.

O zero é usado nesta regra porque existem menos dessas cartas no baralho de UNO, tornando a incidência desse evento de troca de cartas menor. Dessa forma, não é recomendável alterar esta carta para esta regra.

\paragraph{Possíveis Contradições e Conexões}

\begin{description}
\item[Obrigatoriedade]{Os efeitos dessa regra são obrigatórios; assim, aquele que jogou o zero não pode \emph{escolher} entre aplicar o efeito ou não; a troca de cartas \emph{vai} acontecer. Isso é estrategicamente importante, pois jogar a carta não é \textit{sempre} bom, e portanto não se pode sempre esperar muito para usar a carta, ou talvez ela acabe atrapalhando a própria tática.}
\item[Fim de Jogo]{O Zero não pode ser descartado para vencer o jogo.}
\item[UNO]{Quando um jogador, depois do efeito da regra ser aplicado, ficar com uma carta na mão, ele deve dizer ``UNO''. Assim, o jogador A fica com apenas uma carta na mão e fala ``UNO''. O jogador B joga uma carta 0 e passa a ter uma carta na mão (a carta que tinha ficado na mão do jogador A) --- o jogador B deve dizer ``UNO'' assim que jogar a carta 0, ou qualquer pessoa pode acusá-lo de não dizer UNO e ele deverá comprar duas cartas.}
\item[Dois Jogadores]{Em jogos com dois jogadores, o zero faz com que os dois jogadores troquem de cartas.}
\item[Parcerias]{Em jogos com parcerias e quatro jogadores ao todo, um zero significa que uma dupla fica com as cartas da outra e vice-versa (embora na mesma ordem do jogo, como se cada jogador jogasse individualmente; os jogadores de uma dupla não podem trocar de mãos.}
\item[Regra do Corte]{Zeros não possuem cartas iguais, logo não podem ser cortados. Entretanto, se uma carta diferente for usada, a cada carta descartada (seja ela por corte ou não, em sequência ou não), as mãos devem trocar.}
\item[Regra da Troca]{A não ser que as duas regras usem cartas diferentes, elas não podem ser aplicadas simultaneamente.}
\end{description}

\paragraph{Indicado}

para grupos que queiram um jogo mais imprevisível e surpreendente.

\subsubsection{Regra da Troca [T]}

Esta é uma variação mais pessoal e estrategicamente controlável da Regra do Carrossel; nela, o jogador que descartar um zero pode escolher um jogador, e então trocar de cartas com ele. As cartas do jogador escolhido passam a ser as cartas do jogador que descartou o 0, e as cartas deste último passam a ser as cartas daquele.

\paragraph{Consequências}

Como dito no resumo da regra, essa é uma variante mais \emph{estratégica}: nem todos os jogadores são afetados, e é algo mais controlável: o jogador \emph{pode} escolher entre trocar ou não de mão (assim, se tiver boas cartas, ele não precisa se desfazer delas).

Esta regra também possibilita o Roubo de Cena, discutido mais adiante.

\paragraph{Possíveis Contradições e Conexões}

\begin{description}
\item[Fim de Jogo]{O zero, desta vez, pode ser usado para vencer o jogo, uma vez que a troca de mãos é apenas opcional.}
\item[UNO]{Quando um jogador, depois do efeito da regra ser aplicado, ficar com uma carta na mão, ele deve dizer ``UNO''. Assim, o jogador A fica com apenas uma carta na mão e fala ``UNO''. O jogador C joga uma carta 0 e escolhe ter a carta do jogador A; o jogador C deve dizer ``UNO'' assim que declarar que escolheu o jogador A, ou qualquer pessoa pode acusá-lo de não dizer UNO e ele deverá comprar duas cartas.}
\item[Dois Jogadores]{Em jogos com dois jogadores, o zero faz com que os dois jogadores troquem de cartas se aquele que descartou o zero quiser.}
\item[Parcerias]{É permitido que um jogador troque de mão com o seu parceiro se descartar um zero. Isso permite uma estratégia chamada Bússola, discutida adiante.}
\item[Regra do Corte]{Zeros não possuem cartas iguais, logo não podem ser cortados. Entretanto, se uma carta diferente for usada, a cada carta descartada (seja ela por corte ou não, em sequência ou não), as mãos devem trocar.}
\item[Regra do Carrossel]{A não ser que as duas regras usem cartas diferentes, elas não podem ser aplicadas simultaneamente.}
\item[Descarte Rápido]{Se a Regra do Descarte Rápido também for usada no jogo, e dois zeros forem trocados, a pessoa pode trocar de mãos duas ou mais vezes (dependendo do número de zeros jogados).}
\end{description}

\paragraph{Indicado} 

para quem deseja um jogo mais imprevisível, mas não exageradamente, e que tenha um componente tático mais preciso e inteligente.

\subsubsection{Regra do Monte [M]}

Nesta regra, toda vez que a carta 8 for descartada, todos os jogadores devem pôr a mão sobre a pilha de descarte. O último a bater (o que ficar com a mão por cima de todas as outras) deve comprar uma carta e o jogo prossegue normalmente.

\paragraph{Consequências}

Como muitas outras regras similares, esta regra aumenta o nível de atenção por parte de todos os jogadores. Não é particularmente crucial do ponto de vista estratégico.

\paragraph{Possíveis Contradições e Conexões}

\begin{description}
\item[Duplicatas e Descarte Rápido]{Duas ou mais cartas 8 de uma vez não significam que todos devem bater duas ou mais vezes; estas cartas acabam contando como uma.}
\item[Vencendo o Jogo]{Se ao vencer o jogo, um jogador descartar um 8 e não se lembrar de bater com a mão no monte, e for, no fim das contas, o último a bater, ele deve comprar uma carta e não vence o jogo --- o jogo prossegue.}
\item[Regra do Corte]{Se jogada junto com a regra do Corte: a partir do momento em que um dos jogadores bate no monte depois que um 8 for jogado, um corte não pode ocorrer até que alguém compre uma carta por ser a última a bater na pilha de descarte. Depois que as mãos foram recolhidas alguém pode cortar.}
\end{description}

\paragraph{Indicado} 

para grupos que queiram um jogo mais divertido.

Grupos grandes podem enfrentar dificuldades para proporcionar a todos a mesma chance de bater no monte de cartas, uma vez que os jogadores ficarão mais afastados do centro. Grupos com crianças muito jovens ou pessoas mais velhas podem enfrentar problemas também, uma vez que algumas pessoas não serão tão ágeis.

\subsubsection{Regra dos Elogios [E]}

Jogando com esta regra, a cada 5 verde jogado, um elogio (por adversário) deve ser feito à pessoa que descartou o 5. Se o elogio não for feito, o adversário compra uma carta.

\paragraph{Consequências}

Não muitas, do ponto de vista tático. Como é relativamente simples evitar a punição, é uma regra bem inexpressiva, estrategicamente.

É preciso ficar atento, contudo: nas primeiras vezes que a regra é aplicada, pode-se muito facilmente esquecer de elogiar o jogador. Nesse caso, o esquecimento deve render ao jogador uma carta extra a ser comprada.

\paragraph{Possíveis Contradições e Conexões}

\begin{description}
\item[Duplicatas e Descarte Rápido]{Se dois cincos verdes forem jogados ao mesmo tempo, dois elogios devem ser feitos por adversário.}
\item[Regra do Corte]{Se um cinco verde é cortado, apenas os elogios do segundo cinco verde devem ser feitos; o primeiro pode ser ignorado.}
\item[Repetição]{Os elogios não podem se repetir.}
\end{description}

\paragraph{Indicado} 

para jogadores que se divertiriam com essa regra --- afinal, ela não se declara contra ironias, por exemplo.

\subsubsection{Regra Invertida dos Elogios [IE]}

Com esta regra, a cada 5 vermelho jogado, uma crítica por adversário deve ser feitas à pessoa que descartou o 5.

\paragraph{Consequências}

As mesmas da Regra dos Elogios (ver acima).

\paragraph{Possíveis Contradições e Conexões}

Equivalentes às da Regra dos Elogios (ver acima).

\paragraph{Indicado}

para jogadores que se divertiriam com essa regra.

\subsubsection{Regra das Doações [DOA]}

Com esta regra, o jogador que descartar uma carta 2 escolhe duas cartas de sua mão para entregá-la a uma outra pessoa a ser escolhido pelo próprio jogador.

\paragraph{Consequências}

A regra não muda muito do ponto de vista estratégico, exceto que a carta 2 passa a ter primordial importância (ver mais adiante) no jogo. O jogo, contudo, fica mais divertido com a quantidade de ``generosas doações'' e ``dadivosas retribuições'' que ocorre.

É importante notar, também, que nem tudo é ótimo: às vezes, dar duas cartas da própria mão pode minar a própria estratégia, e as próprias cartas dadas podem servir para atacá-lo.

\paragraph{Possíveis Contradições e Conexões}

\begin{description}
\item[Obrigatoriedade]{Os efeitos dessa regra são obrigatórios; assim, aquele que jogou o 2 não pode \emph{escolher} entre aplicar o efeito ou não; a doação \emph{vai} acontecer.}
\item[Dois Jogadores]{Em um jogo com dois jogadores apenas, aquele que jogou o dois necessariamente doará duas cartas para o outro jogador.}
\item[Momento da Doação]{É possível ``ficar UNO'' e vencer o jogo doando cartas. No primeiro caso, assim que a doação terminar a pessoa precisa se declarar UNO; no segundo caso, se a pessoa só tiver duas cartas (além do 2 que acabou de jogar) ou uma outra carta, ela vence o jogo. As cartas, uma vez que doadas para adversários, contam pontos.}
\item[Regra do Corte]{Se um 2 for cortado, isso não anula o efeito do primeiro 2, e a pessoa que o jogou ainda tem direito a doar duas cartas suas para outra pessoa. A pessoa que cortou também tem o mesmo direito, já que jogou o dois.}
\item[Descarte Rápido e Duplicatas]{Se alguma dessas regras for aplicada, isto significa que a pessoa deve receber o número de cartas relativo ao número de cartas 2 jogadas. Se três cartas 2 foram jogadas, seis cartas devem ser doadas de sua mão para outros jogadores. Um jogador pode receber mais de duas cartas.}
\end{description}

\paragraph{Indicado}

para jogadores que se divertiriam com essa regra.

\subsubsection{Segunda Regra da Paz [PAB]}

Esta outra regra visa criar um período de ``cessar-fogo'' durante a partida, e acaba criando um elemento estratégico fundamental: a possibilidade de impedir, sem muito esforço, que os adversários possam atacar uns aos outros (e, mais importantemente, a você). Esta regra determina que, quando uma carta 1 é jogada, nenhuma \emph{carta de ação} pode ser jogada até outra carta 1 ser jogada.

\paragraph{Consequências}

Assim como a Primeira Regra da Paz, ela cria um ambiente de segurança em que cartas de ação não podem ser jogadas. Diferentemente da primeira regra, contudo, ela não visa equilibrar os jogadores e proporcionar a mesma chance de jogar a todos, pelo contrário: a ideia é obter uma vantagem tática gigantesca. Se não for possível usar cartas de ação, o adversário não pode combatê-lo caso você esteja quase vencendo a partida.

\paragraph{Indicado} 

para jogadores que queiram adicionar ao jogo uma complexidade muito interessante ao jogo, do ponto de vista tático.

\subsection{Regras que Modificam o Fim da Partida}

A cada fim de partida, os pontos são contados, mas existem outras coisas a serem feitas?

\subsubsection{Regra do Movimento [MOV]}

Seguindo esta regra, as pessoas devem trocar de lugar depois do fim de cada partida. Existem variações que pretendem ordenar esta troca; a minimalista regra original prevê, contudo, apenas que os jogadores não continuem na mesma ordem que estavam antes.

\paragraph{Consequências}

Por vezes, certas relações entre jogadores que sentam lado a lado podem começar a se desgastar devido ao fato de que um jogador acaba mais frequentemente atrapalhando outro. Para evitar que a escolha de lugares influencie demais a dinâmica de jogo, esta regra pode ser interessante. 

\paragraph{Indicado} 

para todos os tipos de grupos, por ser uma regra, em geral, benéfica. Não é muito interessante, contudo, num jogo de parcerias, e completamente inócua num jogo de apenas duas pessoas. 

\subsection{Regras que Modificam ou Inserem Penalidades}

Ah, as penalidades! Ver os adversários encherem-se de cartas é maravilhoso, mas ver-se com uma dupla de cartas a mais pro esquecer-se de dizer ``UNO'' é terrível. Estas regras buscam inserir novas punições no jogo ou modificar a natureza de algumas penalidades já existentes.

\subsubsection{Regra do Desafio Polivalente [DP]}

O Desafio é um mecanismo interessante do UNO: Permite que um jogador possa ser penalizado por tentar enganar os outros ao obrigar que alguém compre quatro cartas propositalmente. Por que não estender a ideia do Desafio para o jogo todo?

Esta regra determina que qualquer jogador pode desafiar qualquer jogador que compre uma carta: a ideia é que um jogador seja punido com duas cartas quando compre uma carta do monte \emph{mesmo possuindo}\footnote{razões para tal decisão na seção ``O Consumismo'' acima.} na mão cartas para jogar. De forma semelhante, o jogador que desconfiou de outro, mas estava errado, deve comprar cartas também --- desta vez, quatro.

\paragraph{Consequências}

Essa regra desencoraja qualquer tática que passe pela lógica do Consumismo (ver acima), em que comprar, mesmo sendo desnecessário, possa ser uma boa ideia. De certa forma a regra empobrece o jogo, do ponto de vista estratégico: Ao restringir a pluralidade de estratégias encorajadas (embora não a de estratégias \emph{possíveis}), o jogo torna-se mais previsível, e a compra, que já é um risco, torna-se um risco ainda maior. É claro que um jogador dissimulador pode usar esta regra como uma arma, mas ainda assim este é algo mais difícil de fazer.

Por outro lado, esta regra cria um ambiente em que muito se debate, dentro da cabeça de cada jogador, se alguém está mentindo ou não. Esse jogo de adivinhações quanto às intenções dos outros acaba desenvolvendo um tipo de jogo com uma dinâmica focada na questão de verdades e mentiras (muito similar ao já descrito na Regra do Duplo Descarte acima), que pode ser divertido para grupos de amigos.

\paragraph{Possíveis Contradições e Conexões}

\begin{description}
\item[Curinga Comprar Quatro]{No caso do Curinga Comprar Quatro, as regras oficiais continuam valendo.}
\end{description}

\paragraph{Indicado} 

para grupos que queiram um jogo mais divertido, baseado na questão de descobrir mentiras e verdades nas jogadas alheias.

\subsubsection{Regra de Dori [MEM]}

A Regra de Dori pega pesado com os esquecidos de plantão ao obrigar aqueles que se esquecem de dizer UNO a comprar quatro, e não duas, cartas.

\paragraph{Consequências}

Basicamente, esta regra apenas faz dobrar a atenção dos jogadores para que não se esqueçam de dizer UNO. De certa forma, isso pode acabar tornando o jogo levemente mais chato, uma vez que vai ser agora mais raro alguém se esquecer de dizer UNO. 

\paragraph{Indicado} 

para grupos que percebem que esquecem vezes demais de dizer UNO e queiram punir aqueles que se esquecem --- o que implica que a palavra vai ser esquecida menos vezes, resultando assim num processo de quase treinamento, de condicionamento para que não se esqueça de dizer UNO.

\subsubsection{Regra da Prescrição [PRE]}

No UNO jogado com as regras originais (variação 11211), um jogador só pode ser punido com duas cartas por não dizer UNO se for pego antes que o próximo jogador comece seu turno. A Regra da Prescrição determina que um jogador que esqueceu de falar UNO pode ser pego a qualquer hora, mesmo dois ou mais turnos depois que sua vez já passou.

\paragraph{Consequências}

Esta regra acaba tornando o jogo levemente mais ``justo'', uma vez que é muito mais fácil escapar da punição por não ter falado UNO se ela só dura pouco tempo --- e o tempo geralmente levado desde a jogada de alguém até que o próximo jogador compre uma carta costuma ser muito curto.

Como consequência tangencial, esta regra pode aumentar a atenção dos jogadores.

\paragraph{Indicado} 

para grandes grupos ou para formar jogos em que mais atenção é requerida por parte dos jogadores.

\subsubsection{Regra do Orgulho [ORG]}

Esta regra obriga o jogador que pedir desculpas, por qualquer razão, a comprar duas cartas.

\paragraph{Consequências}

Essa regra tem poucas consequências estratégicas, mas pode ser divertida quando aplicada a um grupo com grande afinidade, em que a personalidade da maioria das pessoas é menos tolerante a ações consideradas ``violentas'' (como obrigar o outro a comprar quatro cartas, por exemplo). Curioso é que esta regra é encontrada na internet, geralmente, fazendo menção ao fato de que alguns jogadores odeiam quando outros pedem desculpas por suas jogadas, e isto é tido como um traço cultural que é visto como irritante durante um jogo de UNO.

\paragraph{Indicado}

para grupos que percebem ter certa predisposição a se importar demais com os sentimentos dos adversários (e não gostam muito disso).

\subsection{Estratégias Específicas}

Nessa seção discutiremos estratégias específicas que surgem a partir da aplicação das regras alternativas apresentadas acima.

\subsubsection{Grupo Surpresa}

Uma consequência tática da Regra do Descarte Rápido (e também da Regra das Duplicatas) é o surgimento de um novo tipo de estratégia específica que não está presente na versão normal do jogo. Análogo ao Curinga Surpresa (ver acima), o Grupo Surpresa aparece como uma dupla de cartas de mesmo número ou símbolo que é guardada por um jogador para ser utilizada no final do jogo, fazendo com que um jogador que tenha duas (ou mesmo três, quatro) cartas na mão vença o jogo sem precisar se declarar UNO.

Essa estratégia é extremamente eficaz por diversos motivos: em primeiro lugar, carrega consigo um elemento surpresa muito maior do que o Curinga Surpesa: quando você possui apenas uma carta na mão, todos de alguma forma desejam impedi-lo de vencer e as atenções se voltam, na maioria das vezes, para você. Com o grupo surpresa, você não é o alvo prioritário dos jogadores que gerenciam o jogo. Além disso, a estratégia possui riscos mínimos de que, ao trilhar o caminho que lhe deixará apenas com essas cartas iguais na mão, você tenha que comprar cartas. Acaba sendo um caminho mais flexível e fácil para a vitória.

\subsubsection{Valorização das Armas}

Um resultado muito interessante de muitas regras alternativas é a valorização de determinadas cartas, que se tornam mais do que cartas; são, antes, armas e escudos de grande valor estratégico.

Se no jogo normal uma carta Comprar Duas e, principalmente, uma Comprar Quatro, são muito importantes, elas se tornam ainda mais úteis e primordiais em um jogo com a Regra do Passe. Contudo, outras regras criam uma situação em que guardar certas cartas para um momento oportuno, criando uma estratégia geral mais bem arquitetada e meticulosa, é muito proveitoso.

Com a Regra da Troca, por exemplo, o zero ganha uma importância fundamental. Se algum dos jogadores surpreender e repentinamente ficar UNO antes de você, você pode usar o zero para reverter o quadro mais facilmente. Outra estratégia em que o zero também acaba sendo mais valorizado (embora não na mesma medida) é a aplicada sob a Primeira Regra da Aritmética, uma vez que a carta zero pode ser descartada junto com qualquer outra como parte de uma adição ou subtração.

Outros exemplos de cartas valorizadas são a carta Inverter na Regra do Batismo e a carta Dois na Regra da Doação.

\subsubsection{Roubo de Cena}

O Roubo de Cena pode ser efetuado quando se está jogando com a Regra do Carrossel ou com a Regra da Troca. Basicamente, o que se faz é trocar de cartas com alguém que já está UNO ou está com menos cartas do que você. A troca deve ocorrer num momento apropriado, para que pouco precise ser feito para que a vitória seja sua.

Uma clara desvantagem do Roubo de Cena é que no mínimo um dos adversários saberá quais cartas você tem e é provável que o adversário em questão vá tentar pará-lo; com sorte não haverá tempo para arquitetar tal contra-ataque.

\subsubsection{Abstenção}

No jogo de Truco, um jogador pode se decidir por não mostrar sua carta. Esse tipo de jogada seria muito bem vinda no UNO, e se você gosta de jogar com a regra do Duplo Descarte, tal estratégia é possível e pode ser muito útil.

A Abstenção consiste em esconder uma carta --- ou seja, descartá-la de forma oculta --- que \emph{poderia} ser jogada, mas que caso o fosse, acabaria por causar consequências indesejáveis. Dessa forma, você diminui seu número de cartas e o jogo permanece do jeito que mais lhe beneficia.

Não se esqueça que a jogada, apesar de oculta, deve ser legítima; um blefe não caracteriza uma Abstenção.

\subsubsection{Bússola}

Essa é uma estratégia usada no jogo em duplas, junto com a Regra da Troca. Se os dois parceiros tiverem zeros na mão (um cada parceiro, pelo menos), eles podem trocar de mão para salvaguardar a mão do parceiro que possui menos cartas, e descartar dela ainda mais cartas.

Suponha um jogo exemplo entre os jogadores A, B, C e D, duplas A e C, B e D. O jogador A possui 8 cartas, o jogador C possui 3 cartas. O jogador A tem relativa certeza de que o jogador B tentará forçar o jogador C a comprar cartas ou tentará impedi-lo de jogar, de alguma forma. Com um zero, o jogador A troca de mão com o jogador C, e aogra o jogador A possui 3 cartas; o ataque do jogador B não surtirá o efeito de diminuir a possibilidade da equipe de vencer mais rápido, pois o menor número de cartas ainda se manterá. O jogador B joga; agora o jogador C joga algo, e depois disso o jogador D.  O jogador A pode se manter jogando ou, se for possível, pode trocar de mão novamente. E agora o jogador C terá duas cartas na mão.

\subsection{Variações Prontas}

Aqui são apresentados códigos de variações prontas de UNO de acordo com objetivos diferenciados. Todas as regras alternativas são aplicadas tendo como base o UNO 11211.

\subsubsection{Escolhas Populares}

Existem algumas regras alternativas comuns que se popularizaram pelo modo como deixam o jogo mais divertido ou mais interessante, taticamente falando, ou apenas porque os jogadores negligenciavam certas regras do jogo oficial. As regras mais populares são:

\begin{itemize}
\item{TV}
\item{C}
\item{DD}
\item{P (todas as variações)}
\item{S}
\item{T}
\item{PRE}
\end{itemize}

\subsubsection{Amigos}

Com amigos costumamos jogar jogos mais descompromissados, divertidos, focados em uma dinâmica mais solta e querendo se divertir acima de tudo.

\begin{description}
\item[UNO TV-BU-C-DR-P3-PRE]{Um jogo mais simples e tranquilo, para grupos pequenos, mas ainda assim divertido.}
\item[UNO DR-P4-PAA-SU-CAR-T-M-DOA-PRE]{Um jogo pra grupos maiores, com mais regras relativas a cartas individuais (Regra do Silêncio, da Troca, do Monte), com a quarta variante da Regra do Passe para garantia de diversão quando alguém acabar comprando mais de 20 cartas\ldots}
\item[UNO DR-AA-AB-GT-P3-BA-T-M-PRE]{Um jogo pra grupos maiores, mas com jogadores com um maior senso de tática e mais ansiosos por vencer de modo engenhoso e pela sensação das reviravoltas do jogo. Com regras um pouco mais complexas, como a Regra do Batismo e as Regras da Aritmética.}
\end{description}

\subsubsection{Poker Face}

O jogo de UNO pode ser customizado de modo a se tornar um jogo em que descobrir quem está mentindo torna-se mais engraçado e estimulante do que ganhar.

\begin{description}
\item[UNO FL-DED-DP-MEM-P2]{Duplo Descarte e Desafio Polivalente, além de extras para dificultar ainda mais.}
\end{description}

\subsubsection{As Estratégias do UNO}

Quem gosta de UNO de verdade sabe que a parte mais interessante é achar um modo de vencer, apesar de todas as armadilhas do azar e o sucesso das estratégias inicias dos adversários. Uma variante que valoriza essa característica do jogo é apresentada a seguir.

\begin{description}
\item[UNO DD-P1-PAB-CHU-T]{Duplicatas, Regra do Passe, Segunda Regra da Paz, Regra de Chukcki e Regra da Troca.}
\item[UNO DR-P1-PAB-CHU-T]{Mesma variação, mas com Descarte Rápido para uma dinâmica ainda mais acelerada.}
\item[UNO AA-AB-P1-PAB-CHU-T]{Mesma variação, mas com as Regras da Aritmética.}
\item[UNO DR-P1]{Regra minimalista que apenas insere as Duplicatas e a primeira variação da Regra do Passe.}
\end{description}

\subsubsection{Sala de Espera}

Às vezes não temos muito o que fazer ou estamos esperando por algo, mas não temos tempo ou espaço suficientes para jogar um campeonato de UNO como se deve. Sendo assim, existem algumas variações para ajudar a criar jogos mais rápidos.

\begin{description}
\item[UNO TV-FU-C-DR-V-P3]{Menos cartas, sem precisar contar pontos, descartando o mais velozmente possível.}
\item[UNO MS-TV-C-DR-V-P3]{Mesma variante, porém com o número normal de cartas e com a Regra da Morte Súbita pra tornar o jogo mais rápido.}
\end{description}


\chapter{Apêndices}

\appendix

\section{Resumo das Estratégias do Jogo Oficial 1121}

\begin{itemize}
\item{O jogo é um constante equilíbrio pessoal entre jogar pela continuidade (para livrar-se das cartas na mão e vencer o mais rapidamente possível) e jogar pelo gerenciamento (para controlar as jogadas dos outros e evitar a vitória deles).}
\item{Quanto menor o número de cartas, menores as possibilidades de gerenciar o jogo, e vice-versa.}
\item{Guardar cartas importantes para um outro momento do jogo é uma boa estratégia, mas guardá-las por tempo demais é perigoso: Curingas contam 50 pontos no jogo e cartas de ação, 20.}
\item{Ao jogar de acordo com o princípio da continuidade, jogue sempre para tentar manter o maior número de possibilidades para a próxima rodada na sua mão. Dessa forma, se você pode escolher entre jogar uma carta A ou a carta B, jogue de forma a preservar suas possibilidades; se ao jogar a carta A você acaba com todas as cartas de uma determinada cor na sua mão, e isso não ocorre com a carta B, jogue a carta B.}
\item{Se com nenhuma delas você prejudica a variedade de cores na sua mão, você precisará ter uma boa noção do que já foi jogado antes para forjar uma estratégia mais consciente: jogue uma cor que já tenha sido bastante jogada (menos chance de que ela volte na próxima rodada) ou jogue um número já bastante jogado (menos chance de que eventualmente ele apareça de novo). Prefira sempre acabar com a variedade de cartas de mesmo número na mão do que com a variedade de cores.}
\end{itemize}

\newpage

\section{Arranjo das Cartas na Mão}

Aqui vai um conselho pessoal quanto à ordenação de cartas na mão:

\begin{itemize}
\item{Em primeiro lugar, cuidado com a postura, a posição da mão e com superfícies espelhadas: não mostre suas cartas aos oponentes!}
\item{Depois, ordene as cartas numéricas em ordem crescente da esquerda para a direita. Se houver mais de uma carta do mesmo número, procure deixar as cartas de um jeito que de uma carta para outra fique a mesma cor. Por exemplo, se você tiver um 1 Amarelo, um 1 Verde, um 3 Amarelo, um 3 Azul e um 4 Azul, deixe as cartas na ordem 1 Verde, 1 Amarelo, 3 Amarelo (para que duas cartas amarelas fiquem juntas), 3 Azul, 4 Azul (aqui temos as cartas azuis juntas também).}
\item{Insira as cartas de ação coloridas próximo a um grande grupo de cartas com a mesma cor (se possível), mas nunca entre duas cartas de mesmo número. Ou, se preferir, insira todas no início ou no final da mão.}
\item{Insira os Curingas entre duas cartas de valores numéricos diferentes e cores diferentes. Ou, se preferir, insira todas no início ou no final da mão.}
\end{itemize}

\newpage

\section{Resumo das Regras Alternativas}

\begin{center}
\begin{longtable}{|p{4cm}|p{5cm}|}
\hline
Torneio de Vitórias (TV) & Conta-se o número de vitórias, não os pontos. \\ \hline
Pontos Reversos (PR) & Todos começam com 500 e vão perdendo pontos. \\ \hline
Morte Súbita (MS) \emph{e variações} & Quem precisar comprar uma carta está eliminado da partida. \\ \hline
Final Limpo (FL) \emph{e variações} & É proibido vencer o jogo com Curingas ou cartas que forcem outros jogadores a comprar cartas. \\ \hline
Flash UNO (FU) \emph{e variações} & Todos começam o jogo com 4 cartas. \\ \hline
Battlefield UNO (BU) \emph{e variações} & Todos começam o jogo com 9 cartas. \\ \hline
Corte (C) & Uma carta idêntica à da pilha de descarte pode ser jogada por qualquer jogador, na vez de jogar ou não. O jogo prossegue a partir do jogador que fez o corte. \\ \hline
Descarte Rápido (DR) & Duas cartas de número ou símbolo iguais podem ser descartadas na mesma vez. \\ \hline
Duplicatas (DD) & Duas cartas idênticas podem ser descartadas na mesma vez. \\ \hline
Primeira Regra da Aritmética (AA) & Duas cartas numéricas de mesma cor podem ser descartadas se sua soma ou subtração (mais que duas cartas no caso da soma) resultarem no número da carta (necessariamente numérica) que está na pilha de descarte. \\ \hline
Segunda Regra da Aritmética (AB) & Uma carta numérica pode ser jogada se for a soma das duas últimas cartas na pilha de descarte se estas forem de mesma cor. \\ \hline
Satisfação (SAT) & É obrigatório comprar até ser comprada uma carta que possa ser jogada. \\ \hline
Consolo (CON) & Cartas compradas por causa de cartas Comprar Duas ou Comprar Quatro podem ser jogadas na mesma vez. \\ \hline
Insatisfação (INS) & Uma carta comprada nunca pode ser jogada. \\ \hline
Duplo Descarte (DED) & Os jogadores podem escolher descartar em uma pilha de descarte em que as cartas são visíveis para todos os jogadores, ou ocultar sua jogada em uma pilha de descarte em que as cartas são viradas para baixo. Punições são aplicadas se os jogadores descobrirem que uma jogada não-legítima foi feita. \\ \hline
Gêmeo Traidor (GT) & Se uma carta numérica idêntica à que está na pilha de descarte for jogada, o jogador que jogou a primeira carta compra o número de cartas da carta jogada. \\ \hline
Passe (P) \emph{e variações} & Cartas Comprar Duas Cartas e Curinga Comprar Quatro Cartas podem ser jogadas uma em cima da outra, acumulando. Variações decidem quais cartas podem ou não ser acumuladas. \\ \hline
Velocidade (V) \emph{e variações} & Todas as jogadas precisam ser feitas em 5 segundos. \\ \hline
Primeira Regra da Paz (PAA) \emph{e variações} & Ninguém pode descartar uma carta de ação nas primeiras rodadas do jogo. \\ \hline
Inferno (I) & Qualquer carta de ação pode ser jogada em cima de qualquer carta de ação. \\ \hline
Regra de Chukcki (CHU) & Se duas cartas de mesmo valor numérico forem jogadas consecutivamente por jogadores distintos, o terceiro jogador deve mudar a cor na pilha de descarte. \\ \hline
Batismo (BA) & Depois de duas cartas de valor numérico igual jogadas consecutivamente por jogadores diferentes, se o terceiro jogador jogar uma carta Inverter todos compram 1 carta. \\ \hline
Revelação (REV) & O MONTE fica virado para cima. \\ \hline
Silêncio (S) \emph{e variações} & A cada 7 jogado, os jogadores não podem falar até que outro 7 seja jogado. \\ \hline
Carrossel (CAR) & A cada 0 jogado, os jogadores trocam de mão, de forma que cada jogador passa sua mão para o próximo jogador na sequência. \\ \hline
Troca (T) & A cada 0 jogado, o jogador que o jogou pode escolher trocar de mão com outro. \\ \hline
Monte (M) & Quando um 8 é jogado os jogadores devem pôr a mão sobre a pilha de descarte. \\ \hline
Elogios (E) & Um elogio por adversário é feito a quem jogar um 5 Verde. \\ \hline
Regra Invertida dos Elogios (IE) & Uma crítica por adversário é feita a quem jogar um 5 Vermelho. \\ \hline
Doações (DOA) & O jogador que descartar uma carta 2 escolhe um adversário para o qual doar duas de suas cartas. \\ \hline
Segunda Regra da Paz (PAB) & Se uma carta 1 é jogada, não se pode jogar cartas de ação até que outro 1 seja jogado. \\ \hline 
Regra do Dado (DAD) & Se uma carta 4 é jogada, o próximo jogador deve jogar um dado e comprar o número de cartas equivalente ao resultado do dado. \\ \hline
Regra da Trindade (TRI) & Se uma carta 3 é jogada, o próximo jogador deve jogar uma outra carta 3 ou uma carta de mesma cor, sob pena de comprar 3 cartas. \\ \hline 
Regra do Backlash (BACK) & Se uma carta 8 é jogada, o jogador anterior deve comprar uma carta. \\ \hline 
Regra da Clarividência (CLA) & O jogador que joga uma carta 1 escolhe um jogador que deve mostrar a ele suas cartas. \\ \hline 
Movimento (MOV) & A cada fim de partida os jogadores devem mudar de ordem. \\ \hline
Desafio Polivalente (DP) & Os jogadores podem desconfiar quando um jogador compra uma carta. Existem penalidades para o caso de não haver necessidade de fazê-lo. \\ \hline
Regra de Dori (MEM) & Quem esquecer de dizer UNO deve comprar quatro cartas. \\ \hline
Regra da Prescrição (PRE) & Uma pessoa pode ser acusada de não ter dito UNO a qualquer momento após ter ficado UNO. \\ \hline
Regra do Orgulho (ORG) & Comprará duas cartas aquele que pedir desculpas durante a partida. \\ \hline
\end{longtable}
\end{center}

\newpage

\section{Regras Oficiais Comumente Ignoradas}

Durante uma partida casual de UNO, uma variedade de regras oficiais costuma ser ignorada por diversos motivos: ignorância quanto às regras, confusão entre regras oficiais e alternativas, má interpretação de regras oficiais ou a deliberada exclusão de determinadas regras. Esta seção documenta algumas dessas regras que são frequentemente deixadas de fora, seja para lembrar-se delas com facilidade ou para escolher conscientemente por sua exclusão\footnote{Regras que são ignoradas porque uma regra alternativa entra em vigor serão deixadas de fora deste apêndice: tratar-se-á exclusivamente das regras que não entram em vigor sem uma substituição.}.

\begin{description}
\item[Comprar uma carta e jogar uma que já estava na mão]{Segundo as regras oficiais (ver página \pageref{oficiais}), ``Se [uma carta comprada] puder ser jogada, a carta pode ser descartada na mesma vez, mas o jogador não pode jogar uma carta que estiver na sua mão depois de ter feito a compra''. Geralmente os jogadores compram uma carta como forma de estratégia (ver página \pageref{consumismo}), para obter uma carta melhor que as que já têm, mesmo quando alguma destas pode ser jogada. Depois que a compra se revelou infrutífera, ainda assim jogam uma das cartas que poderiam ter jogado antes.}
\item[Determinação da cor no começo do jogo]{Segundo as regras oficiais, ``Se um Curinga for aberto no começo do jogo, o jogador à esquerda do distribuidor determina a cor que dará sequência ao jogo''. Em geral, a pessoa que determina a cor costuma depender do modo como as cartas são distribuídas, e a regra é ignorada em favor de alguma preferência do grupo.}
\item[Sem duvidar]{Uma das minúcias do UNO é que um jogador que jogar uma carta Curinga Comprar Quatro Cartas não pode fazê-lo se tiver em mãos uma carta de \emph{mesma cor} que a pilha de descarte. Segundo as regras oficiais, ``um jogador que possuir um Curinga Comprar Quatro Cartas pode tentar blefar e jogar uma carta inapropriada, mas certas regras se aplicam se ele for apanhado''. Muitas vezes pode-se esquecer desta propriedade do jogo, e nenhum jogador duvidar da carta jogada.}
\item[Comprar depois do fim]{Segundo as regras oficiais, ``se a última carta jogada em uma partida for a carta Comprar Duas Cartas ou o Curinga Comprar Quatro Cartas, o próximo jogador precisa comprar 2 ou 4 cartas, respectivamente'', e essas cartas são contadas ``na totalização dos pontos''. Por vezes os jogadores pensam que após o fim da partida mais nenhuma penalidade se aplica.}
\item[Palpites]{Segundo as regras oficiais, ``jogadores que derem palpites aos outros jogadores precisam comprar 2 cartas do MONTE''. Esta regra é comumente ignorada, pois os jogadores costumam colaborar entre si com informação --- especialmente quando um dos jogadores está com apenas uma carta na mão.}
\end{description}

\newpage

\section{Ensinando UNO de forma rápida}

Aqui está um guia prático sobre como ensinar UNO para leigos de forma rápida:

\begin{enumerate}
\item{Explique que cada jogador começa com 7 cartas e vai descartando até não sobrar nenhuma na mão, e que quando a pessoa fica com uma carta na mão, deve dizer ``UNO''. Se alguém vê que ela não o fez e a acusa, ela deve comprar duas cartas. Quando o jogo acaba, os pontos das mãos dos adversários são contados e vão para o vencedor: o jogo acaba quando alguém chega aos 500 pontos.}
\item{Explique que você pode descartar uma carta de mesmo número, mesmo símbolo ou mesma cor que a da pilha de descarte. Se não houver nada a ser jogado, você compra uma carta; se puder jogá-la, você pode jogá-la. Se não, o próximo jogador inicia a vez.}
\item{Explique a função das Cartas de Ação e dos Curingas (e sua pontuação equivalente no término de uma partida).}
\item{Diga que existem punições caso alguém jogue um Curinga Comprar Quatro tendo uma carta de mesma cor para jogar, caso alguém dê palpites ou caso jogue equivocadamente.}
\item{Explique eventuais regras alternativas e clarifique que nem todos jogarão com estas regras. É uma boa ideia também jogar primeiro o jogo oficial e depois introduzir regras alternativas.}
\end{enumerate}

Explicar isso fazendo um jogo-exemplo entre pessoas que já sabem jogar e ir explicando os mecanismos do jogo enquanto o aprendiz observa é também uma excelente ideia.

\newpage

\section{Paciência UNO}

Os jogos de Paciência são famosos; são jogos de cartas em que apenas uma pessoa joga, e há uma maneira de usar o baralho de UNO para jogar uma espécie de jogo no estilo de Paciência --- uma mistura dos conhecidos jogos Klondike e Aranha, que ficaram famosos a partir da popularização de computadores pessoais.

Em primeiro lugar, é preciso tirar todos os Curingas e todas as cartas 0 do Baralho. Depois, você deve ver a carta Comprar Duas como o Valete (o \emph{J} do baralho comum), o Inverter como a Dama ou Rainha (o \emph{Q} do baralho comum), e o Pular como o Rei (o \emph{K} do baralho comum). Para se lembrar disso, pense que que Comprar Duas (C), Inverter (I) e Pular (P) estão, nessa ordem, em ordem alfabética por suas primeiras letras, e estão representando as cartas do baralho normal em ordem crescente.

Depois, é preciso entender também que não há carta 10, e portanto logo após o 9 vem o Comprar Duas.

Depois, distribua as cartas na mesa em uma linha horizontal da seguinte maneira: uma carta virada para cima na primeira coluna, uma virada para baixo e outra virada para cima na segunda coluna, duas viradas para baixo e uma virada para cima na terceira coluna, e assim sucessivamente, até que a última coluna tenha seis cartas viradas para baixo e uma carta virada para cima (o número de cartas viradas para baixo começa com 0 e aumenta 1 a cada coluna até o fim; há sempre apenas uma carta virada para cima em cada coluna). A primeira coluna fica à extrema esquerda; a última, à extrema direita.

As cartas que sobrarem são o Monte. Se você consegue ver alguma carta 1, você pode colocá-la em um dos 8 espaços acima das colunas que devem ficar reservados para as sequências de cartas. Assim que você tira uma carta visível de cima de sua pilha de cartas, você pode revelar a primeira das cartas viradas para baixo.

O objetivo é criar oito sequências de cartas de mesma cor (indo do 1 até o Pular) e empilhá-las acima das colunas. Primeiro você deve colocar lá a carta 1, e assim que o fizer, pode colocar a 2, e depois a 3, e assim por diante, até concluir uma sequência.

Para ir fazendo isso, você pode mover uma carta nas colunas para outra nas colunas, se elas forem subsequentes e de mesma cor: ou seja, um 7 azul pode ser colocado embaixo de um oito azul, e estas duas cartas podem, conjuntamente, serem colocadas debaixo de um nove azul. Com essa movimentação você vai oganizando as cartas e descobrindo as cartas que estão viradas para baixo.

Quando você não tiver mais o que mover nas colunas (ou mesmo se o tiver, é claro), você vira uma carta do Monte. Se ela puder ser colocada nos espaços acima das colunas ou debaixo de alguma carta das colunas, você pode fazê-lo. Se não, vire mais uma carta do Monte. Quando o Monte acabar, você pode virá-lo todo para baixo novamente e recomeçar a virar suas cartas. Isso pode ser feito quantas vezes forem necessárias.

Assim que uma coluna ficar vazia, apenas uma carta Pular pode preencher seu lugar. O jogo termina quando você consegue organizar oito sequências de mesma cor.

\newpage

\vspace{5cm}

\begin{center}
\begin{small}
Este livro tem como objetivo ser o documento definitivo da estratégia e das variações do UNO. Se você não concorda com algo que está aqui, se encontrou algum erro de digitação, ortografia, gramática ou organização, se você joga UNO com uma regra diferente não listada aqui, se encontrou alguma contradição ou tem simplesmente ideias para torná-lo mais organizado e quer contribuir, mande um e-mail para \\

\href{mailto:peterson.235@gmail.com}{peterson.235@gmail.com} \\

ou visite \\

\url{http://petercast.net/projetos/projects/show/uno} \\

para saber como ajudar. Você receberá créditos por sua contribuição, é claro, e ajudará esta obra coletiva a ganhar ainda mais credibilidade e confiabilidade.
\end{small}
\end{center}

\end{document}
